% 本文件由 LaTeX 排版 Agent 根据有效 translation.json 自动生成。
% author=John Chrysostom; work_id=2303; title=哥罗森书讲道集

\chapter{哥罗森书讲道集}

% structure: p0001 source=h1 target=omitted reason=page-title-already-rendered-by-parent

讲道一\newline{}讲道二\newline{}讲道三\newline{}讲道四\newline{}讲道五\newline{}讲道六\newline{}讲道七\newline{}讲道八\newline{}讲道九\newline{}讲道十\newline{}讲道十一\newline{}讲道十二\par

\SourceRecord{Translated by John A. Broadus.；Nicene and Post-Nicene Fathers, First Series；Vol. 13.；Edited by Philip Schaff.；Buffalo, NY: Christian Literature Publishing Co.,；1889.；<http://www.newadvent.org/fathers/2303.htm>.}

% structure: p0001 source=h1 target=section reason=page-title

\section{哥罗森书讲道（一）}

哥 1:1,2\par

\BibleQuote{保禄，因天主的旨意作基督耶稣的宗徒，与弟茂德弟兄，致书给在哥罗森的圣者及在基督内忠信的弟兄们：愿恩宠与平安由天主我们的父赐与你们。}\par

\Person{保禄}的所有书信确实都是神圣的；但那些他在被囚后所写的书信有一些优势：例如，致\WorkTitle{厄弗所书}和致\WorkTitle{费肋孟书}，致\WorkTitle{弟茂德书}、致\WorkTitle{斐理伯人书}以及我们面前的这一封。因为这一封也是在他在押时写的，因为他在信中写道：\BibleQuote{为这缘故我也被束缚，好使我照我所应当的显明它。} \BibleRef{哥 4:3} 但这封书信似乎写在致\WorkTitle{罗马书}之后。因为致\WorkTitle{罗马书}是他未曾见过他们时写的，但这封是在之后，而且接近他宣讲的尾声。由此可知，他在致\WorkTitle{费肋孟书}中说：\BibleQuote{我\Person{保禄}，这样一个老人} \BibleRef{费 1:9}，并为\Person{敖乃息摩}求情；但在这封中他派遣了\Person{敖乃息摩}本人，如他所说：\BibleQuote{同忠信可爱的弟兄\Person{敖乃息摩}} \BibleRef{哥 4:9}，称他为忠信的、可爱的、弟兄。因此他也在这封中放胆说：\BibleQuote{从你们所听到的、那传于天下万有的福音的希望} \BibleRef{哥 1:23}。因为那时福音已经传了很长时间。我想致\WorkTitle{弟茂德书}是在此之后写的，那时他已经到了生命的尽头，因为那里他说：\BibleQuote{因为我已被奠祭} \BibleRef{弟后 4:6}；但这封比致\WorkTitle{斐理伯人书}晚，因为在那封中他刚开始在罗马的囚禁。\par

\Person{保禄}为何说这些书信在这一点上有一些优势，因为它们是在被囚时写的？就好像一个斗士在杀戮和胜利中书写；实际上他就是这样做的。因为他自己也知道这是一件大事，所以他在致\Person{费肋孟}书中说：\BibleQuote{我在锁链中所生的} \BibleRef{费 1:10}。他说这话，是为了我们不应当在逆境中灰心，反而应当喜乐。这个地点是\Person{费肋孟}和这些（\Place{哥罗森}人）所在的地方。因为在致他的书中他说：\BibleQuote{和我们的战友\Person{阿尔希普}} \BibleRef{费 1:2}；在这封中：\BibleQuote{告诉\Person{阿尔希普}} \BibleRef{哥 4:17}。这个人似乎被委托在教会中担任一些职务。\par

\Person{保禄}却未曾见过这些人们，或者罗马人，或者希伯来人，当他给他们写信时。关于其他人，他在许多地方显示了这一点；关于哥罗森人，请听他所说：\BibleQuote{以及凡没有在肉身见过我面的人} \BibleRef{哥 2:1}；又：\BibleQuote{虽然我肉身不在，但精神上却与你们同在}。他以为自己无处不在是一件多么伟大的事。而且，即使他不在，他常常使自己临在。所以，当他惩罚那个淫乱者时，看看他如何将自己置于审判台上；\BibleQuote{因为}，他说，\BibleQuote{我身体虽不在，但心理上却已在，好像我已在场一样判定了他} \BibleRef{格前 5:3}；又：\BibleQuote{我要到你们那里去，并且要知道那些自高自大之人的，不是言语，而是能力} \BibleRef{格前 4:19}；又：\BibleQuote{不但在我与你们同在的时候，而且更是在我不在的时候} \BibleRef{斐 2:12}。\par

\Quote{保禄，因天主的旨意作基督耶稣的宗徒。}\par

从这封书信我们也可以说说，通过考虑，我们发现它的起因和主题是什么。那么它是什么呢？他们惯于通过天使接近天主；他们持守许多犹太和希腊的习俗。因此他在纠正这些。所以在开头他就说：\Quote{因天主的旨意。} 所以这里他再次用了\InnerQuote{通过}这个表达。\Quote{和弟兄\Person{弟茂德}，}他说；当然他也是宗徒，大概也为他们所熟知。\Quote{致在哥罗森的圣者。}这是弗里吉亚的一个城，由\Place{劳狄刻雅}靠近它而可知。\Quote{并在基督内的忠信弟兄。} \BibleRef{哥 4:16} 那么，他说，你从哪里成为圣人？告诉我。你从哪里被称为忠信者？岂不是因为你借着死亡而成圣？岂不是因为你在基督内有信仰？你从哪里成为弟兄？因为无论事实、言语或成就，你都没有显明自己是忠信者。告诉我，你怎么被托付了如此伟大的奥迹？岂不是因为基督？\par

\BibleQuote{愿恩宠与平安由天主我们的父赐与你们。} 恩宠从哪里来？平安从哪里来？\BibleQuote{从天主，}他说，\BibleQuote{我们的父。} 虽然他在此处没有使用基督的名字。\par

我要问那些贬低圣神的人，天主从哪里是仆人的父？谁完成了这些伟大的成就？谁使你成为圣人？谁使你成为忠信者？谁使你成为天主的儿子？那使你配得上被信赖的那一位，也是你被委托一切的原因。\par

因为我们被称为忠信者，不仅因为我们有信仰，而且因为天主将连天使在我们之前也不知道的奥迹委托给了我们。不过，对\Person{保禄}来说，是否这样表述无所谓。\par

第3节。\BibleQuote{我们感谢天主，我们的主耶稣基督的父。} \BibleRef{哥 1:3}\par

在我看来，他将一切归于圣父，以便他所说的不会立即冒犯他们。\par

\BibleQuote{常常为你们祈祷。}\par

他不仅通过感恩展示自己的爱，也通过不断的祈祷，因为他将那些他未见的人持续放在自己心里。\par

第4节。[ \BibleQuote{听到你们的信仰在基督耶稣内，} \BibleRef{哥 1:4}\par

[ 稍早他说了\Quote{我们的主}。\Quote{祂}，他说，\Quote{是主，不是仆人。} \Quote{耶稣基督的。} 这些名字也是他对我们恩惠的象征，因为\Quote{祂}意思是，\BibleQuote{要将自己的民族从他们的罪恶中拯救出来。} \BibleRef{玛 1:21} ]\par

第4节。\BibleQuote{听到你们的信仰在基督耶稣内，以及你们对所有圣者的爱。} \BibleRef{哥 1:4}\par

他已经安抚了他们。将消息带给他的乃是\Person{厄帕夫辣狄托}。但他却打发\Person{提希苛}送去这封信，而把\Person{厄帕夫辣狄托}留在了自己身边。他说：\Quote{论到你们对众圣人所有的爱德，}并不是对这一位或那一位：当然也包括对我们了。\par

第五节。\BibleQuote{因为那为你们存留在天上的望德。} \BibleRef{哥 1:5}\par

他谈论将来的美事。这样做是为了他们的试探，好叫他们不要在此世寻求安息。因为恐怕有人会说：\Quote{他们自己对圣人的爱德又有什么益处呢，如果他们本身在困苦中的话？}他就说：\Quote{我们喜乐，因为你们正在为自己在天上预备高贵的接待。}\Quote{因那望德，}他说，\Quote{那已存留的。}他显示这望德的稳妥。\Quote{你们从前在真理之言中听过。}这里的语气仿佛是在责备他们，好像他们长期持有之后竟改变了。\par

\Quote{关于这望德，}他说，\Quote{你们从前在福音真理之言中听过。}他为这真理作证。这是有充分理由的，因为在其中没有虚假。\par

\Quote{论到福音。}他并没有说\Quote{论到宣讲，}而是称它为\Quote{福音，}不断地提醒他们天主的恩惠。他先赞美他们，然后提醒他们这些恩惠。\par

第六节。\BibleQuote{这福音传到你们那里，如同在全世界上一样。} \BibleRef{哥 1:6}\par

他现在称赞他们。\Quote{传到}是比喻的说法。他的意思是，福音不是来了又去，而是存留并住在那里。又因为对许多人来说，教义最有力的确证就是他们与许多人共同持守，所以他加了\Quote{如同在全世界上一样。}\par

它处处临在，处处得胜，处处坚立。\par

\Quote{并结出果实，不断增长，如同在你们中间一样。}\par

\Quote{结出果实。}指的是在行为上。\Quote{增长。}指的是因许多人加入，变得更坚固；因为植物在变得坚固时才开始变粗。\par

他说：\Quote{如同在你们中间一样。}\par

他先以赞美赢得听者，好叫即使听者不情愿，也不会拒绝听他。\par

\Quote{自从你们听见的日子。}\par

何等奇妙！你们很快便归向它并相信了；并且立刻，从最初，就显出了它的果实。\par

\Quote{自从你们听见并真正认识天主的恩宠的日子。}\par

他说：不是只在言语上，也不是在欺骗中，而是在实际的行为中。那么，他所说的\Quote{结出果实}要么是指这个，要么是指神迹奇事。因为你们一领受它，就认得天主的恩宠。那么，那些立刻显明其内在德能的事物，如今竟然被怀疑，岂不是很令人难以接受吗？\par

第七节。\BibleQuote{正如你们从我们所亲爱的同仆\Person{厄帕夫辣}所学到的。} \BibleRef{哥 1:7}\par

他大概曾在那里宣讲过福音。\Quote{你们学到了}福音。然后为显示此人的可信，他说：\Quote{我们的同仆。}\par

\Quote{他为你们是基督忠信的仆役；他也向我们报告了你们在圣神内的爱德。}\par

他说：不要怀疑那将来的望德：你们看见世界正在被归化。何必引用别人的例子呢？发生在你们身上的事本身就是一个充分的相信基础，因为\Quote{你们真正认识了天主的恩宠，}也就是说，在行为上。所以这两件事，即众人的相信和你们自己的相信，都证实了将来的事。事实与\Person{厄帕夫辣}所说的并无不同。他说：\Quote{他是忠信的，}即真实的。怎么是\Quote{为你们的仆役}呢？因为他曾到他那里去。他说：\Quote{他也向我们报告了你们在圣神内的爱德，}即你们对我们所怀的属神的爱。如果这人是基督的仆役，你们怎能说你们是藉着天使接近天主呢？他说：\Quote{他也向我们报告了你们在圣神内的爱德。}因为这爱德是奇妙而稳固的；其他一切都空有其名。有些人不属于这一类，但这样的不是友谊，因此也容易瓦解。\par

产生友谊的原因有很多；我们要略过那些可耻的原因（因为没有人会为它们辩护，显然它们是邪恶的）。但如果你愿意，让我们反思那些自然的原因，以及那些从生活关系中产生的原因。属于社会性的原因如：一个人受人之惠，或从祖先继承一位朋友，或曾同席或同游；或与另一人为邻（这些是好的）；或同业，但最后一种不真诚，因为它伴随着某种竞争和嫉妒。而自然的原因如：父亲对儿子、儿子对父亲、兄弟对兄弟、祖父对后代、母亲对儿女，如果愿意，还可以加上妻子对丈夫；因为所有婚姻的纽带也属于今生，属于尘世。这些后者似乎比前者更强——我说\Quote{似乎，}是因为它们常常被前者超越。因为朋友有时会表现出比兄弟或儿子对父亲更真诚的善意；当亲生儿子不肯援助时，素不相识的人却会站在他身边援助他。但属神的爱高于一切，如同一位女王统治她的臣民；她的形态光辉灿烂：因为她不像其他爱那样，有任何尘世间的东西作为她的根源；既非习惯性的交往、恩惠、天性，也非时间；而是从天而降，出自天上。你为什么惊奇她不需要恩惠就能存在，既然连伤害也不能推翻她？\par

既然这爱比那爱更大，请听\Person{保禄}说：\Quote{我宁愿自己\OriginalTerm{受诅咒}{\GreekText{ἀνάθεμα}}，与\Person{基督}隔绝，为我的弟兄们。}\BibleRef{罗 9:3} 哪个作父亲的会这样希望自己受苦呢？又说：\Quote{去与\Person{基督}同在}是\Quote{更好得多的，然而留在肉身内}是\Quote{为你们更需要。}\BibleRef{斐 1:23} 哪个作母亲的会这样说话，不顾自己呢？再听他说：\Quote{因为你们暂时离别，是面离而心不离。}\BibleRef{得前 2:17} 在这里，确实在世上，当一位父亲受辱时，他会收回他的爱；但在那里却不是这样，他反而走向那些用石头砸他的人，设法对他们行善。因为没有什么，没有什么比圣神的纽带更坚固。因为那因受惠而成为朋友的人，一旦恩惠停止，就会成为仇敌；那因习惯交往而难舍难分的人，一旦习惯被打破，他的友谊就会消逝。同样，妻子若发生争吵，会离开丈夫，收回情感；儿子见父亲活到很大年纪，就会不满。但论到属神的爱，这些事一样也没有。因为没有任何东西能使它破碎，因为它不是由这些东西组成的。时间、长途旅行、恶待、毁谤、愤怒、侮辱，以及任何其他事，都不能侵袭它，也不能瓦解它。为使你们知道，\Person{梅瑟}曾被石头砸过，但他仍为他们恳求。\BibleRef{出 17:4} 哪个父亲会为砸他的人这样做，而不是也把他砸死呢？\par

那么，让我们追求这些属于圣神的友谊，因为它们坚固而难以解散，而不是那些从筵席产生的友谊，因为这些是我们被禁止带往那里的。请听\Person{基督}在福音中说：\Quote{你设宴时，不要请你的朋友或你的邻居，而要请瘸子、残废的。}\BibleRef{路 14:12} 有理由，因为对这些人的报答是大的。但你却不能，也不忍与瘸子、瞎子一同宴乐，反而认为它可厌可憎，并且拒绝。现在，你最好不拒绝，但也不一定要做。如果你不让他们与你同席，就差人把你自己桌上的菜肴送给他们。那请自己朋友的人，并没有做大事，因为他已在此得了赏报。但那请残废的和穷人的，却有天主做他的债主。那么，当我们在此不得赏报时，不要烦恼，但当我们就此得赏报时，却应烦恼；因为将来我们将一无所得。同样，如果人报答，天主就不报答；如果人不报答，天主要报答。那么，不要为我们的好处而寻求那些能回报我们的人，也不要存这样的期望施恩于他们：那是冷冰冰的想法。如果你请一位朋友，感激只持续到傍晚；因此，片刻的友谊比所付的开销消逝得更快。但如果你请穷人和残废的，感激绝不会消失，因为天主永远记念，永不忘记，你甚至有了祂做你的债主。请问，这是何等的小心眼，你竟不能与穷人同席？你说什么？他不洁净、污秽？那就给他洗，领他到你的桌前。但衣服污秽？那就给他换，给他洁净的衣服。你没有看见收益有多大吗？\Person{基督}藉着他来到你跟前，你却为这些琐事斤斤计较？当你邀请君王到你的桌前时，你因这些事而畏惧吗？\par

让我们设想两桌筵席：一桌坐满这些人，有瞎子、瘸子、手或脚残废的、赤脚的、只穿着一件破旧单衣的；另一桌则有显贵、将军、总督、高官，穿着昂贵的袍服和细麻衣，腰束金带。在穷人这桌，既没有银器，也没有丰盛的酒，只够使他们舒爽喜乐，酒杯和其余器皿都是玻璃制的；而在富人那桌，所有器皿都是金银的，半圆形的桌子，不是一个人能抬动的，而是两个壮汉都难以搬动的，酒罐排列整齐，黄金的光辉远胜银子，半桌上铺满柔软的帷幔。这里又有许多仆人，他们的衣着不比宾客逊色，穿着华丽的衣服和宽松的裤子，相貌俊美，正当青春年华，健壮丰润；而那里只有两个仆人，轻看一切骄傲的虚荣。一桌有昂贵的肴馔，另一桌只有足够充饥和振奋精神的。我说得够多了吗？两桌的摆设是否足够详尽？缺了什么吗？我想没有。因为我已提到了宾客、器皿的昂贵、亚麻布和肴馔。然而，如果我们遗漏了什么，我们会在继续讲论时发现。\par

那么，既然我们已经正确地按各自的轮廓画出了每张桌子，让我们看看你们将坐在哪一张。就我而言，我要去那瞎子和瘸子的桌子，但你们大多数人可能会选择另一张，就是将军们的，那桌子如此欢快华丽。那么让我们看看哪一张更充满快乐；因为目前我们还不检查将来的事，因为至少在那些方面，我的桌子有优越性。为什么？因为这一张有\Person{基督}坐席，另一张有人；这一张有主人，那一张有仆人。但我们暂时不提这些；让我们看看哪一张有更多的现世快乐。即使在这方面，我的快乐也更大，因为与君王坐席比与他的仆人坐席更快乐。但我们再排除这个考虑；让我们只检查事情本身。那么，我和那些选择我所选桌子的人，将极其自由和安心地说话和听一切；但你们战战兢兢、害怕，在你们同席的人面前羞愧，甚至不敢伸手，就好像你们去了学校而不是宴席，就好像你们在可怕的主人面前发抖。但他们不是这样。但是，有人说，尊荣很大。不，我更是尊荣；因为你们的卑下显得更宏大，当你们即使同席时，说出的话也是奴隶的话。\par

因为仆人最能显出他是仆人，就是当他与主人同坐时；因为他是在他不该在的地方；他并未因这样的亲密而获得多少尊严，反而得到卑下，因为他那时极其卑下。一个人可以看到仆人独自一人时显得勇敢，穷人独自一人时显得光彩，而不是当他与富人同行时；因为卑下者靠近崇高者时，就显得卑下，并置使卑下者显得更低，而不是更高。同样，你们与他们同坐也使你们显得更加卑下。但我们并非如此。那么，在这两件事上，我们有优势：自由和尊荣；这在快乐方面无可比拟。因为我至少宁愿要有自由的一块干面包，也不要有奴役下的千种美味。因为有人说：\BibleQuote{吃素菜而彼此相爱，强如吃肥牛而彼此相恨。}\BibleRef{箴 15:17} 因为无论那些人说什么，在场的人必须称赞它，否则就得罪人；因此他们充当了奉承者的角色，或者更糟。因为奉承者即使有羞辱和侮辱，仍有言论自由；但你们连这个也没有。但你们的卑下确实很大，（因为你们害怕和畏缩，）但你们的尊荣却不是这样。那么那张桌子确实被剥夺了一切快乐，但这一张充满了灵魂的一切喜悦。\par

但让我们进一步检查食物本身的性质。因为在那边，确实必须被迫喝大量酒，甚至违背自己的意愿；但在这里，任何不情愿的人都不必吃或喝。所以那边的食物品质带来的快乐被先前的羞辱和过饱带来的不适抵消了。因为过饱与饥饿一样摧毁折磨我们的身体，甚至更严重；随便你给谁，我都能更容易地用撑饱来毁灭他，而不是用饥饿。因为后者比前者更容易忍受，一个人可以忍受饥饿二十天，但过饱连两天也忍受不了。而且那些经常与饥饿斗争的乡下人是健康的，不需要医生；但另一种，我说的是过饱，没有人能不经常请医生就能忍受；甚至，它的暴政常常挫败了医生试图救援的努力。\par

那么，就快乐而言，我的这张桌子具有优势。因为如果尊荣比羞辱更有快乐，权柄比服从更有快乐，刚毅的信心比战栗和恐惧更有快乐，满足比淹没在奢侈的浪潮中更有快乐；那么在快乐方面，这张桌子比另一张更好。此外，在花费方面也更好；因为另一张昂贵，但这一张不是。\par

但怎样呢？这张桌子只是对客人更快乐，还是比另一张更给邀请者带来快乐呢？因为这是我们更想探究的。邀请那些人的主人许多天前就做准备，被迫有烦恼、焦虑和牵挂，夜晚不睡，白天不休息；而自己制定许多计划，与厨师、糖果师、摆桌师交谈。然后当那一天到来时，人们看到他比准备打拳击比赛的人更加害怕，担心任何事情出乎意料，担心被嫉妒的眼神击中，担心因此招来一群指责者。但另一位通过即兴准备桌子而避免了所有这些焦虑和麻烦，不在许多天前就操心。然后，确实，之后，前者立即失去了感恩的回报；但另一位有\Person{天主}作为他的债务人；并被美好的希望滋养，每天从那张桌子得享盛宴。因为食物确实被消耗了，但感恩的思念永远不会消耗，他每天比那些暴饮暴食的人更加欢喜快乐。因为没有什么比德行上的希望和美好事物的期望更能滋养灵魂。\par

但现在让我们思考接下来的事。那里确实有笛子、竖琴和风笛；但这里却没有不合适的音乐；有什么呢？有赞美的诗歌。在那里，确实是魔鬼被歌颂；但在这里，是万有的主、天主。你看到这一位是多么充满感恩，而那一位是多么忘恩负义和麻木不仁吗？因为，请告诉我，当天主用祂的美物养育了你，而你吃饱后本该感谢祂时，你甚至引进了魔鬼？因为这些里拉琴的歌曲，无非是献给魔鬼的歌曲。当你应该说：\Quote{主啊，你当受赞美，因为你用你的美物养育了我}时，你却像一条无用的狗，甚至不记得祂，反而引进了魔鬼？不，狗无论是否得到什么，都会向它们认识的人摇尾乞怜；但你连这也没有做到。狗即使什么也没得到，也向主人摇尾乞怜；但你即使得了恩惠，却向祂吠叫。再者，狗即使被陌生人善待，也不会因此消除对他的仇恨，不会被引诱与他为友；但你即使遭受魔鬼无数的祸害，却在你的宴席上引进他们。这样看来，你在两个方面比狗更坏。我刚才提到狗，对于那些只在得到恩惠时才感恩的人来说，是恰当的。我恳求你们，向那些即使在饥饿时仍向主人摇尾乞怜的狗学习羞耻吧。但你，如果碰巧听说魔鬼治好了某个人，就立刻抛弃了你的主人；啊，你比狗还要没有理性！\par

但有人会说：妓女看起来是一种快乐。那是怎样的快乐？不，那又是何等无耻！你的房子变成了妓院、疯狂和暴怒；你难道不羞于称之为快乐吗？如果允许使用她们，那么最大的快乐就是羞耻，以及因羞耻而产生的不安，使一个人的房子变成妓院，像猪在泥潭中打滚？但如果只允许看到她们，那么，看本身并不是快乐，因为不允许使用，反而只会增加欲望，使火焰更猛烈。\par

但你愿意知道结局吗？那些从餐桌起身的人，就像疯子或失去理智的人；鲁莽、好争吵、甚至成为奴隶的笑柄；仆人们确实清醒地退下，而他们却醉了。啊，真可耻！但与另一类人在一起，则没有这种事；他们以感恩结束宴席，然后这样回家，带着快乐入睡，带着快乐醒来，远离一切羞耻和谴责。\par

如果你也考虑客人本身，你会发现一类人在内里正像另一类人在外表一样：瞎眼、瘸腿、跛足；正如这些人的身体，那些人的灵魂也是如此，患着水肿和炎症。因为骄傲就是这样：在奢华的满足之后，残废就发生了；暴饮暴食和醉酒也是这样，使人跛足和残废。而且你也会看到，这些人拥有像那些人的身体一样光辉灿烂的灵魂。因为那些生活在感恩中、不追求超出所需之物、具有这种哲学的人，都在一切光明之中。\par

但让我们看看两边的结局。在那里，确实有淫荡的快乐、放荡的笑声、醉酒、粗俗的玩笑、污秽的言语；（因为既然他们自己羞于说污秽的话，就通过妓女来实现这一点；）但在这里，是爱人之心、温柔。站在邀请那些人之人的旁边的是虚荣，武装着他；但站在另一个人旁边的是爱人之心和温柔。因为一张桌子由爱人之心准备，另一张则由虚荣和残忍，出于不义和贪婪。那张桌子以我所说的告终：失去理智、谵妄、疯狂；（因为虚荣就是这样产生的；）但这张桌子则以感恩和天主的荣耀告终。而且来自人的赞美也更多地伴随着后者；因为前者甚至被人嫉妒，但后者被所有人视为共同的父亲，即使那些未从他手中得到任何好处的人也是如此。正如在受到伤害时，即使未受伤害的人也同情，所有人都成为共同的敌人（对伤害者）；同样，当一些人得到善待时，那些未得到任何好处的人也不亚于那些得到的人，赞美和钦佩施恩者。在那里，确实有许多嫉妒；但在这里，有许多温柔的关怀，许多来自众人的祈祷。\par

在今生如此之多；但到了那时，当基督来临时，这一位将带着极大的胆量站立，并在全世界面前听到：\BibleQuote{你看见我饿了，就给了我吃的；我赤身裸体，就给了我穿的；我做客旅，就收留了我}\BibleRef{玛 25:35}；以及所有类似的话；但另一位将听到相反的话：\BibleQuote{又恶又懒的仆人}\BibleRef{玛 25:26}；并且，再次：\BibleQuote{祸哉，那些在床榻上享乐，躺在象牙床上，喝精酿的酒，用上等的油膏抹自己的人；他们把这些事当作永恒的，而非短暂的。}\BibleRef{亚 6:4}\par

我说这话并非没有目的，而是为改变你们的心思，叫你们不要做任何无果之事。有人会说：\Quote{那么，我两者都做又如何呢？}这个论点所有人都常用。请告诉我，既然每件事都可以做得有益，何必分出一部分用在不需要甚至毫无目的的事上，另一部分才用在有益处的事上呢？告诉我，如果你撒种时，一些撒在石头上，一些撒在极好的地里，你难道会满足于此，说：\Quote{若有些撒在无用之地，有些撒在极好之地，又有何妨呢？}为何不全部撒在好地里呢？为何要减少收益呢？假如你要攒钱，你不会那样说，而是会从各处收集；但在另一件事上你却不如此。假如你要放债取利，你不会说：\Quote{我们为何要将一些给穷人，一些给富人？}而是全部给前者；然而在此处，收益如此之大，你却不这样计算，也不肯最终停止那无目的的消耗和无回报的支出。\par

有人说：\Quote{但这也有收益。}什么收益，请告诉我？\Quote{它增进友谊。}没有什么比因这些事——即筵席和暴食——而结交的朋友更冷淡的了。寄生虫的友谊只源于那里。\par

不要侮辱像爱德这样奇妙的事，也不要声称这是它的根源。就好像有人说，一棵结出金子与宝石的树，其根并非同类，而是由腐烂而生的；你正是如此。因为即使友谊由此而生，也没有比这更冷淡的了。但那另一张桌子所产生的友谊，不是与人的，而是与天主的；那是深厚的友谊，只要你们专心预备它。因为若有人一部分这样用，一部分那样用，即使他施舍了很多，也算不得什么大事；但若有人全部这样用，即使只给了少许，也成就了全部。因为我们被要求的不是施舍很多或很少，而是不要少于我们的能力。请想想那领了五个塔冷通的，和那领了两个的。\BibleRef{玛 25:15} 请想想那投入两个小钱的寡妇。\BibleRef{谷 12:41} 请想想\Person{厄里亚}时代的那位寡妇。那投入两个小钱的妇人并没有说：\Quote{我为自己留一个小钱，给另一个，有何妨呢？}而是把她全部的生活费都投了进去。\EditorialNote{列上 17} 而你，在如此丰裕之中，竟比她更吝啬。那么，我们不要对自己的救恩漠不关心，而要致力于施舍。因为没有什么比这更好的了，将来的时日会显明；同时，现在也显明了。为此，我们当为天主的荣耀而生活，做祂喜悦的事，以使我们堪当承受那预许的福分。愿我们众人，因我们主耶稣基督的恩宠和慈爱，获得这一切；愿光荣、权能和尊威归于祂，自今至永远，及世之世。阿们。\par

\SourceRecord{Translated by John A. Broadus.；Nicene and Post-Nicene Fathers, First Series；Vol. 13.；Edited by Philip Schaff.；Buffalo, NY: Christian Literature Publishing Co.,；1889.；<http://www.newadvent.org/fathers/230301.htm>.}

% structure: p0001 source=h1 target=section reason=page-title

\section{哥罗森书讲道 第二篇}

\BibleRef{哥 1:9-10}\par

\BibleQuote{为此，我们自从听见的那天起，也不断为你们祈祷，恳求你们充满对祂旨意的知识，藉着各种属神的智慧和明达，好使你们行事相称于主，事事中悦，在一切善工上结果实，并不断增加对天主的认识。}\par

\Quote{为此。}什么原因？因为我们听说了你们的信德和爱德，因为我们抱有美好的希望，我们也希望祈求未来的祝福。因为在比赛中，我们最鼓励那些接近胜利的人，同样，保禄也最劝勉那些已经完成了大部分的人。\par

\Quote{自从我们听见的那天起，}他说，\Quote{我们没有停止为你们祈祷。}不是一天两天，也不是三天。他在这里既显示了他的爱，又温和地暗示他们尚未到达终点。因为\Quote{使你们充满}这句话有这样的含义。请看我祈祷，这位蒙福者的谨慎。他从未说他们一无所有，而是说他们有所欠缺；到处都有\Quote{使你们充满}的话表明这一点。再者，\Quote{事事中悦，在一切善工上} \BibleRef{哥 1:11}，并且再次\Quote{加增一切力量，}又再次\Quote{以各种忍耐和宽容}；因为不断添加\Quote{一切}，证明他们部分做得不错，尽管可能不是全部。并且\Quote{使你们充满}，他说；不是\Quote{使你们领受，}因为他们已经领受了；而是\Quote{使你们充满}那尚欠缺的。这样，责备既无冒犯，赞美也不让他们消沉懈怠，好像已经完成了似的。但是\Quote{使你们充满对祂旨意的认识}是什么意思？就是藉着圣子我们被带到祂面前，不再藉着天使。现在你们必须被带到祂面前，你们已经学会，但还需要学习这一点，以及祂为何派遣圣子。因为如果我们本应被天使拯救，祂就不会派遣圣子，不会交出祂。\Quote{在一切属神的智慧中，}他说，\Quote{和明达。}因为哲学家们欺骗了他们；我希望你们，他说，在属神的智慧中，而不是按人的智慧。但是，如果为了认识天主的旨意，需要属神的智慧；那么认识祂的本质是什么，就需要不断的祈祷。\par

保禄在这里表明，从那时起他一直在祈祷，尚未成功，但也没有停止；因为\Quote{自从我们听见的那天起}这话就显示了这一点。但如果从那时起，即使有祈祷帮助，他们仍未改正自己，那对他们意味着谴责。他说，\Quote{并且恳求，}带着极大的热诚，因为\Quote{你们知道}这句话显示。但仍然有必要知道其他一些事。他说，\Quote{行事相称，}\Quote{于主。}他在这里谈论生命及其作为，因为他到处都这么做：他总是将信德与行为结合在一起。\Quote{事事中悦。}那么如何\Quote{事事中悦}呢？\Quote{在一切善工上结果实，并不断增加对天主的认识。}他说，既然祂已经将自己完全启示给你们，你们也领受了如此伟大的知识；那么你们就当显出行事相称于信德的行为；因为这需要高尚的行为，远超过旧约时代。因为那认识天主，并被算为配作天主仆人的，不，甚至祂的儿子，请看需要多么大的德行。\Quote{加增一切力量。}他在这里谈论考验和迫害。我们祈祷你们充满力量，不至于因忧愁而昏厥，也不致绝望。\Quote{依照祂荣耀的权能。}但你们要重新振作，如同祂荣耀的权能所给予的那样。\Quote{以各种忍耐和宽容。}他说的就是这些。简言之，他说，我们祈祷你们过一种有德行的生活，相称于你们的公民身份，并且站立得稳，藉着天主所赐的合理力量得以坚固。因此，他还没有触及教义，而是专注于生活，在这方面他没有什么可指责他们的，在应当称赞的地方称赞了他们之后，他于是转到了责备。而且他处处如此：当他将要写信给谁，既有可责备之处，又有可称赞之处时，他先称赞他们，然后再转到他的指责。因为他先抚慰听者，使他的指控免于任何怀疑，并表明就他自己而言，他本乐意全程称赞他们；但因情况所迫，不得不说出他所说的话。他在\WorkTitle{格林多前书}中也这样做。因为他在极为称赞他们爱他之后，甚至从淫乱者的事例，转过来指责他们。但在\WorkTitle{迦拉达书}中却不是这样，而是相反。然而，如果仔细看，即使在那里，指责也是在称赞之后。因为看到他们当时没有什么善行可说，而且指控非常严重，他们每人都败坏了；并且他们能承受，因为他们刚强，他先从指责开始，说：\Quote{我诧异。}\BibleRef{迦 1:6}所以这也是一种称赞。但后来他称赞他们，不是因为他们当时的状况，而是因为他们过去的情况，说：\Quote{若是可能，你们会挖出眼睛，给我。}\BibleRef{迦 5:15}\par

\Quote{结果实}，他说：这指向工作。\Quote{得坚固}：这指向试炼。\Quote{抵达一切的\OriginalTerm{忍耐}{patience}和\OriginalTerm{含忍}{longsuffering}}：对彼此含忍，对外人忍耐。因为含忍是对我们能够回报的人，而忍耐是对我们不能回报的人。因此，\OriginalTerm{忍耐}{patience}一词从未用于天主，但\OriginalTerm{含忍}{longsuffering}却常用；正如这位有福者在别处写道：\Quote{或者你轻视祂良善、宽恕和含忍的丰富吗？} \Quote{抵达一切所喜悦的。}不是一时这样，过后又不这样。\Quote{借着一切属神的智慧，}他说，\Quote{和悟性。}因为否则不可能知道祂的旨意。虽然他们以为自己知道祂的旨意，但那智慧并非属神的。\Quote{行走，}他说，\Quote{相称于主。}因为这是最佳生活之路。因为凡理解了天主对人的爱（他若看见圣子被交付，就理解了），就会有更大的进取心。此外，我们不仅为此祈祷，使你们知道，而是使你们在行为中显出知识；因为知道而不去做，甚至是在走向惩罚的路上。\Quote{行走，}他说，就是常常，不是一次，而是持续。如同行走对我们来说是必需的，同样正当生活也是必需的。当他谈论这个主题时，他经常使用\Quote{行走}一词，而且有理，表明这就是摆在我们面前的生活。但世俗的生活并非如此。赞美也是伟大的。\Quote{行走，}他说，\Quote{相称于主，}且\Quote{在一切善工中，}以便总是前进，绝不停止，并且用一个比喻说，\Quote{结果实并在认识天主上增长，}以便你们得蒙\Quote{坚固}，按照天主的德能，达到人所能达到的程度。\Quote{借由祂的德能，}安慰是大的。——他没有说力量，而是\Quote{德能}，这是更大的：\Quote{借由德能，}他说，\Quote{祂光荣的，}因为祂的光荣处处具有德能。他这样安慰那受责备的人：又说，\Quote{相称于主行走。}他说论到圣子，祂在天地处处都有权能，因为祂的光荣统治万有。他不仅说\Quote{坚固}，而且这样说，如同人们可以期望那些事奉如此强大主人的人那样。\Quote{在认识天主上。}同时他顺便触及认识的方法；因为这就是错误，没有按应当的方式认识天主；或者他的意思是，以便在认识天主上增长。因为不认识圣子的，也不认识圣父；恰当地需要增长的知识：因为没有这个，生活毫无用处。\Quote{抵达一切的忍耐和含忍，}他说，\Quote{怀着喜乐，感谢} \BibleRef{哥 1:12} 天主。然后，他正要劝勉他们，他没有提及将来为他们存留的；然而他在书信开头暗示了这一点，说：\Quote{因那为你们存留在天上的希望} \BibleRef{哥 1:5}：但在此处他提到他们已经拥有的事物，因为这些都是另一事物的原因。他在许多地方也这样做。因为已经发生的事令人信服，更能带动听众。\Quote{怀着喜乐，}他说，\Quote{感谢}天主。联系是这样的：我们不断为你们祈祷，并为已领受的恩惠而感谢。\par

\Quote{我们满怀喜乐地感谢}，这是一件大事。因为可能只是出于恐惧而感谢，也可能在忧伤中感谢。例如：约伯确实感谢，但是痛苦中；他说：\BibleQuote{主赐予，主取去。} \BibleRef{约 1:21} 因为，不要让任何人说所发生的事没有使他痛苦，也没有使他灵魂沮丧；也不要让那义人的伟大赞美被夺去。但当情况如此时，我们感谢不是因为恐惧，也不是因为祂是主，而是因为事物本身的性质。\Quote{是祂使我们相称，得以分享在光明中圣人的基业。}他说了一件大事。他说，所赐予的是这种性质；祂不仅赐予，还使我们强大到能够领受。现在通过说\Quote{祂使我们\OriginalTerm{相称}{meet}}，他表明这件事分量重大。例如，若某个卑微的人成为国王，他可以随意将管辖职位赐给他所愿意的人；他的能力仅限于赐予尊位：他不能使那人适合那个职位，而且常常尊荣会使被提拔者显得可笑。但如果他同时赐予尊位并使人适合那尊荣，足以担当管理，那么这事才真是尊荣。这就是他在这里所说的：祂不仅赐给我们尊荣，还使我们足够强大以领受它。\par

因为这里有双重的荣耀：赐予，并使人堪当这恩赐。祂没有简单地说赐予，而是\BibleQuote{使我们堪当与圣徒在光明中同得基业}，也就是说，祂为我们安排了与圣徒在一起的位置。祂不是简单地说安置了我们，而是使我们享受完全相同的事物，因为\Quote{分份}就是每人所得的部分。因为有可能在同一个城里，却不享受同样的事物；但拥有相同的\Quote{分份}却不享受相同的事物，是不可能的。有可能在同一个基业中，却没有同样的分份；例如，我们（司铎）都在基业中，但并非都有同样的分份。但这里他说的不是这个，而是连基业带分份一并说了。但他为什么称它为基业（或签分）？为了表明没有人凭自己的成就获得天国，而是如同签分是幸运的结果，这里也实在如此。因为没有人能显出一种如此美好的生活，足以被算为配得天国，而是完全出于祂的恩赐。因此祂说：\BibleQuote{你们做了一切之后，要说：我们是无用的仆人，因为我们做了该做的事。}\BibleRef{路 17:10}\BibleQuote{得以在光明中与圣徒同得基业}——他指的是将来的光明和现在的光明，也就是认识。在我看来，他同时谈论现在和将来。然后他显示了我们被算为配得什么。因为这并不是唯一的奇迹，即我们被算为配得天国；而且还当加上我们这些被如此计算的是谁；因为这并非不重要。他在\WorkTitle{罗马书}中这样做，说：\BibleQuote{为义人死，是少有的；为善人也许有人敢死。}\BibleRef{罗 5:7}\par

第十三节。\BibleQuote{祂拯救了我们}，他说，\BibleQuote{脱离黑暗的权势。}\BibleRef{哥 1:13}\par

一切都是出于祂，赐予这些和那些；因为没有任何事是我们的成就。\BibleQuote{脱离黑暗的权势}，他说，即脱离错谬，脱离魔鬼的统治。祂没有说\Quote{黑暗}，而是说\Quote{权势}；因为黑暗的权势对我们有巨大的能力，紧紧抓住我们。因为即使只是在魔鬼之下也是可悲的，但如此\Quote{在权势之下}则更为可悲。\BibleQuote{并且把我们移置}，他说，\BibleQuote{到祂钟爱之子的王国里}。祂不是仅仅为了将人从黑暗中拯救出来才显示了对他的爱。从黑暗中拯救出来确实是一件大事，但带进王国更是大得多。看这恩赐多么丰富：祂拯救了躺在深坑中的我们；其次，祂不仅拯救了我们，还把我们移置到了天国。\BibleQuote{祂拯救了我们}。祂没有说释放了我们，而是说\Quote{拯救}：表明我们的巨大痛苦，以及他们对我们的掳掠。然后，为了显示天主能力运作的轻易，他说，\BibleQuote{并且把我们移置}，犹如将一名士兵从一个阵地转移到另一个阵地。他没有说\Quote{带领过去}，也没有说\Quote{转移}，因为那样整个行动就全属转移者，而过去的人毫无作用；但他却说\Quote{移置}，这样既与我们有关，也与祂有关。\BibleQuote{进入祂钟爱之子的王国}。祂没有简单地说\Quote{天国}，而是通过说\Quote{圣子的王国}赋予其宏伟，因为没有比这更大的赞美了，正如他在别处也说：\BibleQuote{如果我们坚忍，我们也要与祂一同为王。}\BibleRef{弟后 2:12}祂算我们与圣子同得相同的事物；而且不仅如此，而且更有力的是，与祂的爱子同得。那些曾是仇敌、曾在黑暗中的人，仿佛瞬间被移置到圣子所在之处，与祂同享荣耀。祂并不满足于此，为了显示恩赐的伟大；祂不满足于说\Quote{王国}，还加了\Quote{圣子的}；不仅如此，祂还加了\Quote{钟爱的}；不仅如此，祂还加上了祂本性的尊贵。祂说什么？\BibleQuote{祂是不可见的天主的肖像。}但祂没有立刻说这话，而是插入了祂赐予我们的恩惠。因为免得你听见一切都是出于圣父，就以为圣子被排除在外，祂将一切归于圣子，也将一切归于圣父。因为确实是祂（圣父）移置了我们，但圣子提供了缘由。祂说什么？\BibleQuote{祂拯救了我们脱离黑暗的权势。}但相同的是，\BibleQuote{在祂内我们获得了完全的救赎，即罪过的赦免。}因为如果我们没有被赦免罪过，我们就不会被\Quote{移置}。所以这里又是\Quote{在祂内}这个词。祂没有说\Quote{救赎}，而是说\Quote{完全的救赎}，这样我们就不会再跌倒，也不会再被判死。\par

第十五节。\BibleQuote{祂是不可见的天主的肖像，是一切受造物的首生者。}\BibleRef{哥 1:15}\par

我们在这里遇到了一个异端问题。所以我们最好今天搁置，明天再继续讨论，在你们耳聪目明的时候讲给你们听。\par

但如果有人还要说些什么：圣子的工作更为伟大。如何呢？因为当人仍在罪中时，将天国赐予人是不可能的事；但这样却更容易，所以祂为恩赐预备了道路。你怎么说？祂亲自将你从罪中释放：那么祂也亲自使你接近；祂已经预先奠定了他教义的基础。\par

但我们必须先作一点补充，然后结束这篇讲论。那么，这补充是什么呢？既然我们已得以享如此大的恩惠，就应当常常记念它，不断在心中思量天主的恩赐，反省我们从什么中被解救出来，又得到了什么；这样，我们便会感恩，便会增进对祂的爱。人哪，你说什么？你被召叫进入一个王国，进入天主圣子的王国——而你却满口哈欠、搔首弄姿、昏昏欲睡吗？即使你每天必须跳入千万次死亡，难道你不该忍受一切吗？为了官职，你什么都做；那么，当你将要分享独生子的王国时，你难道不会扑向千万把刀剑？你难道不会跳进火里？但这还不算希奇；更希奇的是，当你甚至将要离开时，你却哀哭，宁愿留在这世上的事物中，做一个贪恋肉身的人。这是什么幻想？你甚至把死亡看作可怕的事？其原因在于奢侈和安逸：因为至少那些过着痛苦生活的人，会渴望翅膀，渴望从这里解脱。但现在，我们就像被宠坏的雏鸟，宁愿永远留在巢中。它们待得越久，就越虚弱。因为现世生命是一个用枝条和泥土粘合而成的巢。是的，即使你让我看到那些宏伟的宅邸，甚至用黄金和宝石闪闪发光的王宫，我也不会认为它们比燕子的巢更好；因为冬天一到，它们都会自行倒塌。我所指的冬天是那日子；并非它对所有人都是冬天。因为天主也称它为黑夜和白昼；黑夜指罪人，白昼指义人。所以现在我也称它为冬天。如果我们夏天没有被好好养育，以致冬天来临时不能飞翔，我们的母亲就不会接纳我们，而是让我们饿死，或当巢穴倒塌时死去；因为在那一天，天主会轻易地像拆毁一个巢——或者更轻易——拆毁一切，重新塑造。但那些羽毛未丰、不能在空中迎见祂、被如此粗鲁地养大以至于没有轻盈翅膀的，将会遭受在这种情形下应有之理的人所应受的苦。现在，燕子的雏燕跌落时迅速死亡；但我们不会死亡，而是永远受罚。那季节将是冬天；或者更确切地说，比冬天更严酷。因为那里滚下的不是冬天的洪流，而是火河；不是云中升起的黑暗，而是无法驱散的黑暗，没有一丝光线，以至于他们看不见天，也看不见空气，比那些被埋在地下的人更加困迫。\par

我们常常说这些事，但有些人我们无法使他们相信。然而，如果我们这些小人物在谈论这样的事时受到这样的待遇，并不奇怪，因为先知们也遇到过同样的事；他们不仅谈论这样的事，还谈论战争和掳掠。\BibleRef{耶 21:11} \Person{漆德克雅}受到\Person{耶肋米亚}的责备，却不觉得羞愧。因此先知们说：\BibleQuote{祸哉，那些说：让天主快快加速祂的工作，使我们看见，让以色列圣者的谋略来临，使我们知道的人。} \BibleRef{依 5:18} 我们不必对此感到惊奇。因为在方舟时代的人也不相信；但当他们的相信已无济于事时，他们相信了；索多玛人也没有预料到他们的命运，然而他们相信时，也毫无所得。我何必说将来的事呢？谁会预料到如今在各个地方发生的事：这些地震，这些城市的倾覆呢？然而，这些事比那些——我指的是方舟时代发生的事——更容易相信。\par

从何看出这一点呢？因为那时的人没有别的例子可以借鉴，也没有听到圣经；而另一方面，我们却有无数发生在我们自己年代和往年的事例。但这些人的不信是从何而来的呢？来自一个软化的灵魂；他们又喝又吃，因此不信。因为一个人所愿望的，他就认为并期待；而那些反对他的人，对他来说只是笑谈。\par

但愿我们不要如此；因为此后将不再有洪水，惩罚也不仅限于死亡；对于不信有审判的人，死亡将是惩罚的开始。有人问：谁从那边来这样说过？如果现在你这样说只是开玩笑，那也不妥当；因为这类事不可戏言；我们开玩笑，不是在当玩笑的地方，而是冒着危险；但如果你真是这样认为，以为身后一无所有，那你怎么还自称基督徒呢？因为我不考虑那些教外的人。你为何领受圣洗？为何踏入教会？难道我们应许你官职吗？我们的一切盼望都在将来的事。那么，如果你不信圣经，为何来呢？如果你不信基督，我不能称这样的人为基督徒；决不可以！他甚至比外邦人更差。在哪方面？在此：当你认为基督是天主时，你却不信祂如同天主。因为在那另一种不敬中至少有一贯性；认为基督不是天主的人，必然也不信祂；但这种不敬甚至没有一贯性：承认祂是天主，却不认为祂所说的话值得相信；这是醉话、奢靡、放荡的话。\Quote{我们吃喝罢，因为明天我们就要死了。}\BibleRef{格前 15:32} 不是明天；而是现在你已经死了，当你这样说时。那么，我们将与猪和驴毫无区别吗？请告诉我。因为如果没有审判，没有报应，没有法庭，为何我们蒙受理智这样的恩赐，万物都服在我们脚下？为何我们统治，而它们被统治？看魔鬼如何从各方面迫切地说服我们无视天主的恩赐。他像拐卖人口者和忘恩负义的仆人一样，把奴隶与主人混在一起；他竭力把自由人贬低到罪犯的层次。他表面上是推翻审判，实际上是在推翻天主的存有。\par

魔鬼的方式总是这样：他用诡诈而不直接的方式提出一切，好让我们不加提防。如果没有审判，天主就不公义（我按人的说法说）；如果天主不公义，那就根本没有天主；如果没有天主，一切都出于偶然，美德无价值，恶习也无价值。但他并不公开这样说。你明白这撒殚式论调的用意吗？他如何希望把我们变成畜牲，甚至野兽，甚至魔鬼？因此，我们不要被他说服。因为审判是存在的，你这悲惨可怜的人！我知道你为什么说出这样的话。你犯了许多罪，你得罪了，你没有信心，你以为事物的本性会随你的论点而改变。他说：在此期间，我不愿以地狱的期待折磨自己，如果真有地狱，我要说服自己它不存在；在此期间，我要在此享受奢靡！你为什么罪上加罪？如果你犯罪时相信有地狱，你只会为你自己的罪付代价；但如果你再加上这种不敬，你将为你的不敬和这种想法受最重的惩罚；而那对你而言是冷酷短暂的安慰，将成为你永远受罚的缘由。你犯了罪：就算这样吧：你为什么还鼓励别人犯罪，说没有地狱？你为什么误导单纯的人？为什么使百姓的手发软？就你而言，一切都颠倒了；好人不会变得更好，反而懈怠；恶人也不会停止作恶。因为，如果我们败坏别人，我们的罪会得到宽恕吗？你没有看见魔鬼如何试图拉亚当堕落吗？那对他有过宽恕吗？不，这反而成为更大惩罚的缘由，使他不仅为自己的罪，也为别人的罪受罚。因此，我们不要以为把别人拉入同样的毁灭会使审判座对我们更宽容。这反而会使它更严厉。我们为何自趋毁灭？这一切都来自撒殚。\par

人啊，你犯了罪吗？你有一位爱人的主。恳求、哀求、哭泣、呻吟；并恐吓他人，求他们不要陷入同样境地。如果在某家，一个得罪了主人的仆人对儿子说：\Quote{我的孩子，我得罪了主人，你要小心讨他喜欢，免得像我一样}：告诉我，他不会得到一些宽恕吗？他不会软化主人的心吗？但如果他抛开这种说法，反而说这样的话：主人不会按各人的行为报应；善恶都混在一起，不分彼此；在这家中没有感谢；你认为主人会对他怎么想？他不会因自己的恶行受更重的罚吗？理所当然；因为在前一种情况，他的情感会为他辩护，虽薄弱；但在后一种情况，无人会为他辩护。那么，即使没有别人，至少效法地狱中那个富翁，他说：\Quote{父亲亚巴郎，请派人到我兄弟那里，免得他们也来到这地方}，因为他自己不能去，好使他们不陷于同样的定罪。让我们停止这类撒殚式的话吧。\par

那么，他说，当希腊人向我们提出问题时，难道我们不该尝试医治他们吗？但你在医治希腊人藉口下，使基督徒陷入困惑，实则意在确立你撒殚的教理。因为当你独自与你的灵魂论及这些事时，你未能说服她；于是你便想引入他人作为证人。但若必须与希腊人辩论，不该从此开始；而应讨论基督是否为天主，是否为天主子；他们那些神祇是否是魔鬼。若这些要点得以确立，其余一切都随之而来；但在尚未打牢开端之前，就争论结局是徒劳的；在未学习首要原理之前，就讨论结论是多余而无益的。希腊人不信审判，他与你自己同处一种情况，因为他也有许多人在其哲学中论述过这些事；尽管他们如此谈论时，他们认为灵魂与身体分离，但他们仍设立了一个审判的宝座。这事如此清楚，几乎无人不知，甚至诗人们与所有人都同意，存在一个法庭和一个审判。因此，希腊人连他自己的权威也不信；而犹太人并不怀疑这些事，简言之无人怀疑。\par

那么我们为何自欺呢？看，你向我说这些话。你将向那\Quote{塑造我们每一颗心}\BibleRef{咏 32:15}、知晓心思中一切、那\Quote{是生活的，是有效的，比双刃剑还锐利}\BibleRef{希 4:12}的天主说什么呢？因为你请老实告诉我：难道你不定自己的罪吗？如此大的智慧——即犯罪者定自己的罪——岂是偶然发生的？因为这是伟大智慧之工。你定自己的罪。那赐予你这些思想的天主，会任由一切随意运行吗？因此，下列原则将是普遍而严格的：生活于德行中的人，没有一个完全不信审判的道理，即使他是希腊人或异端者。除非极少数生活于大罪恶中的人，无人接受复活之道。这正是圣咏作者所说：\Quote{你的判决已从他面前被夺去。}\BibleRef{咏 9:5}为何如此？因为\Quote{他的行径常是亵渎的}；因他说：\Quote{我们吃喝吧，因为明天我们就要死了。}\par

你看见吗？如此讲话正是卑俗之人的标记。吃喝产生这些颠覆复活的言论。因为灵魂无法忍受，我说，它无法忍受良心所提供的法庭，于是它便如同一个凶手，先暗示自己不会被察觉，然后就去杀人；因为若他的良心作他的法官，他就不至于仓促犯下那种大胆的恶行。他明知，却假装不知，免得被良心和恐惧折磨；因为确然，若他明知，他对那大胆的行为就不会如此坚决。同样，无疑地，那些犯罪、日复一日在同样邪恶中打滚的人，不愿知道它，虽然他们的良心在拉扯他们。\par

但我们不要理会这样的人，因为将有，必定有审判和复活，天主不会让如此伟大的工程失去指引。因此，我恳求你们，让我们离开邪恶，紧紧抓住德行，好使我们在我们的主耶稣基督内接受真道。然而，哪一样更容易接受？复活之道，还是命运之道？后者充满不义、荒谬、残忍、不人道；前者充满公义，按功过赏报；然而人们仍不接受它。但过错在于懈怠，因为没有有悟性的人接受那另一道。因为在希腊人中，那些接受那道理的人，是那些在享乐的定义中宣称它为\Quote{终向}的人，但那些热爱德行的人不接受它，反而把它当作荒谬加以排除。但若在希腊人中已是如此，那么在复活之道上更会成立。并且请观察，我祈求你，魔鬼是如何确立了两件相反的事：为了使我们忽视德行，并尊崇魔鬼，他引进了这种必然性，并借着两者他使两者都得到相信。那么，那固执地不信如此可赞美之事，并被那些如此妄谈的人说服的人，能提出什么理由呢？不要用你将获得宽恕的安慰来支撑自己；但要让我们聚集我们全部的力量，激励自己走向德行，真正活于天主，在我们的主耶稣基督内，等等。\par

\SourceRecord{Translated by John A. Broadus.；Nicene and Post-Nicene Fathers, First Series；Vol. 13.；Edited by Philip Schaff.；Buffalo, NY: Christian Literature Publishing Co.,；1889.；<http://www.newadvent.org/fathers/230302.htm>.}

% structure: p0001 source=h1 target=section reason=page-title

\section{哥罗森人书第三篇讲道}

\BibleRef{哥 1:15-18}\par

\BibleQuote{祂是不可见的天主的肖像，是万物的首生者；因为在天上和地上，看得见的和看不见的，或是宝座，或是宰制者，或是率领者，或是掌权者，一切都是祂内受造的：一切都是藉着祂，并且是为了祂而受造的；祂在万有之先，万有都赖祂而存在。祂是身体——教会的头。}\par

今天我必須清償昨日所延遲的債務，為的是當你們的心思處於飽滿狀態時，能對你們講論。\Person{保祿}在論及聖子的尊威時，說了這些話：\Quote{祂是不可見的天主的肖像。}那麼，你願意祂成為誰的肖像？天主的？那麼祂就與你所歸屬於祂的那位完全相似。因為如果祂是人的肖像，就這麼說吧，我就會當你是瘋子而離開你。但如果是天主、是天主子、是天主的肖像，祂就顯示出完全的相似。為什麼在任何地方都沒有天使被稱為\Quote{肖像}或\Quote{子}，而人卻兩者都是？為什麼？因為在前一種情形下，他們本性的崇高可能會立刻將許多人導入這種不敬；但在另一種情形下，卑下低微的本性就是對此的一種安全保障，不允許任何人——即使他們想要——懷疑這樣的事，也不允許將聖言貶低到如此地步。為此原因，在極度卑下的地方，聖經大膽地肯定這榮譽，但在本性較高的地方，它則克制。\Quote{不可見者的肖像}本身也是不可見的，而且是同樣不可見，否則它就不是肖像。因為肖像，就其為肖像而言，即使在我們中間，也應是完全相似的，例如在輪廓和相似方面。但在我們中間，這絕無可能；因為人的藝術在許多方面都失敗，或者說，如果你仔細檢查，它在所有方面都失敗。但在天主那裡，沒有錯誤，沒有失敗。\par

但如果是受造物，祂如何是造物主的肖像？因為馬也不是人的肖像。如果\Quote{肖像}不是指與不可見者的完全相似，那又有什麼妨礙天使也成為祂的肖像？因為他們也是不可見的，但不是彼此不可見；而靈魂也是不可見的：但因為它不可見，就僅僅因此是肖像，而不是像祂和天使那樣是肖像。\par

\Quote{一切受造物的首生者。}\Quote{那么，}有人说，\Quote{看哪，祂是一个受造物。}从何说起？请告诉我。\Quote{因为他说了\InnerQuote{首生者}。}然而，他说的不是\Quote{首先受造的}，而是\Quote{首生者}。那么，他应当被称为许多事物是合理的。因为他也必须被称为弟兄\Quote{在一切事上}。\BibleRef{希 2:17}我们难道必须剥夺祂的造物主身份，并坚持无论在尊威上还是其他任何事上，祂都不比我们优越吗？有谁理解力会这么说？因为\Quote{首生者}这个词并不表达尊威和尊荣，也不表达任何其他事物，而只表达时间。\Quote{首生者}是什么意思？回答说：祂是被造的。好吧。如果真是这样，它也有同类的表达。但另一方面，首生者与那些他所首生者具有同一性体。因此，祂将是万物的首生子——因为经上说\Quote{每一个受造物}；因此，天主圣言也是石头和我的首生者。但是，请再告诉我，\Quote{死者中的首生者}\BibleRef{哥 1:18}这些词表达什么？不是祂首先复活；因为他不只是简单地说\Quote{死者}，而是\Quote{死者中的首生者}，也不是说\Quote{祂首先死去}，而是祂作为死者中的首生者复活了。因此，这些词只表明一件事：祂是复活的初果。那么，在我们面前的段落中肯定也不是这样。接下来，他继续讨论教义本身。为了使他们不认为祂是较晚存在的（因为以前是通过天使接近，而现在是通过祂），他首先表明他们没有能力（否则祂就不会带他们\Quote{出于黑暗}\EditorialNote{第13节}），其次，祂也在他们之前。他使用这一点作为祂在他们之前的证据：万有都是由祂创造的。他说：\Quote{因为万有都是在祂内受造的。}\Person{撒摩撒塔的保禄}的跟随者在这里说什么？\Quote{天上的事物。}他将有疑问的事放在前面；\Quote{和地上的事物。}然后他说：\Quote{有形可见的与无形不可见的事物}；无形不可见的，如灵魂，以及在天上存在的一切；有形可见的，如人、太阳、天空。\Quote{无论宝座，}他放弃已经承认的事，而坚持有疑问的事。\Quote{无论宝座，或宰制者，或率领者，或掌权者。}\Quote{无论}、\Quote{或}这些词包含了所有事物；但通过更大的事物也表明较小的事物。但圣神并不在\Quote{掌权者}之列。他说：\Quote{万有都是藉着祂并为了祂而受造的。}看哪，\Quote{在祂内}就是\Quote{藉着祂}，因为说了\Quote{在祂内}之后，他加上了\Quote{藉着祂}。但\Quote{为了祂}是什么意思？这就是：万有的持存依赖于祂。不仅祂自己将他们从无中带出而存在，而且祂现在也亲自持守他们，以致如果他们与祂的眷顾分离，就会立刻毁灭和消亡。但祂没有说\Quote{祂持续他们}——那是一种更粗鲁的说法——而是更微妙的说法，即他们依赖\Quote{在祂上}。仅仅依靠祂就足以持续任何事物并紧紧维系它们。同样，\Quote{首生者}这个词，在根基的意义上也是如此。但这并不表明受造物与祂同体；而是说万有都是藉着祂并在祂内被支持。既然\Person{保禄}在别处也说：\Quote{我立定了根基}\BibleRef{格前 3:10}，他说的不是关于性体，而是关于行动。因为，为了不使你认为祂是一个仆役，他说祂持续它们，这并不比创造它们更少。确实，对我们来说这甚至更大：因为对于前者，手艺引导我们；但对于后者，则不是这样，它甚至不能防止事物腐朽。\par

\Quote{祂在万有之先，}他说。这适合天主。\Person{撒摩撒塔的保禄}在哪里？\Quote{万有在祂内得以存立，}即，它们被造成归于祂。他以紧密的顺序重复这些表达；以其紧密的连续，如同用快速的重击，将致命的教义连根拔起。因为，即使如此伟大的事已经被宣告，经过如此长的时间后\Person{撒摩撒塔的保禄}仍然出现，那么如果这些事没有提前被说，岂不更[会如此]？\Quote{万有在祂内，}他说，\Quote{得以存立。}如何在一个不存在者内\Quote{得以存立}？所以通过天使所做的事也是出于祂。\par

\Quote{祂是身体，教会的头。}\par

然后，在论及祂的尊威之后，他随后也谈到祂对人的爱。\Quote{祂是，}他说，\Quote{身体，教会的头。}他没有说\Quote{出于圆满}（虽然这也被表明），而是出于愿望显示祂对我们的巨大友谊，因为祂是如此崇高且超越一切，却与低微者相连。因为处处祂是首位：在上为首先；在教会为首先，因为祂是头；在复活中为首先。即，\par

\BibleRef{哥 1:18}\Quote{为使祂在一切事上居首位。}这样，在受生方面祂也是首先的。这正是\Person{保禄}主要竭力证明的。因为如果这一点成立，即祂在众天使之前；那么，随之而来的结论就是，祂通过命令他们而完成了他们的工作。而真正奇妙的是，他着重表明祂在后来的受生中也是首先的。虽然在别处他称亚当为首\BibleRef{格前 15:45}，事实上他确是如此；但在这里他把教会当作全人类。因为祂是教会的首；也是按肉身的人的首，如同是受造物的首一样。因此，他在这里使用了\Quote{首生者}一词。\par

此处所谓的\Quote{首生者}是何意指？是指那最先受造的，或是在万有之前复活的；正如前一处经文所指，是那在万有之先的。而且此处他确实用了\Quote{初果}一词，说：\Quote{祂是初果，是从死者中首生的，为叫祂在一切事上居首位}，表明其余的人也如祂一样；但前一处经文却不是受造的\Quote{初果}。那处经文说：\Quote{是不可见的天主的肖像}，然后才是\Quote{首生者}。\par

第19、20节：\Quote{因为父乐意叫一切的圆满居住在他内。藉着他十字架的血，缔造了和平，藉着他，无论是地上的，或是天上的，都与自己和好了。}\par

凡属于父的，他也说也属于圣子，而且更强烈地，因为祂不但为我们成了\Quote{死的}，而且将自己与我们结合。他说\Quote{初果}，如同果实的初果。他没有说\Quote{复活}，而是说\Quote{初果}，表明祂圣化了我们众人，并将我们如同祭品献上。\Quote{圆满}一词有人用以指天主性，如同若望所说：\Quote{从他的圆满中我们都领受了。}即是说，凡是圣子所是的，整个圣子就居住在那里，不是一种能力，而是一个实体。\par

他没有给出别的原因，只有天主的旨意：因为这就是\Quote{乐意…在他内。并…藉着他与自己和好一切}的含义。免得你以为祂只担任仆人的职务，他说\Quote{与自己}。\BibleRef{格后 5:18}然而他在别处又说，祂使我们与天主和好，就如他在写给格林多人的书信中所说的。他恰当地说\Quote{藉着他完成和好}；因为那时他们已经和好了；但他要说的是完全地、以这样的方式，使他们不再与祂为敌。如何呢？因为不仅和好被显明了，而且和好的方式也被显明了。\Quote{藉着祂十字架的血缔造了和平}。\Quote{和好}一词显明了敌意；\Quote{缔造了和平}一词显明了战争。\Quote{藉着祂十字架的血，藉着祂自己，无论是地上的，或是天上的。}和好实在是一件大事；但这件事是藉着祂自己完成的，则更大；还有更大的—如何藉着祂自己？藉着祂的血。祂的血，而且他不只说祂的血，而是更伟大的，藉着十字架。所以奇迹有五件：祂和好我们；与天主；藉着祂自己；藉着死亡；藉着十字架。令人赞叹！祂如何将它们混合起来！免得你以为只是一件事，或十字架本身算不得什么，他说\Quote{藉着祂自己}。他深知这何等伟大。因为不是藉着言语，而是藉着献出自己为和好，祂成就了一切。\par

但什么是\Quote{天上的}？有理由说\Quote{地上的}，因为那些充满了敌意，并且被多种方式分裂，我们每个人与自己、与众人全然不和；但祂如何使\Quote{天上的}和平呢？那里也有战争和打斗吗？那么我们如何祈祷说：\Quote{愿你的旨意奉行在人间，如同在天上}？\BibleRef{玛 6:10}那么这是什么意思呢？地曾与天分离，天使因看见主受辱而成了人的仇敌。\Quote{总之}，他说，\Quote{使万有都在基督内统一起来，无论是天上的，还是地上的。}\BibleRef{弗 1:10}如何呢？天上的方面是这样：祂将人迁移到那里，将仇敌、被憎恨者带到他们面前。祂不仅使地上的事物和平，而且将他们的仇敌和敌人带到他们面前。这里是深沉的和平。此后天使又在地上显现，因为人也出现在天上。依我看，保禄被提到天上也是为此\BibleRef{格后 12:2}，并表明圣子也已被接升到那里。因为在地上，和平是双重的：与天上之事，与自身；但在天上，和平是单纯的。因为如果天使为一个悔改的罪人而欢喜，那么为这么多罪人悔改，他们岂不更欢喜呢？\par

这一切都是天主的德能所成就的。那么你们为何还要依赖天使呢？他说。因为他们不仅没有使你们亲近，反而曾是你们的仇敌，除非天主亲自使你们与他们和好。那么你们为何投奔他们呢？你们要知道天使对我们的憎恨有多深，以及他们向来多么敌对吗？他们被派遣在以色列人、达味、索多玛人、哭泣山谷的事上施行惩罚。\BibleRef{出 23:20}但现在却不然，相反，他们以极大的喜悦在地上歌唱\BibleRef{撒下 24:16}。祂将天使带到人间\BibleRef{创 19:13}，又将人带到他们那里。\par

请看，我求你们，这里面的奇妙：祂先把这些带下到这里，然后把人提升到他们中间；大地成了天堂，因为天堂将要接纳地上的事物。因此，当我们感恩时，我们说：\BibleQuote{天主受享光荣于高天，平安归于世人，善意临于人。}看，他说，此后连人也蒙祂悦纳。\Quote{善意}是什么？\BibleRef{弗 2:14} 和好。天堂不再是隔墙。最初，天使的数目与万民相同；但现在，不是与万民数目相同，而是与信徒数目相同。这从哪里显明呢？请听基督说：\BibleQuote{你们小心，不要轻视这些小子中的一个，因为他们的天使在天上常常见到我父的面。}\BibleRef{玛 18:10} 因为每个信徒都有一位天使；从起初，每个蒙悦纳的人都有自己的天使，正如雅各伯说：\BibleQuote{那牧养我、救我脱离幼年一切灾难的天使。}\BibleRef{创 48:15} ，几乎如此。如果我们有天使，就让我们清醒，如同在导师面前；因为也有魔鬼在场。因此我们祈祷，祈求和平的天使，处处祈求和平（因为没有比这更珍贵的）；平安，在教会中，在祈祷中，在恳求中，在问候中；一次、两次、三次，许多次，那管理教会的人赐予它说：\BibleQuote{愿平安与你们同在。}为什么？因为这是万善之母，喜乐的基础。因此基督也命令宗徒们一进房屋就立刻说这话，作为福物的标记；因为祂说：\BibleQuote{你们进房屋时，要说：愿平安与你们同在。}因为缺少这一点，一切都无益。基督又对祂的门徒说：\BibleQuote{我把平安留给你们，我将我的平安赐给你们。}\BibleRef{若 14:27} 这为爱预备了道路。那管理教会的人不说简单的\BibleQuote{愿平安与你们同在}，而是说\BibleQuote{愿平安与众人同在}。因为如果我们与这人有平安，却与另一人有争战争斗，又有什么益处呢？在身体上，如果某些肢体安宁，另一些却失调，健康不可能维持；只有当全部肢体都处于良好秩序、和谐、平安，并且全部安宁，各守其限，否则一切都会颠覆。再者，在我们心中，除非所有思想都平安，否则平安不存在。平安是如此大的善，以至于缔造和平的人被称为天主的子女\BibleRef{玛 5:9}，有道理；因为天主子为此来到世上，为使地上的和天上的得以和平。但如果缔造和平的人是上主的子女，那么制造纷争的人就是魔鬼的子女。\par

你说什么？你挑动争斗和争战吗？有人问谁如此不幸？有许多人幸灾乐祸，他们撕裂基督的身体，比士兵用枪刺穿、或犹太人用钉子钉穿更为甚。那恶比这小；那些被刺穿的肢体重新结合了，但这一切被撕裂的，如果在这里不能结合，将永远不能结合，而与\OriginalTerm{圆满}{Fullness}分离开来。当你意欲与你的兄弟争战时，要想到你是在与基督的肢体争战，停止你的疯狂吧。若他是被遗弃的又如何？若他是卑贱的又如何？若他是可轻视的又如何？祂说：\BibleQuote{我父的旨意不是要这些小子中的一个丧亡。}\BibleRef{玛 18:14} 又说：\BibleQuote{他们的天使在天上常常见到我父的面。}\BibleRef{玛 18:10} 天主为他的缘故，也为你的缘故，甚至成了仆人，又被杀害；你竟视他为无物吗？确实，在这方面你也在与天主争战，因为你作出了与祂相反的判断。那管理教会的人进来时，立即说：\BibleQuote{愿平安与众人同在}；他讲道时说：\BibleQuote{愿平安与众人同在}；他祝福时说：\BibleQuote{愿平安与众人同在}；他请人问候时说：\BibleQuote{愿平安与众人同在}；祭献结束时说：\BibleQuote{愿平安与众人同在}；有时在中间说：\BibleQuote{愿恩宠与平安赐予你们。}既然我们多次听到应当拥有平安，却彼此处于敌对状态；领受平安，又回赐平安，却与赐平安的人争战，这岂不荒谬吗？你说：\BibleQuote{也与你的心灵同在。}你却在外面诋毁他吗？我有祸了！教会庄严的礼仪竟成了仅仅的形式，而非真实。我有祸了！这军队的口号只停留在言语上。因此你们也不知道为何说\BibleQuote{愿平安与众人同在}。但请听接下来基督所说的话：\BibleQuote{你们不论进了哪一座城或哪一个村庄……进了房屋，要请安；如果这家是堪当的，你们的平安就降到它身上；如果是不堪当的，你们的平安仍归于你们。}\BibleRef{玛 10:11} 因此我们无知，因为我们只看成是言词的修辞，心里并不认同。因为我赐平安吗？是基督屈尊借我们说话。即使在其它时候我们没有恩宠，但现在为了你们，我们并非没有。因为如果天主的恩宠曾在驴和占卜者身上运作，为了一个\OriginalTerm{救恩计划}{economy}，并为以色列人的益处\EditorialNote{户籍纪22章}，那么显然，祂也不会拒绝在我们身上运作，但为了你们，祂会忍受这一切。\par

但愿没有人说我是卑微、下贱、不配被重视，并以这样的心态对待我。因为我确实如此；但天主的行事方式始终是：为了众人的缘故，甚至与这样的人同在。为使你们知道这事，祂为亚伯尔的缘故俯允与加音说话\EditorialNote{创世纪4}，为约伯的缘故与魔鬼说话\EditorialNote{约伯传1}，为若瑟的缘故与法郎说话\EditorialNote{创世纪41}，为达尼尔的缘故与拿步高说话\EditorialNote{达尼尔2和达尼尔4}，为同一人的缘故与贝耳沙匝说话\EditorialNote{达尼尔5}。而且贤士们获得了启示\EditorialNote{玛窦福音2}；盖法虽为杀害基督者，不配的人，却因司祭职的尊贵而说了预言\BibleRef{若 11:49}。据说正是为此，亚郎没有患癞病。为何，请告诉我，当两人都说了反对梅瑟的话，却只有她受了惩罚？\EditorialNote{户籍纪12}不要惊奇：因为在世俗的职位上，即使有千万指控加于一人，在他卸任之前也不会受审，为的是不让职位与他一同受辱；何况在神职上，无论他是谁，天主的恩宠在他内运作，否则一切都丧失了；但当他卸任后，或离世后，甚至在此世，那时，他必受更重的惩罚。\par

我恳求你们，不要以为这些话是出于我们；而是天主的恩宠运作在不配的人身上，不是为了我们，而是为了你们。那么请听基督所说的话。\Quote{如果这家堪当，你们的平安就临到它。}\BibleRef{玛 10:13} 它如何成为堪当的呢？如果\Quote{他们接待你们}\BibleRef{路 10:8}，祂说。\Quote{但若他们不接待你们，也不听你们的话，……我实在告诉你们：在审判的日子，索多玛和哈摩辣地所受的惩罚，比那城还要容易忍受。}那么，如果你们接待我们，却不听我们所说的话，有什么益处呢？如果你们服事我们，却不留心所讲给你们的话，有什么收获呢？这将是我们的荣耀，这可钦崇的服事，如果你们听我们，对你们和我们都有益处。也听保禄所说：\Quote{弟兄们，我不认识他是大司祭。}\BibleRef{宗 23:5} 也听基督所说：\Quote{凡他们所吩咐你们遵守的}\BibleRef{玛 23:3}，你们要\Quote{遵守并实行。}你们所轻视的不是我，而是司祭职；当你们看见我脱去这职务时，那时再轻视我吧；那时我不再勉强发号施令。但只要我们坐在这座位上，只要我们居首位，我们就拥有尊威和权能，即使我们不配。如果梅瑟的座位如此可敬，以致为了它他们该被听从，何况基督的座位呢？我们由继承得到了它；从它我们发言；自从基督将修好的职务托付给我们以来。\par

无论怎样的使节，因为使节的尊严，都享有极大的荣耀。因为请看：他们独自进入蛮邦腹地，穿过如此众多的敌人；由于使节的法律具有强大力量，所有人都尊敬他们，都以敬意注视他们，都安全地护送他们离去。我们现在也领受了一个使节的话语，我们是从天主来的，因为这就是主教职的尊严。我们来到你们这里，作为使节，请求你们结束战争，我们讲明条件；不是应许给城市，也不是若干斗的谷物，也不是奴隶，也不是金子；而是天国、永生、与基督的共融，以及其他美善，只要我们还在这肉身和现世生命中，我们甚至无法告诉你们。所以我们作为使节，希望享有荣耀，不是为了我们自己——绝无此意，因为我们知道它的虚无——而是为了你们；使你们能热切地听我们所说的话，使你们受益，使你们不是漫不经心或漠不关心地注意所讲的话。你们难道不见使节们，所有人都如何向他们致敬吗？我们是天主对人的使节；但如果这冒犯你们，不是我们，而是主教职本身，不是这人或那人，而是主教。愿没有人听我，只听这尊严。那么，让我们事事遵行天主的旨意，使我们能生活于天主的荣耀，并堪当领受那为爱祂的人所预许的美善，藉着恩宠和慈爱，等等。\par

\SourceRecord{Translated by John A. Broadus.；Nicene and Post-Nicene Fathers, First Series；Vol. 13.；Edited by Philip Schaff.；Buffalo, NY: Christian Literature Publishing Co.,；1889.；<http://www.newadvent.org/fathers/230303.htm>.}

% structure: p0001 source=h1 target=section reason=page-title

\section{\WorkTitle{哥罗森书讲道词 第四篇}}

\BibleRef{哥 1:21,22}\par

\BibleQuote{\Emph{你们从前是仇敌，被隔绝}\Emph{在你们的心意上，在你们的恶行上，如今他却藉着他肉身的死亡使你们与他和好，好使你们在他面前成为圣洁、毫无瑕疵、无可指摘。}}\par

保禄在这里要表明，基督甚至与那些不配和好的人和好了。因为他说他们曾在\Quote{黑暗的权势}之下，由此指明他们所遭受的灾祸。\BibleRef{哥 1:13} 但恐怕你听到\Quote{黑暗的权势}时，会认为那是必然性，他便补充说\Quote{你们这些过去被隔绝的}，这样，虽然他所说的似乎是同一件事，但事实并非如此；因为将因必然性而受苦的人从邪恶中拯救出来，与将自愿忍受的人拯救出来，并不是同一回事。前者固然值得同情，后者却令人憎恶。然而，他说，你们这些并非出于勉强或强迫，而是出于自己的意愿和欲望背离了祂，且不配和好的人，基督却使你们与祂和好了。并且，因为保禄提到了\Quote{天上的事物}，他便表明，一切敌意都起源于这里，而非那里。因为他们早已渴望和好，天主也渴望，但你们不愿意。\par

并且，他在整个过程中表明，在随后的时期里，天使们没有能力，因为人们仍然是仇敌；他们既不能说服他们，即使说服了，也不能将他们从魔鬼手中解救出来。因为说服他们没有任何益处，除非那捆绑他们的被束缚；捆绑那捆绑者的也没有用，除非被拘留者愿意返回。但这两者都是必需的，而天使们一样也做不到，但基督却做到了这两样。因此，甚至比解除死亡更奇妙的，是说服他们。因为前者完全出于祂自身，能力完全在祂自己内，但后者不仅在于祂自身，也在于我们；但我们更容易完成那些能力在于我们自己的事。因此，因为更伟大，他把说服放在最后。而且，他不仅仅说\Quote{曾为仇敌}，而是说\Quote{被隔绝}，这表示极大的敌意，还不只是\Quote{被隔绝}，而是甚至没有任何返回的希望。他又说\Quote{在你们的心意上为仇敌}；那么，隔绝还不止于意图——而是什么？还有\Quote{在你们的恶行上}。他说，你们既是仇敌，也行仇敌的事。\par

\BibleQuote{如今他却藉着他肉身的死亡使你们与他和好，好使你们在他面前成为圣洁、毫无瑕疵、无可指摘。} 接着，他又论及和好的方式，即\Quote{在肉身中}，不是仅仅被鞭打、被鞭笞、被出卖，而是甚至死于最耻辱的死亡。他又提到十字架，又论及另一项恩惠。因为祂不仅\Quote{拯救}，而且，如他上文所说，\Quote{使我们有资格}\BibleRef{哥 1:12}，他在这里也暗指同一件事。他说，\Quote{藉着他的死亡}\Quote{使你们在他面前成为圣洁、毫无瑕疵、无可指摘。} 因为确实，祂不仅从罪恶中拯救了我们，而且还把我们列在蒙悦纳的人中。因为祂遭受如此大的苦难，不仅是为救我们脱离邪恶，而且也是为使我们获得首奖；就如一个人不仅使被定罪的犯人免于处罚，而且还提升他到尊荣的地位一样。并且，祂将你们与那些没有犯罪的人并列，是的，不仅与没有犯罪的人并列，甚至与行了最大义德的人并列；并且，真正伟大的是，祂赐予了在祂面前的圣洁和无可指摘。现在，无可指责比无瑕疵更进一步，就是当我们没有可被定罪或控告的事的时候。但是，既然他将一切都归于祂，因为藉着祂的死亡，祂成就了这些事；\Quote{那么，有人会说，这与我们有什么关系？我们什么都不需要。} 因此他补充说：\par

\BibleRef{哥 1:23} \BibleQuote{只要你们在信德上根基稳固，坚定不移，不致偏离福音的希望。}\par

这里他打击他们的懈怠。他并不仅仅说\Quote{继续}，因为继续也可以摇摆不定、犹豫不决；可以站着并继续，虽然东倒西歪。他说，\Quote{只要你们继续}，\Quote{根基稳固，坚定不移，不致偏离。} 奇妙！他使用了多么有力的比喻；他不仅说不要被抛来抛去，甚至说不要移动。并且注意，到目前为止，他没有提出任何沉重或辛苦的要求，而是信德和望德；也就是说，如果你们继续相信，相信未来之事的希望是真实的。因为这确实是可能的；但是，就德行生活而言，即使只是轻微地，也不可能不被摇动；所以，他所要求的并不沉重。\par

他说，\Quote{从福音的希望，就是你们所听过，且传遍了天下受造物的。} 但福音的希望除了基督之外，还有什么呢？因为祂自己就是我们的和平，成就了这一切；所以，将这些归给其他人的人，就是\Quote{偏离了}；因为他失去了一切，除非他相信基督。他说，\Quote{你们所听过的。} 他又以他们自己为证人，然后以整个世界为证人。他不是说\Quote{正在被传扬的}，而是已经被人相信并传扬了。正如他在开始时也做过的那样\BibleRef{哥 1:6}，希望通过许多人的见证来坚固这些。\Quote{我保禄做了这福音的仆役。} 这也使得它可信；他说，\Quote{我，保禄，做了仆役。} 因为他的权威很大，如今处处被颂扬，且是世界的导师。\par

\BibleRef{哥 1:24} \BibleQuote{现在我为你们受苦而喜乐，并在我肉身上补满基督苦难所欠缺的，为了他的身体，即教会。}\par

这与什么有关？似乎确实没有关联，但其实密切相关。他说：\Quote{\InnerQuote{仆役}}，意思是不从我自身带来什么，而是宣布来自他人的事。我如此相信，以致我甚至为祂受苦，不仅受苦，而且在苦难中喜乐，仰望那未来的希望；我受苦不是为我自己，而是为你们。他又说：\Quote{\InnerQuote{在我身上补满基督苦难所欠缺的}}。这似乎是一句伟大的话，但绝非出于傲慢，恰恰相反，是出于对基督极大的挚爱；因为他不想将苦难归于自己，而是归于祂，以渴望使这些人归向祂。他说：我所受的一切苦都是为祂的缘故；因此，你们不必感谢我，而要感谢祂，因为是祂自己受苦。就像一个人被派往某人处，向另一人请求说：\Quote{我求你，替我去见这人}；然后那人说：\Quote{我是为他而做的}。所以，祂不耻于将这些苦难也算作祂自己的。因为祂不仅为我们而死，甚至在死后，祂仍准备为你们而受苦。祂热切而强烈地表明，为了教会的缘故，祂现在仍在自己的身体中身处险境；他旨在说明这一点：你们不是被我们带到天主面前，而是被祂带到天主面前，即使我们做了这些事，因为我们所从事的不是我们自己的工作，而是祂的工作。这就像有一支军队，有指定的指挥官保护它，它应该站在战场上；当指挥官离开后，他的副手应接替他受伤，直到战斗结束。\par

其次，请听他是如何为祂的缘故做这些事的：他说：\Quote{\InnerQuote{为祂的身体}}，当然意思是说：\Quote{\InnerQuote{我不是取悦你们，而是基督；因为祂本该受的苦，我替祂受了}}。请看，他确立了多么多的要点。他表明他们对他有极大的爱债。正如他在致格林多人的第二封信中所写：\Quote{\InnerQuote{祂将和好的职务赐给了我们}} \BibleRef{格后 5:20}；又说：\Quote{\InnerQuote{我们是代基督作大使，好像天主藉着我们来劝勉你们}}。同样地，他在这里也说：\Quote{\InnerQuote{我为祂受苦}}，以便更将他们吸引到祂那里。也就是说，虽然你们欠债的那位已经离开了，但我来偿还。为此他也说：\Quote{\InnerQuote{所欠缺的}}，以表明他并不认为祂已经受尽了所有的苦。他说：\Quote{\InnerQuote{为你们}}，即使在祂死后祂仍受苦；因为仍有欠缺。他在致罗马人的书信中用了另一种方式做同样的事，说：\Quote{\InnerQuote{祂也为我们转求}} \BibleRef{罗 8:34}，表明祂不满足于祂的死亡，而且在之后仍做无数的事。\par

他这样说不是为了抬高自己，而是渴望表明基督仍然关心他们。他加上\Quote{\InnerQuote{为祂的身体}}来证明他所说的是可信的。因为事实如此，并无不合情理之处，从这些事是为祂的身体而行便可明瞭。请看祂如何将我们与祂自己紧密相连。那么，为什么要在中间引入天使呢？他说：\Quote{\InnerQuote{我作了这福音的仆役}}。为什么还要引入别的天使？\Quote{\InnerQuote{我是仆役}}。然后他表明他自己并没有做什么，尽管他是仆役。他说：\Quote{\InnerQuote{我作了这福音的仆役，是按照天主赐给我为你们而有的\OriginalTerm{上智安排}{dispensation}，以完成天主的圣言}}。\Quote{\InnerQuote{\OriginalTerm{上智安排}{dispensation}}}。他或者意思是：祂愿意在自己离开后，由我们接续这上智安排，免得你们感觉被遗弃（因为是祂自己受苦，是祂自己作大使）；或者他意思是：祂容许我这个曾是迫害者的人（比任何人都更是）受迫害，好使我的宣讲得以令人相信；或者他藉\Quote{\InnerQuote{\OriginalTerm{上智安排}{dispensation}}}意思是：祂不要求行为、行动或善工，只要求信德和圣洗。因为否则你们就不会接受这话。他说：\Quote{\InnerQuote{为你们，以完成天主的圣言}}。他提到外邦人，用\Quote{\InnerQuote{完成}}一词表明他们仍在摇摆不定。因为被遗弃的外邦人能够接受如此高深的道理，不是出于保禄，而是出于天主的\OriginalTerm{上智安排}{dispensation}；他说：\Quote{\InnerQuote{因为我绝不可能有这样的能力}}。在表明那更伟大的事（即他的苦难是基督的）之后，他接着加上更明显的事：这也是出于天主，\Quote{\InnerQuote{在你们内完成祂的圣言}}。他在这里含蓄地指出，这也是出于上智安排，即现在当你们能听的时候，这话对你们说，并非由于疏忽，而是为使你们接受它。因为天主不是突然做一切事，而是因祂对人丰厚的爱而屈就。这就是为何基督在这时候来临，而非在古时。祂在福音中表明，为此祂先派了仆人，以免他们立刻杀死圣子。因为如果他们在仆人之后仍不尊敬圣子，那么祂若更早来，他们岂不更不尊敬？如果他们不听从较小的诫命，怎能听从更大的？那么，有人会反对：难道现在不是仍有犹太人和希腊人处于极不完善的状态吗？然而，这是极端懈怠的证明。因为经过这么长的时间，受了这么大的训诲，仍然停留在不完善的状态，是极大的愚钝。\par

当希腊人说，为什么基督在这时候来临？我们不要让他们这样说，而要问他们，难道祂没有成功吗？因为，假如祂在最初就来临，却没有成功，那么时间就不会成为我们充分的借口；同样，既然祂已经成功了，我们就不能正当地因\Quote{时间}的缘故而受追究。因为没有人会要求一位消除了疾病、使人恢复健康的医生说明他的治疗方法，也没有人会仔细盘问一位取得胜利的将军为什么在这个时候、为什么在这个地方。这些事，如果他没有成功，问一问倒也合适；但他既然成功了，就必须视为理所当然。因为，请告诉我，你的推理和诽谤，或是事物的完善，哪个更值得相信？祂得胜了还是没有得胜？请指出。祂占了上风还是没有占上风？祂成就了祂所说的还是没有？这些是探究的要点。请告诉我，我恳求你。你完全承认天主存在，即使不是基督？那么我问你：天主是无始的吗？你会说，当然。那么请告诉我，祂为什么不在无数年前造人？因为那样他们就能活更长时间。他们现在因那段不存在的时间而蒙受了损失。不，他们没有损失；但如何，唯有创造他们的那一位知道。再者，我问你，祂为什么不一下子造出所有的人？但是，无论谁是最初被造的，他的灵魂已存在了那么多年，而那尚未被造的人却被剥夺了这些岁月。为什么祂使这一个先被带到这个世界，而另一个后来才被带来？\par

虽然这些事情确实是适合探究的主题，但并非出于多管闲事的好奇心：因为这根本就不是探究。因为我要告诉你们我为何这样说。因为假设人性仿佛是一个连续的生命，早期我们人类处于童年；随后处于成年；而现在临近晚年，则处于老年。如今，当灵魂达至圆满，当身体的四肢松弛，我们的战争结束时，我们便被引向哲学。相反，有人会说，我们在幼年时就教导男孩。是的，但并非伟大的教义，而是修辞学和语言的熟练；其他的则在他们成熟时才教导。请看天主对犹太人也是如此行的。因为正如犹太人曾是小孩子，天主便立梅瑟作他们的\OriginalTerm{启蒙导师}{schoolmaster}，并如同对待小孩子一般，通过\OriginalTerm{预象}{shadowy representations}为他们管理这些事，就像我们教写字一样。\BibleQuote{因为法律只有未来美物的影子，没有事物的真相。} \BibleRef{希 10:1}如同我们既为孩子们买糕点，又给他们零钱，只要求他们一件事：目前能去上学；同样，天主在那时也赐予他们财富和奢侈，藉此祂的大度宽容只向他们换取一件事：就是他们会听从梅瑟。因此，祂把他们交托给一位\OriginalTerm{启蒙导师}{schoolmaster}，免得他们因祂是一位温柔慈爱的父亲而藐视祂。请看，那时他们只惧怕梅瑟；因为他们不说：天主在哪里？而是说：梅瑟在哪里？而且他的出现本身就令人畏惧。所以当他们做错时，请看他是如何惩罚他们的。因为天主确实想抛弃他们，但梅瑟不允许祂这样做。或者不如说，整个都是出于天主；正如一位父亲在威胁，而启蒙导师向祂恳求说：\Quote{求祢因我的缘故宽恕他们，今后我为他们担保。}这样，旷野便成了一所学校。正如那些上学已久的孩子渴望离开学校，那时他们也不断渴望埃及，哭泣着说：\BibleQuote{我们完了，我们完全被消灭，我们彻底完了！} \BibleRef{出 16:3}梅瑟就摔碎了他们的\OriginalTerm{石板}{tablet}，因为他为他们写了一些话，\BibleRef{出 32:19}；正如一位启蒙导师所做的：他拿起石板，发现写得很差，就把石板扔掉，想表示极大的愤怒；如果他摔碎了石板，父亲也不会生气。因为他确实忙着写字，但他们却未注意他，反而转向别处，制造混乱。如同在学校里，他们互相击打，同样，在那场合，他命令他们互相击打并杀戮。此外，他好像给了他们功课学习，然后查问，发现他们没有学会，就惩罚他们。例如。那些表明天主能力的文字是什么？在埃及的事件？是的，有人说，但这些文字代表的是灾祸，即祂惩罚祂的仇敌。对他们来说，这是一所学校。惩罚你们的仇敌除了是你们的益处之外，还能是什么？在其他方面，祂也使你们受益。这就像一个人说自己认识字母，但被反复询问时，却出错而被责打。同样，他们确实说他们知道天主的能力，但当反复询问他们的知识时，他们答不出来，因此被责打。你们看见过水吗？你们应当想起埃及的水。因为那将水变成血的，也必有能力这样做。正如我们也常对孩子们说：\Quote{当你在一本书中看到字母A时，要记得你在石板上也学过它。}你们见过饥荒吗？要记得是祂毁坏了庄稼！你们见过战争吗？要记得那次淹溺！你们见过那地居民强大吗？但他们并不比埃及人强大。那将你们从他们中间领出来的，岂不更能拯救你们出来吗？但他们不知道如何按顺序回答他们的字母，因此被责打。\BibleQuote{他们吃了，喝了，就踢跳。} \BibleRef{申 32:15}当他们以玛纳为食时，本不应寻求奢华，因为他们已知道由此产生的恶果。他们的行为恰如一个自由的孩子，被送去上学时，却要求与奴隶为伍，并伺候他们——他们寻求埃及也是如此——当获得一切必需的、合乎自由人身份的给养，坐在父亲的桌前时，却渴望那气味难闻、吵吵闹闹的仆人的饭食。他们对梅瑟说：\BibleQuote{是的，主，凡你所说的，我们都必遵行，并听从。} \BibleRef{出 24:7}正如极其顽劣的孩子的情况：当父亲要处死他们时，启蒙导师坚持不懈地为他们求情，那时的情况也是如此。\par

为什么我们说了这些？因为我们与孩童毫无差别。你要听听他们的道理，也是孩童的道理吗？\BibleQuote{以眼还眼}，经上说，\BibleQuote{以牙还牙}。\BibleRef{肋 24:20} 因为没有比孩童的心更急于报复的了。因为看到它是一种非理性的激情，在那年龄段有许多非理性和极大的欠缺考虑，无怪乎孩童被愤怒所支配；而且这支配如此之大，常常他们绊倒后又站起来，会因愤怒而捶打自己的膝盖，或踢翻脚凳，以此来减轻痛苦，平息暴怒。天主也以类似的方式对待他们，祂允许他们实行\BibleQuote{以眼还眼，以牙还牙}，并毁灭了曾使他们忧伤的埃及人和阿玛肋克人。祂还应许了这样的事；就像有人对父亲说：\Quote{父亲，某某人打了我}，父亲就回答说：\Quote{某某人是个坏人，我们恨他。} 同样，天主也说：\BibleQuote{我要作你们敌人的敌人，恨你们的人的仇人。}\BibleRef{出 23:22} 再者，当巴郎祈祷时，天主对他们的迁就也是孩童式的。因为如同对待孩童，当他们被某件并不可怕的东西吓到时，比如一团羊毛或其他类似的东西，他们会突然惊慌；为了不让恐惧持续，我们就把那东西拿到他们手上，让保姆给他们看：天主也是如此；祂见那先知成了他们的恐惧，就把对先知的恐惧化为信心。又如同断奶的孩童，篮子里有各色各样的东西，祂也赐给他们一切，大量的美味。但孩童仍渴望乳房；这些百姓也渴望埃及和那里的肉食。\par

所以称梅瑟为老师、乳父和领路者并不为过\BibleRef{出 16:3}；这人的智慧极大。然而引导已为哲学家的人和统治无理性的孩童是两回事。如果你还愿意听另一个例子：正如乳母对孩子说，当你方便时，要提起你的衣服，并且只要你坐着，也要如此；梅瑟也是这样做的。\BibleRef{申 23:13} 因为所有情欲在孩童中都是专横的（因为他们还没有克制它们的工具）；虚荣、欲望、非理性、愤怒、嫉妒；就像在孩童中一样，这些也占上风；他们唾弃梅瑟，殴打他。又如同当孩童拿起一块石头，我们都大喊：不要扔！他们也是如此拿起石头对抗他们的父亲；梅瑟就逃离了他们。而且，就像一位父亲有什么饰物，孩子喜爱饰物，就向他索要；同样，达堂和阿彼兰一党为了司铎职位而反叛时，也是如此。\EditorialNote{户 16} 此外，他们是所有人中最嫉妒、最小气、各方面都不完美的。\par

那么，请问，基督当时应该显现吗？当时应该给他们这些真智慧的教导吗？那时他们正被情欲激怒，像疯狂追逐母马的马，是金钱和肚腹的奴隶？不，祂只是把智慧的道理浪费在与那些无理解力的人的谈话中；他们既不会学到一样，也不会学到另一样。正如一个在教字母之前就教阅读的人，永远连字母也教不会；当时也正是如此。但现在不是这样了，因为靠着天主的恩宠，许多忍耐、许多德行已被栽种在各地。让我们为一切感恩，不要过分好奇。因为我们不知道那合适的时间，只有祂——时间的创造者、万世的造物主——知道。\par

那么，凡事我们都当顺服祂：因为这就是光荣天主，而不是向祂要求祂所做的事的理由。亚巴郎也是这样光荣天主的；我们读到：\BibleQuote{且深信祂所应许的，祂也必能成就}。\BibleRef{罗 4:21} 他甚至不问将来；但我们却连过去的事也要追究。请看这是何等大的愚昧、何等大的忘恩！但让我们从今以后停止这样做，因为这样得不到任何益处，反而有大害；且让我们的心感恩地归向我们的主，并向天主献上光荣，好使我们因一切事物奉献感谢的祭，而得配祂的慈爱，借着祂独生子的恩宠与对人的爱，与祂一同，等等。\par

\SourceRecord{Translated by John A. Broadus.；Nicene and Post-Nicene Fathers, First Series；Vol. 13.；Edited by Philip Schaff.；Buffalo, NY: Christian Literature Publishing Co.,；1889.；<http://www.newadvent.org/fathers/230304.htm>.}

% structure: p0001 source=h1 target=section reason=page-title

\section{哥罗森书讲道辞 第五篇}

\BibleRef{哥 1:26-28}\par

\Quote{就是那从世世代代以来所隐藏的奥迹，但现在却显示给祂的圣徒们；天主愿意使他们知道，这奥迹在外邦人中是何等丰富的光荣，这奥迹就是基督在你们内，作了光荣的希望：我们所传扬的，就是祂，劝诫一切人，用一切智慧教训一切人，好把一切人在基督内成全地呈献给天主。}\par

说了我们来到了哪里，并显示了天主的慈爱和尊荣，藉着所赐之事的伟大，他又引入了另一个提升它们的考虑，即，在我们之前没有人认识祂。正如他在\BibleBook{厄弗所}{书}也说的：无论天使、宰制者，或任何其他受造的能力，都不认识，唯独天主圣子认识。\BibleRef{弗 3:5} 他说，不是简单地隐藏，而是\InnerQuote{完全隐藏}，并且即使现在才实现，却是古老的，从起初天主就愿意这些事，它们是这样计划的；但为什么，他还没有说。\InnerQuote{从万世以来}，从起初，可以这么说。他称那\OriginalTerm{奥迹}{mystery}为除了天主之外无人知道的，这是有理由的。它在何处隐藏？在基督内；正如他在\BibleBook{厄弗所}{书}\BibleRef{弗 3:9}中所说，或者如先知所说：\Quote{从永远到永远，你是天主。}\BibleRef{咏 89:2} 但现在已显示出来，他说，\Quote{给祂的圣徒}。所以这完全是天主的管理。\Quote{但现在已显示出来}，他说。他并没有说\Quote{已成就}，而是说\Quote{已显示给祂的圣徒}。所以它甚至现在仍是隐藏的，因为只显示给了祂的圣徒。\par

因此，不要让其他人欺骗你们，因为他们不知道。为什么只给他们？他说：\Quote{祂所愿意的}。请看，他处处堵住他们的疑问之口。他说：\Quote{天主愿意使他们知道}。然而，祂的意愿并非没有理由。为了让他们对恩宠负责，而不是让他们自高自大，好像是自己成就的，他说：\Quote{祂所愿意使他们知道的}。\Quote{这奥迹在外邦人中是何等丰富的光荣}。他说话崇高，并累积强调，出于极大的热忱，追求层层递进的放大。因为这也是一个放大，即不确定地说：\Quote{这奥迹在外邦人中是何等丰富的光荣}。因为这在\Person{外邦人}中最为明显，正如他在别处也说：\BibleQuote{使外邦人因天主的仁慈而光荣天主}\BibleRef{罗 15:9}。因为这奥迹的伟大光荣在其他人中也显明，但在这\Person{外邦人}中更甚。因为突然之间，将比石头更无知的人提升到天使的尊位，仅仅通过空言和信德，没有任何劳苦，这确实是奥迹的光荣和丰富：就好像一个人取一条狗，完全被饥饿和疥癣耗尽，污秽不堪，连动都不能动，被扔在外面，然后突然把它变成一个人，并让它出现在王座上。他们曾惯于崇拜石头和地，但如今他们知道他们自己比天和太阳更好，而且整个世界都服侍他们；他们是魔鬼的俘虏和囚徒：突然之间，他们被放在魔鬼的头上，指挥它并鞭打它；从魔鬼的俘虏和奴隶，他们成为天使和总领天使之主的身体；从不知道天主是什么，突然之间成为天主宝座的分享者。你愿意看到他们跨越了无数的台阶吗？首先，他们必须学习石头不是神；第二，它们不仅不是神，甚至低于人；第三，低于野兽；第四，低于植物；第五，他们汇集了两极：不仅石头，连地也不是神，动物、植物、人、天都不是；或者重新开始：不是石头、动物、植物、元素、上界、下界、人、魔鬼、天使、总领天使、任何上面的能力——都不应该被人性所崇拜。他们好像从深渊中被拉上，必须学习万有的主就是天主，只有祂是应当崇拜的；美德的生活是好事；现世的死亡不是死亡，现世的生活不是生活；身体复活，变成不朽，升天，甚至获得不死，与天使并列，被移置到那里。但那位在下界者，一举跨越所有这些台阶，被安置在高高的宝座上，使那比石头还低的，在权力上高于天使、总领天使、宝座、宰制者。真的，\Quote{这奥迹的光荣的丰富是什么？}就像有人突然使一个愚人变成哲学家；不，无论说什么，都不足以形容：因为连\Person{保禄}的话也是无定的。他说：\Quote{这丰富是什么？}他说：\Quote{这奥迹在外邦人中的光荣的丰富，就是基督在你们内？}再次，他们必须学习那位在上，统治天使和宰制者以及所有其他能力的，下降到下界，成为人，经历无数事情，复活，被接升天。\par

这一切都属于奥秘；他连同崇高的赞美一同说出，说：\Quote{基督在你们中间吗？}但如果祂在你们内，为什么还寻求天使？\Quote{这奥秘。}因为还有别的奥秘。但这确实是一个奥秘，是无人知道的，是不可思议的，是超出常理的，是隐藏的。他说：\Quote{基督在你们中，} \Quote{是荣耀的希望，我们所传扬的，} 是从高处带来祂。\Quote{我们所传扬的，} 而不是天使：\Quote{教导}和\Quote{劝诫}；不是以命令的方式，也不是用强迫，因为这也是天主对人的慈爱，不以暴君的方式带人归向祂。既然他说了一件大事，\Quote{教导，} 他又加了\Quote{劝诫}，这更像是父亲而非教师。他说：\Quote{我们所传扬的，劝诫每个人，教导每个人，用一切智慧。}所以需要一切智慧。也就是说，用智慧说一切话。因为学习此类事情的能力并非人人都具备。\Quote{为使我们在基督耶稣内将每个人呈献为成全的。}你说什么？\Quote{每个人}？是的，他说，这是我们热切渴望做的。因为如果这不成呢？蒙福的\Person{保禄}曾努力。\Quote{成全。}所以这是成全，其他是不成全的：这样，如果一个人甚至没有全部的智慧，他就是不成全的。\Quote{在基督耶稣内成全，} 而不是在法律中，也不是在天使中，因为那不是成全。\Quote{在基督内，} 也就是说，在认识基督内。因为知道基督做了什么的人，会有更高的思想，而不会满足于天使。\par

\BibleQuote{在基督耶稣内}；第29节。\BibleQuote{我也为此劳苦，奋力挣扎。}\BibleRef{哥 1:29}他没有仅仅说\Quote{我愿意}，也不是随便地说，而是\Quote{劳苦，奋力挣扎}，带着极大的热忱，许多不眠。如果我为了你们的好处这样警醒，你们更应当如此。然后，又表明这是出于天主，他说：\BibleQuote{按着祂在我里面运行的德能。}\BibleRef{哥 1:29}他表明这是天主的工程。现在，那使我为此刚强的，显然愿意这事。因此，他在开头也说：\BibleQuote{借天主的旨意。}\BibleRef{哥 1:1}所以，他这样表达不仅是出于谦逊，也是坚持圣言的真实。\Quote{并且奋力挣扎。}说这话，表明有许多人在与他争战。然后，他的深情是伟大的。\par

第二章第1节。\BibleQuote{我愿意你们知道，我为你们并为在老底嘉的人怎样奋力挣扎。}\BibleRef{哥 2:1}\par

然后，免得这似乎是由于他们特有的软弱，他将别人也与他们一同包括进来，并且尚未定他们的罪。但他为什么说：\BibleQuote{以及那些在肉身上未曾见过我面的人}？\BibleRef{哥 2:1}他在这里以属神的方式表明，他们在灵里时常看见他。他也为他们的伟大爱情作证。\par

第2至3节。\BibleQuote{为使他们的心受安慰，他们在爱中结合，达到一切富足的理解之全备确信，好认识圣父和基督的奥秘；在祂内藏着一切智慧和知识的宝藏。}\BibleRef{哥 2:2-3}\par

从此刻起，他急切地、痛苦地要进入教义，既不定他们的罪，也不为他们开脱。他说：\Quote{我奋力挣扎。}为了什么目的？使他们结合。他的意思是类似这样：使他们在信德中站稳。但他没有这样表达；而是减轻了指控的成分。也就是说，使他们用爱结合，而不是用必需或强迫。因为如我所说，他总是避免冒犯，把事情留给他们自己；因此\Quote{奋力挣扎}，因为我希望这事出于爱，而且心甘情愿。因为我不希望仅用嘴唇，也不希望只是使他们聚集，而是\BibleQuote{使他们的心受安慰。}\BibleRef{哥 2:2}\par

\BibleQuote{他们在爱中结合，达到一切富足的理解之全备确信。}\BibleRef{哥 2:2}也就是说，使他们关于什么也不怀疑，使他们在一切事上全备确信。但我指的是因信德而来的全备确信，因为有一种全备确信是由论证而来，但那是不值得考虑的。他说，我知道你们相信，但我愿你们全备确信：不仅\Quote{达到富足}，而且\Quote{达到一切富足}；使你们的全备确信是强烈的，以及在一切事上。且看这位蒙福者的智慧。他没有说：\Quote{你们做不好，因为你们不确信，}也没有定他们的罪；而是说：你们不知道我何等渴望你们得以确信，并且不仅是这样，而是带着理解。因为他说到信德；他说，不要以为我指的是肤浅无益的，而是带着理解和爱。\BibleQuote{好认识圣父和基督的奥秘。}\BibleRef{哥 2:2}所以，这就是天主的奥秘：借圣子被带到祂面前。\BibleQuote{以及基督，在祂内藏着一切智慧和知识的宝藏。}\BibleRef{哥 2:3}但如果它们在祂内，那么无疑祂也明智地在这时来临。那么为什么有些愚昧的人反对祂，说：\Quote{看祂怎样与更单纯的人交谈。}\BibleQuote{在祂内藏着一切宝藏。}\BibleRef{哥 2:3}祂自己知道一切。\Quote{隐藏着}，因为你们不要以为已经拥有一切；它们连天使也隐藏着，不仅对你们；所以你们应当从祂那里求问一切。祂自己赐予智慧和知识。现在，说\Quote{宝藏}，表明它们的广大；说\Quote{一切}，表明祂一无所不知；说\Quote{隐藏}，表明唯有祂知道。\par

第4节。\BibleQuote{我说这话，免得有人用花言巧语迷惑你们。}\BibleRef{哥 2:4}\par

你看见了吗？他说，我因此说了这话，使你们不从人那里寻求。他说：\Quote{迷惑你们，} \Quote{用花言巧语。}因为如果有人说话，并且说得动听呢？\par

第5节。\BibleQuote{因为我肉身虽不在你们那里，但心灵却与你们同在。}\BibleRef{哥 2:5}\par

此处本来可以直接说：\Quote{即使我在肉身以外，我仍然知道那些欺骗者}；但他却以赞美作结：\Quote{欢喜中看到你们的秩序，并你们在基督内的信德的坚定。}他说\Quote{你们的秩序}，意指你们良好的秩序。\Quote{并在基督内的信德的坚定。}这更是一种嘉许。他没有说\Quote{信德}，而是说\Quote{坚定}，如同对站立整齐、稳固的士兵所说。凡是坚定的，既非欺骗也非考验所能动摇。他说，你们不但没有跌倒，甚至没有人能使你们混乱。他把自己放在他们之上，好让他们如同他在场一样敬畏他；因为秩序就是这样维持的。由稳固产生紧凑，因为当你将许多东西聚集在一起，并将它们紧密而不可分地黏合起来时，你就会产生稳固；正如墙壁那样。但这是爱德特有的工作；因为它将那些单独的人紧密地黏合并结合起来，使之稳固。信德也做同样的事，它不允许推理侵入。因为推理使人分裂、松散，而信德则带来稳固与紧凑。\par

因为既然天主赐予我们超越人类推理的恩惠，祂便适当地引入了信德。当要求推理时，是无法保持稳固的。请看我们所有崇高的教义，它们多么缺乏推理，而唯独依靠信德。天主不在任何地方，又无处不在。还有什么比这更缺乏理性的呢？每个单独的说法本身就充满困难。事实上，祂不在空间中；也没有任何地方是祂所在的。祂不是受造的，祂没有创造自己，祂从未开始存在。如果没有信德，什么推理能接受这一点？它难道不显得完全荒谬，比谜语还难解吗？\par

说祂无始、非受造、无界限、无限，正如我们所说，有明显的困难；但让我们思考祂的非物质性，我们能否藉推理来探究这一点？天主是非物质的。什么是非物质的？只是一个空洞的词，再无其他；因为理智没有接收到任何东西，也没有给自己留下任何印象；因为如果它留下印象，就会落入本性，并构成形体。所以口虽说出，理智却不知道自己在说什么，只知道一点：它不是形体，这就是它所知道的一切。我为什么提到天主呢？关于灵魂——它是受造的、有界限的、有范围的——什么是非物质性？你说！你展示！你不能。它是空气吗？但空气是形体，即使不紧密，并且有许多证据表明它是一种可让性的形体。火是形体，而灵魂的运作是非物质的。为什么？因为它渗透各处。如果本身不是形体，那么非物质的东西存在于空间中，因此它是有范围的；有范围的东西有形状；形状是线性的，而线属于形体。再者，无形的东西，能有什么概念？它没有形状、没有形式、没有轮廓。你看到理智如何变得头晕目眩吗？\par

再者，那本性（即天主的本性）是不受邪恶影响的。但祂凭自己的意志也是善的；因此它是可受影响的。但不能这样说，绝对不可！再者，祂是愿意还是不愿意被带入存在？但也不能这样说。再者，祂包围世界吗？还是包围不了？如果祂不包围世界，祂自己就是有界限的；但如果祂包围世界，祂的本性就是无限的。再者，祂包围自己吗？如果祂包围自己，那么祂对自己就不是无始的，而只是对我们而言；因此祂的本性不是无始的。处处必须承认矛盾。\par

你看到黑暗有多大；处处都需要信德。这才是稳固的。但如果你愿意，让我们来看那些较小的事物。那个实体有一种运作。祂的运作是什么？是某种运动吗？那么祂就不是不变的：因为被运动的东西不是不变的：因为它从不动变成运动。然而祂却在运动，从不静止。但告诉我是什么运动；在我们这里有七种：下、上、内、外、右、左、圆周，或者如果不是这些，还有增加、减少、生成、毁灭、改变。但祂的运动不是这些，而是像心灵运动那样？不，也不是。绝对不可！因为心灵在许多情况下甚至荒谬地运动。意愿是运作吗？还是不是？如果意愿就是运作，而祂愿意所有的人都好，都得救\BibleRef{弟前 2:4}，为什么没有实现？但意愿是一回事，运作是另一回事。因此意愿不足以构成运作。那么圣经怎么说：\Quote{祂成就了祂所愿意的一切}\BibleRef{咏 113:11}？并且癞病人对基督说：\Quote{祢若愿意，就能使我洁净}\BibleRef{玛 8:2}。因为如果运作随着意愿而来，那该怎么说？你们愿意我再说一件事吗？现存之物如何从无中造出？它们将如何化为乌有？天之上是什么？再之上是什么？再再之上是什么？以至无穷。地之下是什么？海洋，再之下是什么？再再之下？此外，右边和左边，不也有同样的困难吗？\par

但这些确实是看不见的事物。你们愿意我把话题引向那些看得见的、已经发生的事物吗？告诉我，那野兽如何将约纳含在腹中而不使他死亡？它没有理性，它的动作不受控制吗？它如何放过义人？热气如何没有使他窒息？如何没有使他腐烂？因为仅仅在深渊里就已经无法设法了，而在动物的腹中又是那样热，就更加无法解释了。如果我们从内部呼吸空气，那呼吸如何足够供养两个生物？它又如何将他完好无损地吐出来？他又如何说话？他又如何镇定自若、祈祷呢？这些难道不是难以置信的吗？如果我们用推理来检验，它们难以置信；如果凭信德，它们是极其可信的。\par

我还要说什么比这更甚的事吗？麦子在土中腐朽，然后复活。看，奇迹，相反，却彼此超越；不朽坏是奇迹，朽坏后复活更是奇迹。那些嘲笑这些事、不信复活、说\Quote{这根骨头如何与那根接合？}并引入诸如此类愚昧故事的人在哪里？告诉我，厄里亚如何乘着火战车升天？火通常焚烧，而不是托举上升。他如何活了这么久？他在什么地方？这事为何发生？哈诺客被接去了哪里？他像我们一样吃同样的食物吗？什么阻止他在这里？不，难道他不吃？那他为何被接去？看，天主如何一点一点地教导我们。祂接去了哈诺客；这不是什么了不起的大事。这为我们预备了厄里亚的被提升天。祂将诺厄关在方舟内\BibleRef{创 7:7}；这也不是什么了不起的大事。这为我们预备了先知被关在鲸鱼腹中。就连古时的事也需要先行者和预象。因为正如在梯子上，第一级通向第二级，从第一级不可能跳到第四级，而第二级通向第三级，这样一级一级地通往下一级；同样，也不可能在到达第一级之前到达第二级。这里也是如此。\par

请注意预象的预象，你会在雅各伯所见的梯子中辨别出这一点。经上说：\BibleQuote{上主立在上面，天使在天梯上上去下来}\BibleRef{创 28:13}。曾预言圣父有一圣子；这需要被相信。你愿意我从哪里给你显示这事的预象？从上往下？从下往上？因为祂无痛苦地生，因此那不育者先怀孕。让我们更进一步。需要相信祂从自己生出。那么怎样？这事确实以一种隐晦的方式发生，如同在预象和影子里，但它确实发生了，并且随着进程变得更为清晰。一个女子单独从男人形成，而他保持完整无恙。再者，需要有一个关于童贞女怀孕的确切记号。因此不育者怀孕，不是一次，而是第二次、第三次，多次。那么，对于祂由童贞女所生，不育者是预象，它把心灵引向信德。再者，这也是一个预象，表明天主能够独自生。因为如果人是主要行动者，而生育在没有他的情况下以一种更卓越的方式发生，那么更不用说，从最高行动者那里诞生一位。还有另一种生育，是真理的预象。我指的是我们藉着圣神的生育。对此，不育者又是一个预象，即它不是从血统来的\BibleRef{若 1:13}；这属于上面的生育。一个——如同预象——表明生育是无痛苦的；另一个表明它可能从上面而来。\par

基督在上，统治万有；这需要被相信。同样的事在地上关于人也发生。\BibleQuote{让我们按我们的肖像，照我们的模样造人}\BibleRef{创 1:26}，使他管理一切走兽。这样，祂不是用言语，而是用行动教导我们。乐园显示了他本性的独特性，以及人是万物中最好的。基督将要复活；看，现在有多少确切的记号：哈诺客、厄里亚、约纳、火窑、诺厄的情形、圣洗、种子、植物、我们自身的生育、所有动物的生育。因为既然一切都系于此，它比其他任何事都拥有更丰富的预象。\par

宇宙并非没有眷顾，我们可以从自身周围的事物推测，因为如果没有照料，什么都不会持续存在；就连畜群和所有其他事物都需要治理。宇宙不是偶然造成的，地狱是一个证明，诺厄时代的洪水、烈火、埃及人在海中的淹没、旷野中所发生的事，也都是证明。\par

同样，需要许多事物为圣洗预备道路；是的，成千上万的事物：例如，旧约中的那些事、水池中的那些事、不健康者的洁净、洪水本身、所有在水中完成的事、若翰的洗礼。\par

需要相信天主交出祂的圣子；一个人预先做了这事，圣祖亚巴郎。因此，如果我们愿意，我们将在圣经中寻找所有这些事的预象。但让我们不要厌倦，而是藉着这些事调谐自己。让我们坚定持守信德，并表现出严谨的生活：这样，藉着万事感谢天主，我们才能被算为配得那应许给爱祂之人的福乐，藉着我们的主耶稣基督的恩宠和慈爱，愿光荣归于祂，等等。\par

\SourceRecord{Translated by John A. Broadus.；Nicene and Post-Nicene Fathers, First Series；Vol. 13.；Edited by Philip Schaff.；Buffalo, NY: Christian Literature Publishing Co.,；1889.；<http://www.newadvent.org/fathers/230305.htm>.}

% structure: p0001 source=h1 target=section reason=page-title

\section{哥罗森人书讲道辞第六篇}

\BibleRef{哥 2:6,7}\par

\BibleQuote{你们既然接受了基督耶稣为主，就该在他内行动生活，在他内生根修建，坚定于你们所学得的信德，并充满感恩。}\par

他再次用他们自己的见证预先提醒他们说：\Quote{你们既然接受了。} 他说，我们并没有引入什么奇怪的添加，你们也没有。\Quote{你们要在他内行走，} 因为他是通往父的道路；不是在天使内；这条路不通往那里。\Quote{生根}，即固定；不是一时走这条路，一时走那条路，而是\Quote{生根}；生根的，永不移动。请看他所用的词语是何等恰当。\Quote{并被建立}，即藉着思想达到他。\Quote{且坚定}于他，即持守他，建造在根基上。他表明他们已经跌倒，因为\Quote{建立}一词有此含义。因为信德确实是一座建筑；既需要坚固的根基，也需要牢固的构造。因为如果有人不建在稳固的根基上，它会摇动；即使建了，如果不牢固，它也不会站立。\Quote{正如你们所学的。} 再次，\Quote{正如}一词。\Quote{充满}，他说，\Quote{感恩}；因为这是善意之人的本分，我所说的不仅是感恩，而且是极其丰盛，若可能，超出你们所学的，以极大的热忱。\par

第8节。\BibleQuote{你们要谨慎，恐怕有人用哲学和虚妄的诡诈来\OriginalTerm{掳掠}{makes spoil}你们。} \BibleRef{哥 2:8}\par

你看他如何表明那人是个贼，是外人，是悄悄进来的？因为他已经表明他进来了。\Quote{谨慎。} 他巧妙地说\Quote{掳掠}。正如有人从下面挖土，可能没有明显迹象，但它逐渐下沉，所以你们也要谨慎；因为这正是他的要点，甚至不让他自己被察觉。如同有人每天偷窃，而他对（屋主）说：\Quote{谨慎，恐怕有人偷窃}；他指明了途径——通过这个途径——如同我们说，通过这个房间；所以，他说，\Quote{通过哲学}。\par

因为\OriginalTerm{哲学}{philosophy}一词看似高尚，他加上了\Quote{和虚妄的诡诈}。因为也有一种好的\OriginalTerm{诡诈}{good deceit}；许多人曾被这种诡诈欺骗，其实根本不该称之为诡诈。耶肋米亚对此说道：\BibleQuote{上主，你欺骗了我，我被你欺骗了} \BibleRef{耶 20:7}；因为这种诡诈根本不该称之为诡诈；因为雅各伯也曾欺骗他的父亲，但那不是诡诈，而是一种\OriginalTerm{经济}{economy}。他说：\Quote{通过他的哲学和虚妄的诡诈，按照人的传统，按照世界的\OriginalTerm{元素}{elements of the world}，而不是按照基督。} 现在他开始责备他们遵守特定日子，用\Quote{世界的元素}指太阳和月亮；正如他在致迦拉达人书中所说：\BibleQuote{你们怎么又回到那些软弱无用的元素那里去呢？} \BibleRef{迦 4:9} 他没有说\Quote{日子的遵守}，而是概括地说\Quote{现在世界}，以表明它的虚无：因为如果现在世界算不得什么，那么它的元素就更算不得什么。他首先指出他们领受了何等大的恩惠和恩宠，然后才提出控告，以表明其严重性，并说服他的听众。先知们也是如此。他们总是先指出恩惠，然后加重控告；如依撒意亚说：\BibleQuote{我养育了子女，使他们成长，他们却背弃了我} \BibleRef{依 1:2}；又如：\BibleQuote{我的百姓，我对你做了什么，或者我在什么事上使你厌烦，或者我在什么事上使你劳苦？} \BibleRef{米 6:3} 达味也是如此；如他说：\BibleQuote{我在暴风的隐蔽处听见了你} \BibleRef{咏 80:7}；又说：\BibleQuote{张开你的口，我必充满它。} \BibleRef{咏 80:10} 你到处都会发现同样的情况。\par

那本来是最应尽的本分，即使他们说得有理，也不该被他们说服；然而，撇开义务不谈，即使这样，我们也应避免那些事。他说：\Quote{而不是按照基督。} 因为如果此事是半心半意的，以至于你能够同时事奉两者，你也不应该这样做；然而，事实上，他不容许你\Quote{按照基督}。那些事物使你远离他。他先抨击了希腊人的遵守，然后推翻了犹太人的遵守。因为希腊人和犹太人都遵守许多规矩，但前者出于哲学，后者出于法律。因此，他首先攻击那些罪责更重的人。怎么是\OriginalTerm{不按照基督}{not after Christ}？\par

第9、10节。\BibleQuote{因为在他内，圆满地住着整个天主性，\OriginalTerm{身体的}{bodily}；你们在他内也得到了圆满，他是\OriginalTerm{一切率领者和掌权者的头}{head of all principality and power}。} \BibleRef{哥 2:9,10}\par

请注意，他在控告一人时如何刺透另一人，先给出解答，再提出异议。因为这样的解答不引人怀疑，听众更容易接受，因为说话者的目的并非为自己辩论。因为若为自己，他定会力求不落下风，但此处并非如此。\Quote{因为住在他内}，意即，因为天主住在他内。但为使你不以为他被局限在身体里，像在身体内一样，他说：\Quote{所有神性的\OriginalTerm{圆满}{fullness}有形有体地居住在祂内，并且你们在他内也得以圆满。} 另有人说，他意指教会因祂的神性而圆满，正如他在别处所说：\BibleQuote{就是那充满万有者充满万有}\BibleRef{弗 1:23}，并且\Quote{\OriginalTerm{有形有体}{bodily}}一词在此处如身体之于头。那么他为何没有加上\Quote{即是教会}呢？又有人说，他所谓神性的圆满居住在他内是指圣父而言，但这不对。首先，因为\Quote{\OriginalTerm{居住}{to dwell}}严格来说不能用于天主；其次，因为\Quote{\OriginalTerm{圆满}{fullness}}并非指那接受者，因为\BibleQuote{大地是属于上主的，地上的万有}\BibleRef{咏 23:1}；再者，宗徒说：\BibleQuote{直到外邦人全数进入}\BibleRef{罗 11:25}。\Quote{\OriginalTerm{全部}{the whole}}意指\Quote{全部}。那么\Quote{\OriginalTerm{有形有体}{bodily}}一词想要表达什么呢？\Quote{如同在头里。} 但为何他又重复同样的话？\Quote{并且你们在他内也得以圆满。} 这有什么含义？即你们不比他少。正如它住在他内，也住在你们内。因为\Person{保禄}总竭力使我们接近基督；正如他所说：\BibleQuote{祂已使我们一同复活，一同坐在天上}\BibleRef{弗 2:6}；并且\BibleQuote{我们若能忍，也必与他一同为王}\BibleRef{弟后 2:12}；以及\BibleQuote{怎不也把万有连同他白白赐给我们呢}\BibleRef{罗 8:32}；并称我们为\Quote{同继承人}。至于他的尊严：他\BibleQuote{是一切率领者、掌权者的元首}\BibleRef{弗 3:6}。他是万有之上者，那原因岂非\OriginalTerm{同质}{Consubstantial}吗？然后他奇妙地加上了恩惠，比在\WorkTitle{罗马书}中更为奇妙。因为在那里他确实说：\BibleQuote{内心的割损，在于神，不在于文字}\BibleRef{罗 2:29}，但此处是在基督内。\par

第 11 节。\BibleQuote{在祂内，你们也受了割损，但不是由人手所行的割损，而是借着脱去肉身的基督的割损。}\BibleRef{哥 2:11}\par

请看他是何等接近此事。他说：\Quote{在脱去}，是完全脱去，不仅是脱下。\Quote{\OriginalTerm{罪恶之身}{the body of sins}}。他意指\Quote{\OriginalTerm{旧生活}{the old life}}。他不断以不同方式提到这一点，如他上面所说：\BibleQuote{祂拯救了我们脱离黑暗的权势，并使从前疏远的人与祂和好，好使我们成为圣洁，没有瑕疵。}\BibleRef{哥 1:13} 他说，不再是用\OriginalTerm{刀行的割损}{circumcision with the knife}，而是在基督自己内；因为这种割损不是人手所行的，如同那里，而是圣神。它割损的不是一部分，而是\OriginalTerm{全人}{the whole man}。两者都是身体，但在一种情况下是\OriginalTerm{属肉体地}{carnally}受割损，在另一种情况下是\OriginalTerm{属神地}{spiritually}受割损；但不像犹太人，因为你们脱去的不是肉身，而是罪恶。何时何地？在\OriginalTerm{圣洗}{Baptism}中。他称割损，又称之为\OriginalTerm{安葬}{burial}。请注意他如何再次转到正义行为的话题；他说：\Quote{属于罪恶}，\Quote{属于肉身的}，即他们在肉身上所行的事。他谈到一件比割损更伟大的事，因为他们不仅抛弃了所割损的，更毁灭它们，消灭它们。\par

第 12 节。\BibleQuote{你们与他一同受埋于圣洗中，也与他一同复活，借着信德那使祂从死者中复活的天主的作为。}\BibleRef{哥 2:12}\par

但这不仅是埋葬；因为请看他说什么：\Quote{你们也与他一同复活，借着信德那使祂从死者中复活的天主的作为。}他说的好：\Quote{\OriginalTerm{信德}{faith}}，因为全是出于信德。你相信天主能使人复活，于是你复活了。然后也注意祂的\OriginalTerm{可信性}{worthiness of belief}：他说：\Quote{那使祂从死者中复活的。}\par

他现在显示复活。\Quote{你们从前因过犯和肉身未受割损而死了，我说，祂使你们与他一同复活。}因为你们已处于死亡的判决之下。但即使你们死了，这也是\OriginalTerm{有益的死亡}{profitable death}。请注意他又如何借他加添的话显示他们应得什么：\par

第 13、14、15 节。\BibleQuote{祂赦免了我们一切过犯；涂抹了那在规条上所写的、反对我们的债券，就是与我们为敌的；又将它除去，钉在十字架上；脱去率领者和掌权者的权势，公开羞辱它们，在十字架上凯旋了它们。}\BibleRef{哥 2:13}\par

他说：\Quote{\OriginalTerm{祂赦免了我们}{Having forgiven us}}，\Quote{一切过犯}，就是那些导致死亡的事物。那么怎样？祂任凭它们存在着吗？不，祂甚至涂抹了它们；并非仅仅刮去，以至于看不见。他说：\Quote{在\OriginalTerm{规条}{ordinances}上}〔条例〕。什么规条？\OriginalTerm{信德}{Faith}。相信就够了。祂并未以行为对抗行为，而是以行为对抗信德。下一步呢？涂抹比赦免更进一步；他又说：\Quote{将它除去。}仍未保存它，反而撕碎了它，\Quote{钉在祂的十字架上。}\Quote{脱去率领者和掌权者的权势，公开羞辱它们，在十字架上凯旋了它们。}他从未用过如此高扬的语调。\par

你看见祂为废除债据是多么热切吗？的确，我们所有人都处在罪与惩罚之下。祂亲自通过承受惩罚，除去了罪和惩罚，并且在十字架上受罚。于是祂将债据钉在十字架上；祂如同有权能者，将它撕碎。什么债据？祂或许是指他们向梅瑟说的话，即：\BibleQuote{凡上主所说的话，我们都要遵行和服从}\BibleRef{出 24:3}；若非如此，就是指我们亏欠天主的服从；或者又不是这个，祂是指魔鬼所持有的债据，即天主为亚当所立的债据，说：\BibleQuote{你吃那果子的那一天，必定要死。}\BibleRef{创 2:17} 这债据原本掌握在魔鬼手中。基督并没有把它交给我们，而是亲自将它撕成两半，这是欢欣宽恕者的作为。\par

\BibleQuote{解除了率领者和掌权者的武装。}祂是指魔鬼的权势；因为人性曾披戴上这些，或者因为他们仿佛有一种掌控，当祂成为人时，祂就解除了这种掌控。\BibleQuote{祂公然羞辱他们}是什么意思？他说得真好；魔鬼从未陷入如此可耻的境地。因为当它指望拥有祂时，反而失去了它所拥有的；当那身体被钉在十字架上时，死人复活了。在那里，死亡受了伤，从一个死去的身体上受到了致命一击。如同一位运动员，当他以为击中了对手时，自己却被致命地抓住；同样，基督也显示，带着信心死是魔鬼的耻辱。\par

因为如果魔鬼有能力，它会想尽一切办法让人相信祂没有死。事实上，看到祂的复活，以后的所有时代都提供了明确的证明；而祂的死，除了发生的那一刻，没有其他时间能提供证明；因此，祂公开地在众人眼前死去，却没有公开地复活，祂知道后来的时期会为真理作证。因为在世界注视之下，蛇被高举在十字架上杀死，这正是奇妙之处。魔鬼为了不让祂秘密死去，没有做什么呢？听比拉多说：\BibleQuote{你们自己把他带去，钉在十字架上吧！我在他身上查不出什么罪来}\BibleRef{若 19:6}，并且用千百种方式抗拒他们。犹太人又对祂说：\BibleQuote{如果你是天主子，从十字架上下来吧！}\BibleRef{玛 27:40} 然后，当祂受了致命伤，祂没有下来，为此祂也被埋葬；因为祂本可以立即复活，但祂没有这样做，好让事实被相信。在私人死亡的情况下，有可能将其归咎于昏厥，但在这里，连这一点也不可能。因为士兵们没有像对待其他两人那样打断祂的腿，好表明祂确实死了。埋葬身体的人也是众所周知的；因此犹太人自己也与士兵一起封了石头。因为最被关注的就是这一点：不能隐秘进行。见证来自敌人，来自犹太人。听他们对比拉多说：\BibleQuote{那个骗子活着的时候曾说：三天以后我要复活。所以请下令把坟墓}\BibleRef{玛 26:63} 交由士兵把守。这事就这样做了，他们也亲自封了石墓。后来又听他们对宗徒们说：\BibleQuote{你们有意把此人的血归到我们身上。}\BibleRef{宗 5:28} 祂不容许祂十字架的样子受到羞辱。既然天使从未受过类似的事，祂便为这件事做了一切，显示祂的死成就了一番大业。这就像一场决斗。死亡击伤了基督；但基督受伤后，却杀死了死亡。那看似不死者，被一个必死的身体消灭了；整个世界都看见了。真正奇妙的是，祂没有把这事交给别人。但又产生了一种与先前不同的第二种债据。\par

那么，我们当心，免得说了\Quote{我弃绝撒殚，并穿上祢，基督}之后，反而因此被定罪。不过，这不应该称为\Quote{契据}，而应称为\Quote{盟约}。因为那是\Quote{契据}，凭此人对债务负责；但这是盟约。它没有惩罚，也不说\Quote{若做了这事或若不做这事}；梅瑟洒盟约之血时所说的，天主也由此应许了永生。这一切都是盟约。那里是奴隶与主人，这里则是朋友与朋友；那里说：\BibleQuote{你吃的日子必定死}\BibleRef{创 2:17}，立即威胁；但这里却没有这样的事。天主来了，这里是赤身，那里也有赤身；然而，那里犯罪的人被赤身露体，是因为他犯了罪；但这里，人被赤身露体，是为了使他获得自由。那时，人脱去了他所拥有的荣耀；现在，他脱去旧人；在上场（比武）之前，像脱衣服一样轻易地脱去。他被傅油，如同即将进入竞技场的摔跤手。因为他立刻重生；正如第一个人那样，不是逐渐地，而是立刻。（他受傅油）不像古时的司祭，只在头上傅油，而是更为丰盛。因为古时的司祭在头上傅油，在右耳和手上傅油\BibleRef{肋 8:23}，以激励他服从和行善；但这人却是全身傅油。因为他来不仅是为受教，而是为了摔跤和锻炼；他被提升到另一个受造界。当人宣认（信）永生时，他就宣认了第二个受造界。\BibleQuote{上主天主用地上的灰土形成了人}\BibleRef{创 2:7}；但现在不再是灰土，而是圣神；他用圣神形成人，用圣神协调人，正如他自己在童贞玛利亚的胎中一样。祂没有说\Quote{在乐园里}，而是说\Quote{在天上}。不要以为因主题是大地，就在地上进行；他被迁到那里，到天上，这些事在天上办理，在天使中间：天主将你的灵魂提升到高处，在上面重新协调它，将你安置在君王宝座附近。他在水中形成，接受神代替灵魂。在他形成之后，天主带到他面前的不是野兽，而是魔鬼和他们的君王，并说：\BibleQuote{践踏在蛇和蝎子上}\BibleRef{路 10:19}。祂没有说：\BibleQuote{让我们按我们的肖像，照我们的模样造人}\BibleRef{创 1:26}，而是说什么？\BibleQuote{他赐给他们权能，成为天主的子女；由天主而生的}\BibleRef{若 1:12}。 然后，为使你不听从蛇，他立刻教你这样宣认：\Quote{我弃绝你}，即\Quote{无论你说什么，我都不听从你。}然后，为使你不通过他人而被他毁灭，又说了：\Quote{并弃绝你的虚荣、你的事奉和你的天使。} 祂不再派他看守乐园，而是使他在天上有公民身份。因为他一上来就念这些话：\Quote{我们的天父，你在天上……愿祢的旨意奉行在人间，如同在天上}。 你眼前不再有平原，你不见树木，也不见泉源，而是立刻领受主自己，与祂的身体相融合，与那高处的身体相混合，魔鬼无法靠近那里。那里没有女人，让他去接近并欺骗较软弱者；因为经上说：\BibleQuote{不再分女性与男性}\BibleRef{迦 3:28}。如果你不下降到魔鬼那里，他就没有能力上升到你在的地方；因为你在天上，而天上魔鬼无法接近。那里没有知善恶树，只有生命树。你的肋骨不再被塑造为女人，但我们都从基督的肋旁成为一体。因为那些被人傅油的人不受蛇的伤害，你也完全不受伤害，只要你受了傅油；这样你就能抓住蛇并扼死它，\BibleQuote{践踏在蛇和蝎子上}\BibleRef{路 10:19}。然而，恩赐越大，惩罚也越大。从乐园堕落的人不可能住在\BibleQuote{乐园前面}\BibleRef{创 3:24}，也不能重新上升到我们堕落的地方。但此后呢？地狱和不死的虫。但愿我们中没有人应受此惩罚！但我们要活出德性，努力一生遵行祂的旨意。让我们中悦天主，使我们能逃脱惩罚，并获得永恒的福乐。愿我们众人，因祂的恩宠和慈爱，堪当获得这一切。\par

\SourceRecord{Translated by John A. Broadus.；Nicene and Post-Nicene Fathers, First Series；Vol. 13.；Edited by Philip Schaff.；Buffalo, NY: Christian Literature Publishing Co.,；1889.；<http://www.newadvent.org/fathers/230306.htm>.}

% structure: p0001 source=h1 target=section reason=page-title

\section{哥罗森书讲道 第七篇}

哥罗森书 \BibleRef{哥 2:16-19}\par

\BibleQuote{所以，不要让任何人在饮食上，或在庆节、新月或安息日的问题上判断你们：这些乃是将来之事的影子，但实体是属于基督的。不要让人因自愿的谦卑和敬拜天使，而夺去你们的奖赏，他沉溺于自己未见之事，凭着血肉的意念无端自大，不紧持元首；全身都从他藉着关节和筋络获得供应和联络，遂以天主的增长而增长。}\par

他先是隐晦地说：\BibleQuote{你们要谨慎，恐怕有人按人的传统掳掠你们}\BibleRef{哥 2:8}；又更早前说：\BibleQuote{我说这话，免得有人以花言巧语迷惑你们}\BibleRef{哥 2:4}；这样预先占据他们的灵魂，在其中培育焦虑的思想；然后，在插入了那些恩惠并增强了这一效果之后，他最后才提出责备，说：\BibleQuote{所以，不要让任何人在饮食上，或在庆节、新月或安息日的问题上判断你们。}你看出他如何贬低这些事吗？他说，如果你们已经获得了这些，为什么还要为这些琐碎的事负责？他又轻看它们，说：\BibleQuote{或是在庆节的一部分}，因为事实上他们并没有完全遵守以前的规则，\BibleQuote{或是新月，或是安息日。}他没有说：\Quote{所以不要遵守它们，}而是说：\Quote{不要让人判断你们。}他表明他们在违犯和破坏，却将指责指向别人。他说：不要忍受那些判断你们的人，不，甚至不止如此，但他与那些人辩论，几乎堵住他们的口，说：你们不应该判断。但他不愿责难这些人。他没有说：\Quote{在洁净和不洁净上，}也没有说：\Quote{在住棚节、无酵节和五旬节上，}而是说：\BibleQuote{在庆节的一部分}：因为他们不敢完全遵守；即使他们遵守了，也不是以庆节的方式。他说：\BibleQuote{一部分，}表明大部分已废除了。因为即使他们守安息日，也没有精确地守。\BibleQuote{这些乃是将来之事的影子}；他指的是新约；\BibleQuote{但实体}是\BibleQuote{基督的。}有些人在这里这样标点：\BibleQuote{但实体}是\BibleQuote{属于基督的，}即真理已随着基督而来；另一些人则这样：\BibleQuote{基督的身体，不要让任何人从你们那里夺去，}即阻止你们获得它。\newline{}\OriginalTerm{抢夺奖赏}{\GreekText{καταβραβευθῆναι}}一词用于胜利属于一方，而奖赏属于另一方的情形，即虽然你胜利了，却被阻止。你已凌驾于魔鬼和罪恶之上，为何又使自己屈服于罪恶？因此他说：\BibleQuote{他必须遵守全部法律}\BibleRef{迦 5:3}；又说：\BibleQuote{基督难道成了罪恶的仆役}\BibleRef{迦 2:17}？这是他在写给迦拉达人的书信时说的。当他以\Quote{夺去你们那里}的话使他们充满愤怒后，便开始说：\BibleQuote{自愿的，}他说，\BibleQuote{在谦卑和敬拜天使中，介入自己未见之事，凭着血肉的意念无端自大。}\Quote{在谦卑中}如何，\Quote{自大}又怎样？他表明这一切都出自虚荣。但所说的总意是什么？有些人主张我们必须通过天使而非基督被领近，那对我们来说是太高了。因此他一再提及基督所成就的，\BibleQuote{藉着祂十字架的血}\BibleRef{哥 1:20}；为此他说：\BibleQuote{祂为我们受苦}；\BibleQuote{祂爱了我们}\BibleRef{伯前 2:21}。此外，同样在这件事上，他们又再次被高举。他没有说\Quote{通过引入，}而是说\Quote{敬拜}天使。\BibleQuote{介入自己未见之事}\BibleRef{弗 2:4}。因为他没有见过天使，却好像见过一样。因此他说：\BibleQuote{凭着血肉的意念无端自大，}并非基于任何真实事实。关于这个教义，他自大，并装出一副谦卑的样子。凭着他的血肉之念，而非属灵之念；他的推理是属人的。\BibleQuote{不紧持元首，}他说，\BibleQuote{从他那里，所有的身体。}所有的身体从他那里得到存在和福祉。为什么放弃元首而去紧贴肢体？如果你从他那里堕落，你就迷失了。\BibleQuote{从他那里，所有的身体。}每一个人，无论他是谁，从他那里不仅得到生命，而且还有联系。整个教会，只要她紧持元首，就增长；因为这里不再有骄傲和虚荣的情欲，也没有人幻想的发明。\par

请注意，\Quote{从他}指的是圣子。他说：\BibleQuote{藉着关节和筋络，获得供应和联络，遂以天主的增长而增长}；他意指那合乎天主的增长，即最佳生活的增长。\par

第20节：\BibleQuote{如果你们与基督同死}\BibleRef{哥 2:20}。\par

他将那话放在中间，两边是更激烈的表达。他说：\BibleQuote{如果你们与基督同死于世界的元素，}\BibleQuote{为什么还像活在世界上一样服从条例？}这并非结果，因为本该说的是：\Quote{如何像活着一样服从那些元素？}但撇开这点，他说了什么？\par

第21至22节：\BibleQuote{不可摸，不可尝，不可触；这一切都随使用而败坏，这是照着人的戒律和教导。}\par

他说：你们不在世界上，怎么还顺从它的元素呢？怎么还顺从它的规条呢？且看他怎样嘲笑他们，\Quote{不可摸，不可拿，不可尝}，好像他们是胆小鬼，远远避开某些大事，\Quote{这一切都要随着使用而朽坏}。他挫败了许多人的傲慢，并加上\Quote{照人的规条和教训}。你说什么？你甚至提到法律吗？从今以后，它不过是人的教训，因为时候已经来到。或者，是因为他们败坏了法律，或者，他指的是外邦人的制度。他说，这教训完全是人的。\par

第23节：\Quote{这些事在私意崇拜、谦卑和苦待身体上，确有智慧的表现，但在克制肉欲放纵上，毫无价值。}\BibleRef{哥 2:23}\par

\Quote{外表}，他说，不是能力，不是真理。所以，即使他们有智慧的外表，我们也要远离他们。因为一个人可能显得虔诚、谦逊，并轻视身体。\par

\Quote{在克制肉欲放纵上毫无价值}。因为天主给了它尊荣，但他们不用尊荣对待它。所以，当它是一种教训时，他知道如何称之为尊荣。他说，他们侮辱肉身，剥夺它、剥夺它的自由，不允许它按自己的意愿治理。天主已经尊荣了肉身。\par

第3章第1节\BibleRef{哥 3:1}：\Quote{如果你们已与基督一同复活。}\par

他把他们结合起来，上面已经确立了基督死了。因此他说，\Quote{如果你们已与基督一同复活，就该追求上面的事。}那里没有规条。\Quote{那里基督坐在天主的右边。}何等奇妙！他使我们的心灵升得多高！他使他们充满何等高远的渴望！说\Quote{上面的事}不够，说\Quote{基督所在之处}也不够，而是什么？\Quote{坐在天主的右边。}从那里他开始准备他们以后俯视大地。\par

第2至4节。\Quote{你们要思念上面的事，不要思念地上的事。因为你们已经死了，你们的生命与基督一同藏在天主内。当基督——你们的生命显现时，你们也要与祂一同在光荣中显现。}\BibleRef{哥 3:2}\par

他说，这不是你们的生命，而是另一个生命。他现在急切地想使他们离开，并坚持表明他们已坐在天上，且已经死了；从这两个考虑确立立场：他们不应该寻求这里的事物。因为无论你们是死了，就不该寻求它们；或者你们是在天上，就不该寻求它们。基督显现了吗？你们的生活也没有。它在天主内，在天上。那么，什么时候我们才能生活？当基督——你们的生命显现时，那时你们寻求光荣，寻求生命，寻求喜乐。\par

这是为引导他们离开安逸快乐铺路。这是他惯常的做法：当确立一个论点时，他会突然转向另一个；例如，当谈论那些晚餐中争先的人时，他一下子转到奥迹的遵守上。因为当斥责出其不意地施与时，它有很大的力量。他说，\Quote{它是隐藏的}，从你们那里。\Quote{那时你们也要与祂一同显现。}所以，现在，你们不显现。看，他把他们移入了天上。因为，如我所说，他总是致力于表明他们拥有基督所拥有的一切；在他的所有书信中，主旨就是表明他们在一切事上都是祂的分享者。因此他使用头、身体等词汇，并尽力向他们传达这一点。\par

所以，如果我们那时要被显现，当我们不享受荣耀时，不要悲伤：如果这生命不是生命，而是隐藏的，我们就该像死了一样过这生命。他说，\Quote{那时你们也要与祂一同在光荣中显现。}他说\Quote{在光荣中}，不仅仅说\Quote{显现}。因为珍珠在蚌壳内时也是隐藏的。所以，如果我们受到侮辱，不要悲伤；无论我们遭受什么；因为这生命不是我们的生命，我们是旅客和寄居者。他说，\Quote{因为你们已经死了。}谁那么愚蠢，竟为了一具死尸，死了又埋葬的，买仆人，建房子，或准备昂贵的衣服？没有人。你们也不该如此；但正如我们只寻求一件事，即我们不要赤身露体，同样，在这里我们也只寻求一件事，不再多求。我们的旧人已被埋葬：不是埋在地里，而是埋在水里；不是被死亡消灭，而是被死亡的毁灭者埋葬，不是按自然规律，而是按那比自然更强有力的统治命令。因为自然所做的事，或许可以撤销；但祂命令所做的事，绝不能撤销。没有比这埋葬更有福的了，一切都为此喜乐，包括天使、人类和天使的主。在这埋葬中，不需要殓衣、棺木或任何类似的东西。你愿看见这预象吗？我将给你看一个池子，在其中，一方被埋葬，另一方被提升；在红海中，埃及人沉没其中，而以色列人从其中上来；在同一行动中，祂埋葬了一方，生出了另一方。\par

不要惊讶在圣洗中发生生成与毁灭；因为，请告诉我，分解与黏合，难道不是相反的吗？这是众人皆知的。火的作用也是如此：火溶解并毁灭蜡，却将金属泥土黏合，将其炼成金子。同样，在这里，火的力量抹去了蜡像，却显现了一座金像取而代之；因为在圣洗之前，我们是泥土的，之后却是金子的。这从何处显明呢？请听他这样说：\BibleQuote{第一个人出于地，属于土；第二个人是主，来自天上。} \BibleRef{格前 15:47} 我提到了泥土与金子之间那般巨大的差别；但我发现属天的与属土的差别更大；泥土与金子的差别，远不及属土与属天之物差别之大。我们曾是蜡制的，是泥土所造的；因为情欲的火焰熔化我们，远超过火熔化蜡；任何偶然的诱惑击碎我们，远胜过石头击碎土器。如果你愿意，让我们勾勒出昔日生活的轮廓，看看是否一切不是土和水，充满了波动、尘埃、不稳定和流逝。\par

如果你愿意，让我们审视的不仅是过去的事，还有现在的事，看看我们是否会发现自己所拥有的一切不过是尘土和水。那么，你会告诉我什么呢？权力和势力？因为在现世生活中，没有什么比这些更令人羡慕的了。但人们更容易发现尘土停留在空中，而非这些事物；尤其是现在。因为谁不受它们支配呢？那些贪爱它们的人，那些阉人，那些为钱什么都肯做的人，那些大众的激情，那些更有权势者的愤怒。昨天还高踞审判席、有传令官以刺耳的声音呼喊、有许多人为他开道、傲慢地穿过广场的人，今天却变得卑微低下，失去了一切，如被风吹起的尘土，如流过的溪水。正如尘土被我们的脚扬起，真正的官位也是由那些从事金钱事务、在整个人生中处于脚的地位和状态的人所产生；正如尘土被扬起时占据了大片空气，虽然它本身只是一个小物体，权力也是如此；正如尘土迷瞎眼睛，权力的骄傲也蒙蔽了理智的眼睛。\par

然而，我们要审视那众所祈求的对象——财富吗？来吧，让我们从多方面审视它。它有奢侈，有尊荣，有权力。那么首先，如果你愿意，让我们审视奢侈。它不是尘土吗？是的，它比尘土更快消逝，因为奢侈享受的愉悦只到达舌头，当肚腹填满时，连舌头也感受不到。但有人说，尊荣本身是令人愉快的。然而，当那种尊荣是出于对金钱的考量时，还有什么比这更不愉快呢？当它不是出于自由选择和心甘情愿时，被尊荣的不是你，而是你的财富。因此，正是这一点使富人最受羞辱。因为，请告诉我：假设所有人因你有一位朋友而尊荣你，同时承认你确实一无是处，但迫于他的缘故才尊荣你；他们还能以其他方式如此羞辱你吗？所以，我们的财富成了我们受羞辱的原因，因为它比其拥有者更受尊荣，并且是软弱的证明，而非能力的证明。那么，我们不被视为与尘土和灰烬（金子也是如此）同等价值，反而因它而受尊荣，这是多么荒谬！有道理。然而，轻视财富的人并非如此；因为不受尊荣总比这样受尊荣更好。因为，请告诉我：如果有人对你说，我认为你根本不配受任何尊荣，但为了你仆人的缘故我尊荣你；还有什么比这更糟糕的羞辱吗？如果因着仆人（他们与我们同有同样的灵魂和本性）而受尊荣是耻辱，那么因着更低贱的事物，如墙壁、庭院、金器、衣服而受尊荣，更是耻辱。这确实是轻蔑和羞耻；宁死也不愿如此受尊荣。因为，请告诉我：如果你在骄傲中遭遇危险，而一个低贱可鄙的人愿意将你从危险中解救出来，还有比这更糟的吗？我想对你们说你们彼此谈论城市的话。有一次，我们的城市得罪了皇帝，他下令将整个城市彻底摧毁，包括人、儿童、房屋等一切。（因为君王的忿怒就是如此，他们随意放纵权力，权力是何等大的恶！）城市便处于极端的危险中。然而，邻近的这座沿海城市却去为我们向国王求情；于是，本城居民说，这比城市被夷为平地更糟。因此，这样受尊荣比受羞辱更糟。因为请看尊荣的根源：厨师的双手使我们受尊荣，因此我们应当感激他们；养猪的人为我们提供丰盛的餐桌；还有织工、纺工、金属匠、糖果匠、摆桌者。\par

那么，难道不受尊荣，反倒比因这些人而得尊荣更好吗？除此之外，我还要努力清楚地证明：财富是一种充满羞辱的状态；它贬低灵魂，还有什么比这更羞辱的呢？因为请告诉我，假如一个人容貌俊美，绝伦超群，财富却去对他说，它能使他变丑，使他从健康变为患病，从清凉变为燥热，使他的每一肢体都充满水肿，面部浮肿，全身膨胀；它使他的双脚肿胀，比木头还沉重；使他的肚腹鼓胀，比任何大桶还大；此后，它甚至不允许那些想要医治他的人去治疗他（权力就是如此），却给他如此多的自由，以至于他可以惩罚任何接近他、要使他脱离有害事物的人；那么，请告诉我，当财富在灵魂中产生这些效果时，它怎么可能是光荣的呢？\par

但这种权力比疾病本身更严重；正如一个人患病而不遵从医生的嘱咐，比患病本身更糟糕；财富也是如此，它在灵魂的每一部分制造炎症，并禁止医生接近它。所以，我们不要因他们的权力而祝贺他们，而要怜悯他们；因为如果我看到一个水肿病人躺着，没有人禁止他随意饮用任何饮料和有害的食物，我不会因他的权力而祝贺他。因为权力并不总是在所有情况下都是好事，尊荣也是如此，因为它们也充满傲慢。但如果你不愿意身体随着财富患上这种疾病，你怎么会忽略灵魂，让它不仅染上这个灾祸，还加上另一个呢？因为它全身燃烧着炽热的高烧和炎症，那高烧无人能熄灭，因为财富不允许这样，它说服灵魂那些真正是损失的事是收益，比如不能忍受任何人，并且随心所欲地做一切事。因为人们找不到其他任何灵魂像那些渴望致富的人的灵魂那样充满如此巨大而放纵的情欲。他们为自己描绘了何等无聊的幻想！人们可以看到他们比画半人马、吐火兽、龙足怪兽、斯基拉和海怪的人更加荒诞不经。如果一个人愿意描绘他们的一种情欲，那么无论是斯基拉、吐火兽还是半人马，在这样一个怪物旁边都算不了什么；你会发现它同时包含了一切野兽。\par

或许有人会认为我自己拥有大量财富，因为我对财富带来的后果如此熟悉。传说中有一个国王（因为我要先以希腊传说来证实我的话），据说他如此骄奢淫逸，以至于造了一棵金梧桐树，并在其上方造了一个天空，坐在那里，而且是在入侵一个精于作战的民族的时候。这种情欲难道不是半人马式的，不是斯基拉式的吗？另一个人，又把人们投入木牛中。这难道不是真正的斯基拉吗？甚至那个我刚刚提到的国王，那个战士，财富使他从男人变成了女人，从女人变成了什么？一头野兽，甚至比野兽更低下，因为野兽如果在树下栖息，满足于自然，不再寻求更多；但那个人甚至超越了野兽的本性。\par

那么，还有什么比富人更愚蠢的呢？这源于他们欲望的贪婪。但是，不是有许多人钦佩他吗？因此他们确实分担了他招致的嘲笑。那展示的不是他的财富，而是他的愚蠢。那棵金梧桐树比大地出产的树木好多少？因为自然之物比非自然之物更令人愉悦。但你的金色天空有何意义，愚昧的人？你难道没有看到巨大的财富使人疯狂？它如何点燃了他们？我想他大概连海都不认识，或许很快就会想在海上行走。这难道不是吐火兽吗？不是半人马吗？但在现今时代，也有一些人并不比他逊色，实际上更加愚蠢。因为在愚蠢方面，那些制作银壶、水罐和香水瓶的人，与那棵金梧桐树有何区别，请告诉我？那些女人，我虽羞于启齿，却不得不说，她们制作银制便器，又有什么不同？你们这些制作这些东西的人应该感到羞耻。当基督在挨饿时，你们却如此奢侈放纵？更确切地说，如此愚弄自己！这些人将受到怎样的惩罚？你们还在问为什么有强盗？为什么有杀人犯？为什么有这些邪恶？当魔鬼使你们变得如此可笑的时候。因为仅仅拥有银盘就已经不符合一个追求智慧的灵魂，完全是奢侈；但连不洁的器皿也用银来制作，这难道是奢侈吗？不，我不称之为奢侈，而是愚蠢；不，也不是这个，而是疯狂；甚至比疯狂更糟。\par

我知道许多人因此嘲笑我，但我并不在意，只要这能带来一些益处。事实上，富有确实使人丧失理智、疯狂。如果他们的权力达到如此极端的地步，他们连大地也要用金子做，墙壁也要用金子做，甚至天空和空气也要用金子做。这是何等疯狂，何等不义，何等炽热的狂热！另一个人，按天主的肖像受造，却正在寒冷中死去，而你却为自己置办这些东西？哦，愚蠢的骄傲！一个疯子还会做什么更多的事？你竟对你的排泄物如此尊敬，以至于要用银器来承接？我知道你听到这些话会震惊，但那些制作这些东西的妇人们才应该震惊，以及那些助长这种病态的丈夫们。因为这是放荡、野蛮、不人道、兽性、淫荡。什么斯库拉、什么喀迈拉、什么龙，甚至什么恶魔、什么魔鬼会这样做？基督的益处有何用处？信仰有何用处？当不得不容忍人们像外邦人一样，甚至不是外邦人，而是魔鬼时？如果装饰头部用金子和珍珠是不对的，那么一个人为了如此不洁净的服务而使用银器，又如何能获得宽恕？难道其余的还不够吗？虽然它们也是不可忍受的，椅子、脚凳全都是银的？虽然这些也是出自愚昧。但到处都是过度的骄傲，到处都是虚荣。何处是实用，处处是奢侈。\par

我担心，在这种疯狂的冲动下，妇女的种族会开始呈现出某种怪异的形态：因为她们很可能希望甚至她们的头发也是金子的。否则就申明你完全没有被所说的话影响，也没有被大大激动，并起了渴望，若不是羞耻制止了你，你也不会拒绝。因为如果她们敢于做出比这更荒谬的事，那么我确信，她们更会渴望她们的头发、嘴唇、眉毛和每个部分都用熔金覆盖。\par

但如果你不信，以为我说笑，我就讲述我所听说的，或者更确切地说，现在还存在的事。波斯王的胡须是金子的；那些精通此类工艺的人将金箔缠绕在他的胡须上，如同绕在织线上一般，这被视为奇事被珍藏起来。\par

光荣归于祢，基督！祢以多少美物充满我们！祢如何为我们的健康预备！祢将我们从多大的怪异、多大的无理中释放出来！注意！我预先警告你们，不再劝告，而是命令和嘱咐：愿意的服从，不愿意的悖逆；如果你们妇女继续这样行事，我将不再容忍，不再接纳你们，也不允许你们跨过这道门槛。我需要一群病态的人做什么？再者，在我的培育中，如果我不禁止那些过分的事，那又算什么？然而，保禄禁止了金子和珍珠。\BibleRef{弟前 2:9} 我们被希腊人嘲笑，我们的宗教显得像是寓言。\par

我给男人这样的建议：你来这里是要受教于属灵的哲学？除去那种骄傲！这是我给男女双方的共同建议；如果有人不这样做，从今以后我不会容忍。门徒们只有十二个，听基督对他们说什么：\Quote{你们也愿意走吗？}\BibleRef{若 6:67} 因为如果我们永远奉承你，何时才能纠正你？何时才能服侍你？有人说：\Quote{但还有其他教派，人们会转过去。} 这是一个冷淡的论点：\Quote{一个遵行主的旨意的人，胜过一万个背命者。}\BibleRef{德 16:3} 因为，告诉我，你宁愿选择什么？是拥有一万个逃跑的、做贼的仆人，还是只有一个爱你的仆人？看！我告诫并命令你们，要打碎那些脸上的华丽装饰和我说过的那种器皿，施舍给穷人，不要如此疯狂。\par

愿意离开我的人，立刻离开吧；愿意控告我的人，尽管控告，我不会在任何一个人身上容忍。当我在基督的审判台前受审时，你们远远站着，你们的恩惠对我无益，而我正要交账。\Quote{这些话毁了一切！他说：\InnerQuote{不要让他离开，转去另一个教派！}不！他软弱，要迁就他！} 到什么程度？到何时？一次、两次、三次，但不是永远。\par

看！我再次警告你们，并按照有福的保禄的方式作证：\Quote{若我再来，必不宽容。}\BibleRef{格后 13:2} 但当你们做了当做的事，你们就会知道收益多大，好处多大。是的！我恳求并劝勉你们，甚至不拒绝抱你们的膝恳求你们在这件事上。这是何等软弱！何等奢侈！何等放荡！这不是奢侈，而是放荡。何等愚昧！何等疯狂！这么多穷人站在教会周围；尽管教会拥有这么多富有子女，她却无法救济甚至一个穷人；\Quote{这个饥饿，那个醉醺醺}\BibleRef{格前 11:21}；一个人甚至用银器排泄，另一个人却没有面包！何等疯狂！何等兽性超过这个？愿我们永远不必考验是否追究悖逆者，也不必因纵容这些行为而带来的愤怒；但愿我们甘心忍耐，避免这一切，好能为天主的荣耀而生活，免于来世惩罚，获得为爱祂的人所预许的美物，藉着恩宠和爱人之心，等等。\par

\SourceRecord{Translated by John A. Broadus.；Nicene and Post-Nicene Fathers, First Series；Vol. 13.；Edited by Philip Schaff.；Buffalo, NY: Christian Literature Publishing Co.,；1889.；<http://www.newadvent.org/fathers/230307.htm>.}

% structure: p0001 source=h1 target=section reason=page-title

\section{论哥罗森书第八篇讲道词}

\BibleRef{哥 3:5-7}\par

\BibleQuote{你们要致死属于地上的肢体，即淫乱、不洁、邪情、恶欲和贪婪，这贪婪就是拜偶像；为了这些事，天主的忿怒必降在悖逆之子身上；你们从前也生活在其中，当你们在这些事中活着的时候。}\par

我\Emph{知道}许多人被前面的话触怒，但我能做什么呢？你们听见了主的命令。难道是我的错吗？我该怎么办？你们没有看见债主们，当欠债的人顽固不化时，他们是如何戴上颈枷的吗？你们没有听见保禄今天宣告的话吗？他说：\BibleQuote{你们要致死属于地上的肢体，即淫乱、不洁、邪情、恶欲和贪婪，这贪婪就是拜偶像。}有什么比这种贪婪更坏呢？它比任何欲望都更坏。这比我刚才所说的疯狂和关于银钱的愚蠢弱点更为严重。他说：\Quote{这贪婪就是拜偶像。}请看这罪恶的结局。我请求你们，不要误解我所说的话，因为我并非出于好意，也不是无端地要结仇；而是我盼望你们能达到那样的德行，使我能够听到关于你们的我应该听到的事。因此，我这样说并非出于权威，也不是出于专横，而是出于痛苦和忧伤。请你们原谅我，原谅我！我不愿意因谈论这些话题而有失体面，但我被迫如此。\par

我这样说，不是因为穷人的痛苦，而是为了你们的救恩；因为他们将要灭亡，必定灭亡，那些没有喂养基督的人。因为即使你周济了一些穷人，只要你仍过着如此放荡奢侈的生活，一切都没有用。因为要求的不是给予很多，而是不要在你的财产上给得太少；因为那只是儿戏。\par

他说：\BibleQuote{你们要致死属于地上的肢体。}你说什么？不是你自己说：\BibleQuote{你们已埋葬，与他同葬；你们已受割损；我们脱去了肉性罪恶的身体} \BibleRef{罗 6:4}；那么你怎么又说：\Quote{致死}？你在开玩笑吗？你这样说，好像那些事情仍在我们身上？这并不矛盾；就像一个人，他彻底擦净了一个肮脏的雕像，或者更确切地说，重新铸造了它，使它焕然一新，然后说锈已经蚀掉并毁灭了，却又再次建议努力除去锈迹，他并没有自相矛盾，因为他所建议除去的不是他之前擦掉的锈，而是后来新生的锈；同样，他所说的也不是从前的那种致死，也不是那些淫乱，而是后来滋生的那些。\par

他说这不是我们的生命，而是另一个生命，即天上的生命。现在请告诉我：当他说\Quote{致死属于地上的肢体}时，难道地也要受谴责吗？或者他是指在地上的事本身就是罪恶？\par

他说：\Quote{淫乱、不洁。}他略过了那些连提都不合适的行为，而用\Quote{不洁}一词概括了一切。\par

他说：\Quote{邪情、恶欲。}\par

看！他用一个类别概括了一切。因为嫉妒、愤怒、忧伤，都是\Quote{恶欲}。\par

他说：\Quote{贪婪，这贪婪就是拜偶像。因为这些事，天主的忿怒必降在悖逆之子身上。}\par

他用许多事情说服他们：藉着已经赐下的恩惠，藉着将来的灾祸，我们从中被拯救出来，我们是谁，为什么；以及所有那些考量，例如：我们是谁，处在什么境况，我们从中被拯救出来，如何、以何种方式、在什么条件下。这些足以使人转变，但这一点比所有都更有力；虽说起来不愉快，但并非无益，而是有益的。他说：\Quote{因为这些事，天主的忿怒必降在悖逆之子身上。}他没有说\Quote{降在你们身上}，而是\Quote{降在悖逆之子身上}。\par

\Quote{你们从前也生活在其中，当你们在这些事中活着的时候。}为了使他们羞愧，他说：\Quote{当你们在这些事中活着的时候}，并且暗示赞扬，因为现在不再这样生活了：以前他们可以。\par

第8节。\Quote{但如今你们要除去这一切。} \BibleRef{哥 3:8}\par

他总是既普遍又具体地说话；这是因为热忱。\par

第8、9节。\Quote{愤怒、暴怒、恶意、毁谤、出自你们口中的秽语。不可彼此说谎。}\par

他说：\Quote{出自你们口中的秽语，}清楚表明它玷污了口。\par

第9、10节。\Quote{因为你们已经脱去了旧人及其行为，穿上了新人，这新人依照创造他者的形象，在知识上不断更新。}\par

于此值得探究：为何他将败坏的生活称为\Quote{肢体}、\Quote{人}和\Quote{身体}，又同样称德行生活为这些。如果\Quote{这人}意指\Quote{罪}，何以他又说\Quote{与他的行为}？因为他曾说过\Quote{旧人}，表明此非那真人之本，而是另一种。道德抉择比本性更能决定一个人，而且比本性更称得上\Quote{人}。因为本性不将人投入地狱，也不领人进入天国，而是人自己如此；我们也不是因某人是人就爱或恨，而是因他是这样或那样的人。若本性就是身体，且无论哪种都不能算账，那又怎能说它是恶？而他说\Quote{与他的行为}是什么意思？他意指抉择连同行为。他特意称其为\Quote{老}，以显示其丑陋、可憎与软弱；而称\Quote{新}，仿佛在说：不要指望这人与那人相同，而是相反；因为人越前行，不是趋向衰老，而是趋向比之前更年轻的青春。因当他获得更圆满的知识时，他既配得更大的事，也处于更完美成熟、更强健的状态；这不仅由于青春，也由于他所肖似的\Quote{模样}。看！最好的生活被称为新造，是按基督的肖像；因为\Quote{按造祂者的肖像}之意在此，因为基督也最终未衰老，而是美得无法言说。\par

第11节：\BibleQuote{在此处没有希腊人和犹太人、割礼与未受割礼、化外人、西古提人、为奴的、自主的，惟有基督是一切，又在一切之内。}\BibleRef{哥 3:11}\par

看！这是对这\Quote{人}的第三项赞美。在他那里，不区分民族、地位或家世，因他毫无外在之物，也不需要它们；因为一切外在之物就是这些：\Quote{割礼与未受割礼、为奴的、自主的、希腊人}，即归化者，\Quote{和犹太人}，因其祖先。只要你拥有这\Quote{人}，你就必与拥有他的其他人得着同样的事。\par

\Quote{但基督}，他说，\Quote{是一切，又在一切之内}：基督将成为你的一切，既是地位、又是出身，\Quote{并}祂自己\Quote{在你们众人之内}。或者他说另一件事，即你们众人都成了同一基督，成为祂的身体。\par

第12节：\BibleQuote{所以，你们既是天主的选民，圣洁蒙爱的人，就要穿上。}\BibleRef{哥 3:12}\par

他显示德行之轻易，好使他们既能恒常拥有它，又能以此作为最大的装饰。这劝勉也伴随赞美，因为那时其力量最大。因他们先前已是圣洁的，但并非选民；如今则既是\Quote{选民、圣洁、又是蒙爱的}。\par

\Quote{怜悯心肠}。他没有说\Quote{仁慈}，而是用这两个字更加强调。他没有说应如对待弟兄，而是如父亲对待儿女。因为不要对我说他犯罪了，因此他说\Quote{心肠}。他没有说\Quote{同情}，以免看轻他们，而是说\Quote{怜悯心肠、仁慈、谦卑、温柔、忍耐；彼此包容，彼此饶恕，若有人对谁有怨言，正如基督饶恕了你们，你们也要照样行。}\par

他再次按类别说话，并且总是如此；因为从仁慈生出谦卑，从谦卑生出忍耐。他说\Quote{彼此包容}，即放过事情。请看，他如何称之为\Quote{怨言}，并说\Quote{正如基督饶恕了你们}，显明这是小事。榜样何其伟大！他常这样做；他劝勉他们效法基督。他称之为\Quote{怨言}。的确，他以这些话表明这是小事；但当他将榜样摆在我们面前，便说服我们，即使我们有严重的指控，也应饶恕。因为\Quote{正如基督}这个说法正意谓此，不仅如此，还要全心；不仅如此，还应相爱。因为基督被引入中间，带来这一切：即使事情重大，即使我们不是先受害的一方，即使我们为大、他们为小，即使他们事后必定会侮辱我们，我们也应为他们舍命（因为\Quote{正如}一词要求如此）；并且，不仅不停止于死亡，若可能，甚至死后仍要继续。\par

第14节：\BibleQuote{在这一切以上，要穿上爱德，它是全德的联系。}\BibleRef{哥 3:14}\par

你看见他这样说吗？因为人可能饶恕却不爱；是的，他说，你也必须爱他，并且他指出一条使人能够饶恕的途径。因为人可以仁慈、温柔、谦卑、忍耐，却仍不爱。因此他首先说\Quote{怜悯心肠}，即是爱和同情。\Quote{在这一切以上，爱德，它是全德的联系。}他所要说的是：那些事毫无益处，因为除非以爱德行之，那些事都会散开；爱德将它们全部紧束；无论你提及什么善事，若缺少爱德，便算不得什么，它就会消融。就如船上，纵使索具很多，若没有捆绳，也无用处；屋内若没有大梁，也是如此；身体中，纵使骨骼粗大，若没有韧带，也无用处。因为无论人有什么善行，若没有爱德，一切都归于无有。他没有说爱德是顶峰，而是更重要的\Quote{联系}；这比前者更必不可少。因为\Quote{顶峰}确是成全的极致，而\Quote{联系}则是使产生成全的事物彼此紧密结合；它就好比是根。\par

\BibleRef{哥 3:15}\BibleQuote{并且要让天主的平安在你们心中掌管，你们也为此蒙召，在一个身体内；并常怀感恩之心。}\par

\BibleQuote{天主的平安}。这就是那固定而坚定不移的平安。如果你们确实因人的缘故而有平安，它会很快瓦解；但若因天主的缘故，则永不瓦解。虽然他曾普遍地谈论爱，但又回到具体的事上。因为有一种爱也是过度的；例如，由于过度的爱，人无缘无故地指责人，陷入争论，并产生厌烦。他说：我所渴望的不是这个，不是这个；不是过分苛责，而是要如同天主与你们和好一样，你们也要与人和好。他是如何和好的？出于他自己的旨意，没有从你们那里得到任何东西。这是什么？\BibleQuote{让天主的平安在你们心中掌管。}如果两个思想在互相争斗，不要让愤怒、不要让怨恨去夺奖赏，而要平安；例如，假设一个人无故受辱；从侮辱中生起两个思想，一个叫他报复，另一个叫他忍耐；它们彼此角力：如果天主的平安作为裁判站出来，它就会把奖赏赐给那个叫人忍耐的思想，并使另一个羞愧。怎样做？通过说服它：天主是平安，他已与我们和好。他并非没有理由地表明此事中的巨大挣扎。他说：不要让愤怒当裁判，不要让争论，不要让人的平安，因为人的平安来自报复、来自忍受不了可怕的伤害。但我所意指的不是这个，而是他所留下的那一种。\par

他描绘了一个内在的竞技场，在思想中，有竞赛、角力、和裁判。接着又是劝勉：\BibleQuote{你们也为此蒙召}，他说，也就是说，你们为此蒙召。他提醒他们平安是许多好事的因由；为此他召叫了你们，为这个召叫了你们，好使你们获得相称的奖赏。因为为什么他使我们成为\BibleQuote{一个身体}？岂非不正是为了使她掌权吗？岂非不正是为了使我们有和平的时机吗？为什么我们都是一个身体？并且现在我们是一个身体吗？因着平安我们是一个身体，且因着我们是一个身体，我们得以平安。但是为什么他没有说\Quote{让天主的平安得胜}，而是说\Quote{作裁判}？他使她更尊荣。他不愿恶念来与她角力，而是让它处于下方。而\Quote{奖赏}这个名号本身就鼓舞了听者。因为如果她把奖赏赐给了善念，那么无论另一个如何无耻地表现，此后便毫无用处。此外，另一个意识到，无论他做出什么事迹，他都不会得到奖赏；无论他如何膨胀，并尝试更猛烈的进攻，都会停止，因徒劳无功。他还恰当地补充道：\BibleQuote{并常怀感恩之心}。因为这就是感恩，而且非常有效，就是要像天主对待自己一样对待同仆，要顺服于主，服从他；要为一切事表达他的感激，即使有人侮辱他、打他。\par

因为确实，凡承认因自己所受的苦而应感谢天主的人，就不会报复那亏待他的人，因为至少那报复的人不承认感恩。但让我们不要效法那索要一百银钱的人，免得我们听到\BibleQuote{你这恶仆}，因为没有什么比这忘恩负义更糟的了。因此，报复的人是不知感恩的。\par

但是为什么他以邪淫开头他的列表？因为说过\BibleQuote{致死你们在地上的肢体}\BibleRef{哥 3:5}，他立刻说\BibleQuote{邪淫}；他几乎到处都这样做。因为这个情欲有最大的影响力。因为甚至当他写给得撒洛尼人书信时，他也做了同样的事。\BibleRef{得前 4:3}这有什么奇怪？既然甚至对弟茂德他说\BibleQuote{保持纯洁}\BibleRef{弟前 5:22}；又在别处说\BibleQuote{你们要追求与众人和平，并追求圣化，}没有这个\BibleQuote{没有人能看见主。}\BibleRef{希 12:14}\BibleQuote{致死}，他说，\BibleQuote{你们的肢体。}你们知道那死的东西是怎样的，即悦目的、可憎的、腐烂的。如果你们把任何东西致死，它死后不会继续死，而是立刻腐烂，如同身体一样。所以要消灭那灼热；死了的东西不会持续。他表明一个人手里拿着同样的东西，就是基督在圣洗池中所成就的；因此他也称它们为\BibleQuote{肢体}，如同引入一位斗士，从而使他的言论更具强调。他恰当地说\BibleQuote{在地上的，}因为它们在这里持续，并且在这里腐烂，远甚于我们的这些肢体。所以，不像身体那样真正属乎地，罪是属地的，因为前者有时甚至显得美丽，但那些肢体从不。而那些肢体贪恋一切在地上的事物。如果眼睛是这样，它看不见天上的事物；如果是耳朵，如果是手，如果你提到任何其他肢体。眼睛看见身体、美丽和财富；这些都是地上的事物，它以此为乐；耳朵听见柔和的曲调、竖琴、笛子和污秽的言谈；这些都是与地有关的事。\par

当他将听众置于高处、靠近宝座时，他就说：\BibleQuote{治死你们在地上的肢体。}因为带着这些肢体不可能站在高处；那里没有它们可作用的东西。这泥土比那泥土更糟，因为那泥土确实变成金子，\BibleQuote{因为这可朽坏的，}他说，\BibleQuote{必须穿上不可朽坏的}\BibleRef{格前 15:53}，但这泥土绝不能再被重塑。所以这些肢体与其说是\Quote{在地上}的，不如说那些是。因此他没有说\Quote{属地的}，而是说\Quote{在地上}，因为有可能这些肢体不在地上。这些肢体必须\Quote{在地上}，而那些则不必。因为当耳朵听不见这里所说的一切，只听见天上的事，当眼睛看不见这里的一切，只看见上面的事，它就不\Quote{在地上}；当口不说这里的事，它就不\Quote{在地上}；当手不做任何恶事——这些就不属于\Quote{在地上}的事，而属于天上的事。\par

同样，基督也说：\BibleQuote{如果你的右眼使你跌倒，}即如果你不洁地观看，\BibleQuote{把它剜出来}\BibleRef{玛 5:29}，即你的恶念。他（\Person{保禄}）在我看来，把\Quote{淫乱、污秽、情欲、贪欲}视为同一，即淫乱；通过这些表达使我们远离那事。因为实际上这是\Quote{一个情欲}；如同身体易受任何影响，无论是热病还是创伤，这情欲也是如此。他没有说\Quote{抑制}，而是说\Quote{治死}（处死），使它们永不再起，并且\Quote{抛弃它们}。死了的东西，我们就抛弃；例如，如果身上有老茧，它们已经死了，我们就抛弃它。现在，如果你切入活的部分，会引起疼痛；但如果切入死的部分，我们甚至感觉不到。的确，情欲也是如此；它们使灵魂不洁；它们使不朽的灵魂成为可受难的。\par

贪婪如何被称为偶像崇拜，我们已多次解释。因为最支配人类的事物，就是这些：贪婪、不洁和恶欲。\BibleQuote{因这些事，}他说，\BibleQuote{天主的愤怒临到悖逆之子身上。}他称他们为悖逆之子，使他们无可推诿，并表明是因为他们不愿顺从，才处于那种状况。\BibleQuote{你们从前也在这当中，}他说，\BibleQuote{行事为人，}后来成了顺从的。他指出他们仍在这些事中，并称赞他们说：\BibleQuote{但现在你们也要脱去这一切：忿怒、暴怒、恶毒、毁谤、可耻的言谈。}但针对其他人，他展开论述。在\Quote{情欲和毁谤}的名义下，他意指辱骂，正如在\Quote{暴怒}下他意指邪恶。在另一处，为使他們羞愧，他说：\BibleQuote{因为我们彼此互为肢体}\BibleRef{弗 4:25}。他仿佛把他们说成是造人者；丢弃这个，接纳那个。他曾提到人的\BibleQuote{肢体}\BibleRef{哥 3:5}；在这里他说\Quote{一切}。他提到了他的心、忿怒、口、亵渎、眼、淫乱、贪婪、手和脚、说谎、理解本身和旧心。它有一个君王的形像，即基督的形像。他所指的人，在我看来似乎主要是外邦人。因为正如泥土只是沙，即使一部分大一部分小，失去原来的形状后变成金子；又如羊毛，不论何种，接受另一种面貌，隐藏以前的面貌：对信者也是如此。\BibleQuote{彼此容忍，}他说，\BibleQuote{互相宽容}；他表明什么是公平的。你容忍他，他也容忍你；因此他在\BibleBook{迦拉达}{人书}说：\BibleQuote{彼此承担重担}\BibleRef{迦 6:2}。\BibleQuote{并且要感恩，}他说。因为这是他在各处特别寻求的；是万善之首。\par

那么，我们在一切事上都要感恩；无论发生什么事；因为这就是感恩。在顺境中这样做确实不是大事，因为环境本身的性质自然促使人如此；但在极端困境中我们感恩，那才是可称赞的。因为当别人在环境中亵渎并抱怨不平时，我们感恩，请看这里有多么伟大的哲学。首先，你使天主喜悦；其次，你使魔鬼蒙羞；第三，你甚至使所发生的事化为乌有；因为同时，你既感恩，天主就减轻痛苦，魔鬼就离去。因为如果你抱怨不平，他（魔鬼）因得偿所愿，就近你，而天主因受亵渎就离开你，你的灾祸就更严重；但如果你感恩，他一无所获，便离去；而天主因受尊荣，就以更大的尊荣报答你。一个为自己的厄运感恩的人，不可能感受到它们。因为他的灵魂因做正确的事而喜乐；立刻他的良心光明，因自我赞许而欢欣；那光明的灵魂，不可能面带愁容。但在另一种情况下，随着不幸，良心也用她的鞭子攻击他；而在这种情况下，她为他加冕并宣告他。\par

没有比那在患难中感谢天主的舌头更圣洁的了；实在，它丝毫不逊于殉道者的舌头；两者都同样获得冠冕，这一个和那一些。因为在这舌头之上，也有刽子手站着，强迫它藉亵渎来否认天主；魔鬼站在它之上，用刽子手般的思虑折磨它，以沮丧使之黑暗。因此，若有人忍受自己的痛苦而感谢天主，他就获得了殉道的冠冕。例如，她的小孩生病，她感谢天主吗？这对她是冠冕。什么折磨比沮丧更糟糕呢？然而它没有迫使她吐出一句苦毒的话。孩子死了：她又感谢了。她成了亚巴郎的女儿。因为她虽未亲手祭献，却以祭献为乐，这是相同的；当礼物被取走时，她没有感到愤慨。\par

再者，她的孩子生病了吗？她没有制作任何护身符。这在她的决心中被计为殉道，因为她以决心祭献了她的儿子。因为什么？即使那些东西是无效的，纯粹的欺骗和戏弄，但仍有那些说服她相信它们有效的人：她宁可看见她的孩子死，也不愿容忍偶像崇拜。因此，正如她是一位殉道者，无论她是在自己的事上，或是在她儿子的事上如此行；或在她丈夫的，或在她任何亲爱者的事上；同样，那另一个女人是偶像崇拜者。因为显而易见，如果她有机会祭献，她就会施行祭献；是的，甚至现在她已经完成了祭献的行为。因为对于这些护身符，虽然那些借此牟利的人不断为其合理化，并说\Quote{我们呼求天主，并没有做什么特别的事}，等等；还有\Quote{那老妇人是基督徒}，他说，\Quote{而且是信徒之一}；但这事是偶像崇拜。你是信徒之一吗？画十字圣号；说：这是我唯一的武器；这是我的疗法；其他的我一概不知。告诉我，如果一位医生来到某人那里，忽略属于他医术的疗法，而使用咒语，我们能称那人为医生吗？绝不：因为我们看不到医术的药物；同样，在这种情况下，我们也看不到基督信仰的药物。\par

还有其他女人在身上系着河的名字，并冒险做无数类似的事情。看，我说，并预先警告你们所有人，如果有人被查出来，我将不再宽恕他们，无论他们是制作了护身符，还是念了咒语，还是任何其他这类技艺的东西。那么，有人说，孩子会死吗？如果他藉此手段活了下来，他那时就是死了；但若是他没有这个而死了，他反而是活了。但是现在，如果你看到他依附于娼妓，你希望他被埋葬，并说\Quote{为什么，他活着有什么好处？}，但当你看到他在救恩的危险中时，你希望看到他活着吗？你难道没有听见基督说：\Quote{谁丧失自己的生命，必得着生命；谁得着生命，必丧失生命}？\BibleRef{玛 16:25}你相信这些话吗？还是它们在你看来像是寓言？现在告诉我，如果有人对你说：\Quote{把他带到偶像庙里去，他就会活}，你会容忍吗？不！她回答。为什么？\Quote{因为}，她说，\Quote{他催促我犯偶像崇拜；但这里没有偶像崇拜，只是简单的咒语}：这是撒殚的计谋，这是魔鬼的诡诈，用来掩盖欺骗，把有毒的药放在蜜里。当他发现不能用别的方式说服你时，他就走了这条路，转向缝制的符咒和老妇人的寓言；十字架确实受到了羞辱，而这些符咒被置于它之上。基督被赶出去，一个醉醺醺的愚蠢老妇人被带进来。我们的那个奥迹被践踏，而魔鬼的欺骗在跳舞。\par

那么，有人说，为什么天主不谴责来自这些源泉的帮助？祂多次谴责，却没有说服你们；祂现在任凭你们陷于错误，因为经上说：\Quote{天主就任凭他们陷于邪恶的心思。}\BibleRef{罗 1:28}而且，这些事，即使是一个有理智的希腊人也不能容忍。据说雅典有一个煽动者曾把这些东西挂在自己身上：当他的一位老师——一个哲学家——看见后，斥责了他，规劝他，讽刺他，嘲弄他。因为我们竟陷入如此可怜的光景，甚至相信这些东西！\par

为什么，有人说，现在没有那些复活死人、施行治愈的人呢？是的，那么，我说：为什么现在没有那些轻视今生的人呢？我们是为了报酬而事奉天主吗？当人的本性还软弱，当信德尚需被栽种时，就有许多这样的人；但现在祂不要我们依赖这些标记，而要我们准备好面对死亡。那么，你为什么依恋今生呢？为什么你不顾及将来呢？为了今生，你甚至能忍受犯偶像崇拜，但为了那将来，你却连抑制悲伤都做不到吗？正因为这个原因，现在没有这样的人了；因为那（将来）的生命在我们看来是没有尊荣的，因为我们不为它做任何事，而为今生却没有什么是我们不肯承受的。还有，那另一出闹剧——灰烬、烟灰和盐，又是为什么？并且又把那老妇人请进来？真是闹剧和耻辱！然后，\Quote{一只眼睛}，他们说，\Quote{盯上了孩子}。\par

这些撒旦般的行径将止于何处？希腊人如何不会嘲笑？当我们对他们说：\Quote{十字架的大能}时，他们如何不会讥讽？当他们看到我们求助那些连他们自己也嗤之以鼻的事物时，他们如何会被说服？难道天主为此赐予了医生和药物吗？那又如何？假设他们不能治愈他，但孩子离世了？他将往何处去？告诉我，你这可怜可悲的人！他将往魔鬼那里去吗？他将往某个暴君那里去吗？他不往天堂去吗？他不往他自己的主那里去吗？那你为何悲伤？为何哭泣？为何哀悼？为何你爱你的婴儿胜过爱你的主？不正是藉着祂你才有了这个孩子吗？你为何忘恩负义？你爱礼物胜过爱施与者吗？\Quote{但我软弱，}她回答说，\Quote{我受不了对天主的恐惧。}然而，如果肉体上的痛苦，大的掩盖小的，那么在灵魂上更是如此：恐惧消灭恐惧，悲伤消灭悲伤。孩子漂亮吗？但不管怎样，他不会比\Person{依撒格}更美：他也是独子。是在你年老时出生的吗？他也是。但他美丽吗？好吧，不管多美，他不会比\Person{梅瑟}\BibleRef{宗 7:20}更可爱，\Person{梅瑟}曾使外邦人的眼目对他生出温柔的爱，而且这还是在美貌尚未展现的年纪；然而，这可爱的孩子却被父母抛入了河中。你确实看见他陈放，将他埋葬，并去他的坟墓；但他们却不知道他是否会成为鱼、狗或海中其他捕食野兽的食物；他们这样做时，对天国和复活一无所知。\par

但假设他不是独子；而是在你失去了许多孩子之后，这一个也离去了。但你的灾祸不像\Person{约伯}的那样突然，他的更为悲惨？不是房顶塌落，不是他们正在宴乐之时，也不是在传来其他灾祸的消息之时。\par

但你爱这孩子吗？但不会超过\Person{若瑟}，那个被野兽吞噬的；然而，父亲承受了这场灾祸，以及随后的灾祸，再后来的灾祸。他哭了，但没有做出不敬之举；他哀悼，但没有发出不满之言，而是停留在这些话上，说：\Quote{若瑟不在了，西默盎不在了，你们还要带走本雅明吗？这一切都与我作对。}\BibleRef{创 42:36}你看见饥荒的困迫如何使他漠视自己的孩子吗？难道对天主的畏惧不能像饥荒一样支配你吗？\par

哭吧：我不禁止你；但不要说也不要做任何亵渎的事。你的孩子无论怎样，他都不像\Person{亚伯尔}；然而\Person{亚当}没有说过此类的话；虽然那场灾祸是严重的，他的兄弟杀死了他。但我也想起其他杀害兄弟的人；例如，\Person{阿贝沙隆}杀死了长子\Person{阿默农}\EditorialNote{撒慕尔纪下13}，而\Person{达味}王爱他的孩子，并且穿上麻衣，坐在灰堆中，但他没有召来占卜者或行邪术的（虽然当时有这样的人，如\Person{撒乌耳}所显示的），而是向天主恳求。你也要照样做：正如那位义人那样，你也照样做；当你的孩子去世时，说同样的话：\Quote{我要到他那里去，但他不会回到我这里来。}\BibleRef{撒下 12:23}这才是真智慧，这才是爱。无论你多么爱你的孩子，你都不会像他那时爱得那样深。因为即使他的孩子是通奸所生，但那位有福之人对那母亲的爱达到了顶点，而你知道，后代分享父母的爱。他对孩子的爱如此之大，甚至希望他活着，尽管他会成为自己的控告者，但他仍然感谢天主。你认为\Person{黎贝加}经历了什么？当他的兄弟威胁\Person{雅各伯}时，她并没有使她的丈夫忧伤，而是叫他打发她的儿子离开？\BibleRef{创 27:46}当你遭受任何灾祸时，想想比它更糟的；你就会有足够的安慰；并且自问，如果他战死呢？如果死于火中呢？无论我们遭受什么，让我们想想更可怕的事，我们就会有足够的安慰，让我们常常看看那些经历过更可怕事情的人，以及我们自己也受过的更重的灾祸。\Person{保禄}也这样劝勉我们；正如他说的：\Quote{你们与罪恶争斗，还没有抵抗到流血的地步。}\BibleRef{希 12:4}又说：\Quote{你们所遇见的试探，无非是人所能受的。}\BibleRef{格前 10:13}因此，无论我们遭受什么，让我们看看更糟的；（因为我们会找到这样的），这样我们就会感恩。最重要的是，让我们为一切事不断地感恩；因为这样，这些事都会减轻，我们将生活于天主的荣耀中，并获得所预许的美物，愿我们众人藉着恩宠和对人的爱，都能达到，等等。\par

\SourceRecord{Translated by John A. Broadus.；Nicene and Post-Nicene Fathers, First Series；Vol. 13.；Edited by Philip Schaff.；Buffalo, NY: Christian Literature Publishing Co.,；1889.；<http://www.newadvent.org/fathers/230308.htm>.}

% structure: p0001 source=h1 target=section reason=page-title

\section{哥罗森书第九篇讲道}

\BibleRef{哥 3:16-17}\par

\BibleQuote{让基督的话充分地在你们内居住，以各种智慧彼此教导和劝勉，用圣咏、圣诗和属神的歌曲，在你们心里以恩宠向天主歌唱。凡你们在言语或行为上所做的，一切都要因主耶稣的名而做，藉着祂感谢天主父。}\par

劝勉他们要感恩之后，他也指出了道路，就是我最近对你们讲过的。他说了什么？\Quote{让基督的话充分地在你们内居住}；或者不如说，不仅仅是这样，还有另一条路。因为我曾说过，我们应该数算那些经历过更可怕事情的人，以及那些遭受了比我们更痛苦磨难的人，并感谢这些事没有临到我们身上；但他说了什么呢？\Quote{让基督的话居住在你们内}；就是说，那教导、道理、劝勉，其中祂说，现世生命算不了什么，它的福乐也算不了什么。如果我们知道这一点，我们就不会向任何艰难屈服。\BibleRef{玛 6:25}；并且又说，\Quote{因我的愚妄，我的伤口发臭流脓。}\BibleRef{咏 37:5}因为还有什么比一个人，他把自己裹在衣服里，却不关心那些赤裸的兄弟；他喂养狗，却不关心天主的肖像正在挨饿；他虽然深信人性事物是虚无的，却仍像不死的人那样紧抓住它们——更愚妄的呢？既然没有比这样的人更愚妄的，那么也没有比成就德行的人更明智的。因为请看：他是多么明智，有人说。他分施自己的财物，他慈悲为怀，他爱人如己，他好好思考了他们与他具有共同的本性；他好好思考了财富的用途，它不值得看重；他应当怜惜与己同类的身体，胜过怜惜财富。轻视虚荣的人完全是明智的，因为他懂得人间事务；对神性和人性事物的认识，就是哲学。所以，他知道哪些是神性的事物，哪些是人性的事物，并且他使自己远离前者，而致力于后者。他也知道怎样在一切事上感谢天主，他把现世生命看作虚无；因此，他不会因顺境而喜悦，也不会因逆境而悲伤。\par

我恳求你们，不要等待别人来教导你们；你们有天主的圣言。没有人能像它们那样教导你们；因为人常因虚荣和嫉妒而吝啬许多。我恳求你们，凡是关心此生的人，请听，请购买那些会成为灵魂良药的书。如果你们不愿要别的，至少也要得到\WorkTitle{新约}、\WorkTitle{宗徒书信}、\WorkTitle{宗徒大事录}和\WorkTitle{福音书}，作为你们常年的导师。若不幸临到你们，就潜入它们，如同潜入药箱；从中获取你们烦恼的安慰，无论那是损失、死亡、还是亲属的离世；或者不如说，不仅潜入它们，更要完全将它们取来，存在你们心中。\par

这就是一切祸患的原因：不认识圣经。我们手无寸铁就上战场，又怎能安然脱身呢？即使有它们，我们也应该满足于能够安全，更何况没有它们呢？不要把一切全推给我们！你们是羊，但仍然不是没有理性的，而是有理性的；保禄也把许多事托付给你们。那些受教导的人，不是永远在学习；因为那样他们就没有被教导。如果你们永远在学习，你们就永远不会学到。不要这样来，好像你们要永远学习；（因为那样你们永远不会知道；）而是要这样来，为的是完成学习，并教导别人。在技艺中，难道不是所有的人都有固定的时间吗？在科学中，总而言之，在一切技艺中？这样，我们都确定一个已知的时间；但如果你们永远在学习，这正证明你们什么也没有学到。\par

天主曾用这个斥责指责犹太人。\Quote{从母腹中被扶持，直到年老仍受教导。}\BibleRef{依 46:3}如果你们不是一直期待这个，一切就不会这样倒退。假使有些已经学完，而另一些人即将学完，我们的工作就会进步；你们就会既给他人让位，也会帮助我们。告诉我，如果有人去一位文法教师那里，永远在学习字母，他们不是会给他们的老师带来许多麻烦吗？我还要在关于生命的事上对你们讲多久？在宗徒时代，情况并非如此，而是他们不断地从一个地方跳到另一个地方，指派那些最先学习的人去教导任何其他正在受教导的人。这样，他们就能环绕世界，因为不受限于一个地方。你们认为，你们在乡间的弟兄们和他们的导师需要多少教导？但你们把我固定在这里。因为，在头尚未摆正之前，就去处理身体的其他部分是多余的。你们把一切都推给我们。你们应该只从我们这里学习，你们的妻子从你们学习，你们的孩子从你们学习；但你们把一切都留给我们。因此，我们的辛劳是过度的。\par

他说：\Quote{教导}，\Quote{并以圣咏、赞美诗与属神的歌曲彼此劝勉}。也要注意\Person{保禄}的体贴。既然读书劳神，烦闷甚大，他并未将他们引向历史，而是引向圣咏，好让你一面以咏唱愉悦心灵，一面温和地消除疲劳。他说：\Quote{赞美诗}和\Quote{属神的歌曲}。但现在你的孩子们会唱出\Person{撒殚}的歌曲和舞蹈，如同厨子、酒席承办人和乐师；没有人知道任何圣咏，反而觉得那是一件可耻、可笑、戏谑的事。这是所有这些邪恶的宝库。因为植物站在何种土壤中，就结何种果实；若在沙质咸土中，果实也是如此；若在甜美沃土中，果实也类似。所以教导的材料就像一口泉源。教他唱那些充满爱智之心的圣咏：比如有关贞洁，或者更确切地说，首先是不要与恶人交往，正如书卷一开始（因此先知就这样开始，说：\Quote{不随从恶人的计谋的人是有福的}；\BibleRef{咏 1:1}，又说：\Quote{我没有坐在虚妄的会议中}；\BibleRef{咏 25:4}，又说：\Quote{在他眼中，作恶的人被藐视，他却尊敬敬畏主的人}；\BibleRef{咏 14:4}）关于与善人交往（这些主题你在那里会找到很多），关于克制肚腹、克制手、避免过度、不欺压人；金钱算不了什么，荣耀也算不了什么，以及其他类似的事。\par

当你从小在这事上引导他，渐渐地你会带领他前进到更高的事。圣咏包含一切，但赞美诗却毫无人性。当他受圣咏教导后，他也会认识赞美诗，如同更神圣的事。因为上层的诸军赞美的是赞美诗，而非圣咏。因为有人说：\Quote{赞美诗}，\Quote{在罪人口中不相宜}\BibleRef{德 15:9}；又说：\Quote{我的眼目要眷顾这地的忠实者，使他们与我同坐}\BibleRef{咏 100:6}；又说：\Quote{行事骄傲的，未曾住在我家中}；又说：\Quote{行走正道的人，他服侍了我。}\BibleRef{咏 100:6}\par

因此，你应稳妥地保护他们，不让他们与朋友甚至仆役混合。因为当我们把腐败的奴隶放在自由人之上时，对自由人造成的伤害不可估量。因为如果他们在享有父亲慈爱和智慧的全部恩惠时，尚且难以保全；当我们把他们交给仆役的无情时，仆役就像对待敌人一样对待他们，以为当他们使他们成为完全的愚人、软弱、不值得尊重时，就会成为他们比较温和的主人。\par

比所有其他事更甚，让我们认真关注这一点。他说：\Quote{我爱慕祢法律的人。}\BibleRef{咏 119:165}因此，让我们也效法这人，并爱这样的人。为使青年进一步受训于贞洁，让他们听先知说：\Quote{我的腰充满幻象}\BibleRef{咏 37:7}；又让他们听他说：\Quote{所有背离祢行淫的，祢必尽行毁灭。}\BibleRef{咏 72:27}为知道应当抑制肚腹，让他们再听：\Quote{当他们口中的食物尚在时，他击杀了他们中的多数。}\BibleRef{咏 77:30}为知道他们应高于贿赂：\Quote{若财富增多，不可将心放在其上}\BibleRef{咏 61:10}；为知道他们应制服虚荣：\Quote{他的荣耀也不会与他一同下降。}\BibleRef{咏 48:17}为不要嫉妒恶人：\Quote{不要嫉妒作恶的人。}\BibleRef{咏 36:1}为视权势如无物：\Quote{我看见恶人高升，如黎巴嫩的香柏树自高自大，我经过，看哪，他已不在了。}\BibleRef{咏 36:35}为视现世之事如无物：\Quote{他们以为那国的人有福；那以上主天主为助手的民族有福。}\BibleRef{咏 143:15}为知道我们并非在不知情中犯罪，而是会有报应，因为他说：\Quote{祢必照各人的行为报应他。}\BibleRef{咏 61:12}但为何他不天天这样报应？他说：\Quote{天主是审判者，正义而强有力，且缓于发怒。}\BibleRef{咏 7:11}为知道谦卑的心是好的，他说：\Quote{主啊，我的心不骄矜}\BibleRef{咏 130:1}；为知道骄傲是恶的，他说：\Quote{因此，骄傲完全抓住了他们}\BibleRef{咏 72:6}；又说：\Quote{主抵抗骄傲的人}；又说：\Quote{他们的不义如同脂油涌出。}为知道施舍是好的：\Quote{他散财，施舍给穷人，他的正义永远常存。}\BibleRef{箴 3:34}为知道怜悯是值得称赞的：\Quote{仁慈待人、借贷与人的人是好人。}\BibleRef{咏 72:7}你在那里还会找到许多比这些更充满真正哲理的道理：比如，不应说恶言：\Quote{暗中诽谤邻舍的人，我要驱逐他。}\BibleRef{咏 111:9}\par

上面的圣歌是什么？信友知道。上面的\OriginalTerm{革鲁宾}{cherubim}说什么？天使说什么？\BibleQuote{天主在天受光荣。}\BibleRef{咏 111:5}因此，圣咏之后是圣歌，作为更完善的事。他说：\Quote{以圣咏，}他说，\Quote{以圣歌，以属神的歌曲，怀着恩宠在心内歌颂天主。}\BibleRef{咏 100:5}他可能指：天主因恩宠赐给我们这些；或恩宠中的歌曲；或彼此在恩宠中劝勉和教导；或他们拥有恩宠中的这些恩赐；或是一个\OriginalTerm{解释性说明}{epexegesis}，指来自圣神的恩宠。\BibleQuote{在心内歌颂天主。}不是简单地用口，他意指，而是用心留意。因为这是\Quote{歌颂天主，}而那是向空气歌唱，因为声音散失无果。他意指不为炫耀。即使你在市场中，你也能收敛自己，歌颂天主，无人听见。因为\Person{梅瑟}也曾这样祈祷，并蒙垂听，因为祂说：\BibleQuote{你为什么向我哀号？}\BibleRef{出 14:15}虽然他什么也没说，但在心意中哀号——因此唯独天主听见了他——怀着痛悔的心。因为即使行走时，在心内祈祷并思念天上的事，也并非禁止。\par

第17节。他说：\BibleQuote{凡你们无论做什么，在言语上或在行为上，一切都因主耶稣之名而行，借着祂感谢天主圣父。}\BibleRef{哥 3:17}\par

因为如果我们这样做，就不会有任何污秽，任何不洁，只要基督被呼求。无论你吃，你喝，你结婚，你旅行，一切都因天主之名而行，即呼求祂帮助你：在一切事上先祈祷祂，然后开始你的事务。你要说什么？把这放在前面。因此，我们在书信前面也放上主的名。无论天主之名在哪里，一切都是吉利的。因为若执政官的名字使文书可靠，基督之名更加如此。或者他意指：说和做一切事都要奉天主，不要引入天使之外。你吃吗？在前后都感谢天主。你睡吗？在前后都感谢天主。你进入集会吗？同样如此——没有世俗的事，没有此生的事。一切都因主之名而行，一切都将为你顺利。无论那名被放在哪里，那里一切都吉利。如果它驱逐魔鬼，如果它驱散疾病，那就更使事务容易了。\par

那么，什么是\Quote{在言语上或在行为上}？无论是要求还是执行任何事情。听，\Person{亚巴郎}如何因天主之名派遣他的仆人；\Person{达味}因天主之名杀死了\Person{哥肋雅}。祂的名奇妙而又伟大。再者，\Person{雅各伯}派遣他的儿子们时说：\BibleQuote{愿我的天主使你们在那人面前蒙恩。}\BibleRef{创 43:14}因为这样做的人有天主作为他的盟友，没有祂他不敢做任何事。因此，既然被呼求而受尊敬，祂将反过来使他们的业务顺利而尊敬他们。呼求圣子，感谢圣父。因为当圣子被呼求时，圣父也被呼求；当祂被感谢时，圣子也被感谢。\par

这些事我们不仅要学习言语，也要以行动完成。没有比这名更伟大的，它在各处都是奇妙的。他说：\BibleQuote{你的名是倾泻的香膏。}\BibleRef{歌 1:3}谁说出它，立刻充满馨香。经上说：\BibleQuote{除非受圣神感动，没有一个人能说\InnerQuote{耶稣是主}。}\BibleRef{格前 12:3}这名做如此伟大的事。如果你带着信心说：\Quote{因父、及子、及圣神之名}，你就完成了所有事。看，你成就了何等伟大的事！你创造了一个人，并成就了圣洗所带来的所有其余部分！因此，当用于命令疾病时，这名是可怕的。因此，魔鬼引入了那些天使，嫉妒我们得荣耀。这样的咒语属于魔鬼。即使它是天使，即使是\OriginalTerm{总领天使}{archangel}，即使是\OriginalTerm{革鲁宾}{cherubim}，也不要允许；因为这些掌权者也不会接受这样的称呼，当他们看到自己的主人受羞辱时，甚至会把它们扔开。他说：\Quote{我荣耀了你，并说：呼求我；}而你却羞辱祂？如果你怀着信德吟诵这咒语，你将驱散疾病和魔鬼，即使你没有驱走疾病，这也不是由于缺少能力，而是因为这样更合适。他说：\BibleQuote{按你的伟大，你的赞美也同样。}\BibleRef{咏 47:10}藉这名，世界被转变，暴政被瓦解，魔鬼被践踏，天被打开。我们因这名而重生。如果我们拥有它，我们便发光；这名称产生殉道者和宣信者；我们要紧握它如同伟大的恩赐，使我们生活在光荣中，中悦天主，并配得给爱祂的人所预许的美善，借着恩宠和慈爱，等等。\par

\SourceRecord{Translated by John A. Broadus.；Nicene and Post-Nicene Fathers, First Series；Vol. 13.；Edited by Philip Schaff.；Buffalo, NY: Christian Literature Publishing Co.,；1889.；<http://www.newadvent.org/fathers/230309.htm>.}

% structure: p0001 source=h1 target=section reason=page-title

\section{哥罗森书讲道集 第十篇}

哥罗森书 3:18–25 \BibleRef{哥 3:18-25}\par

\BibleQuote{\Emph{妻子，要服从丈夫，如同在主内是适宜的。}\Emph{丈夫，要爱妻子，不要对她们怀恨。}\Emph{儿女，要在一切事上听从父母，因为这是主所喜悦的。}\Emph{父亲，不要激怒儿女，免得他们灰心丧气。}\Emph{仆人，要在一切事上服从你们肉身的主人，不要只在眼前事奉，像讨人喜欢的人，却要以纯朴的心，敬畏主。}\Emph{无论做什么，都要从心里去做，如同为主，而不是为人，}\Emph{因为知道你们从主那里必得着基业为赏报；你们服事的是主基督。}\Emph{凡行不义的，必得到他行为的报应，而且主是不偏待人的。}\BibleRef{哥 4:1}\Emph{主人，要公正公平地对待仆人，因为知道你们在天上也有一位主人。}}\par

为什么他不是在每处、在所有书信中给出这些命令，而只是在这里，在给\BibleBook{厄弗所}{书}、\BibleBook{弟茂德}{书}、\BibleBook{弟铎}{书}中？大概是因为在这些城市中有纷争；或者可能他们在其他方面是正确的，所以听到这些事对他们有益。然而，他对这些人说的话，也是对所有人说的。现在，在这些事上，这封书信与\BibleBook{厄弗所}{书}很相似，要么是因为不宜给那些已享平安、尚需在高深教义上受教导的人写这些事，要么是因为那些在试炼中受安慰的人，再听这些话题是多余的。所以我猜想，在此处，教会已经根基稳固，这些事是作为补充而说的。\par

三章18节。\BibleQuote{妻子，要服从丈夫，如同在主内是适宜的。}\BibleRef{哥 3:18}\par

就是说，要为了天主而服从，因为这装饰了你们，他说，不是装饰他们。我不是指那种对主人的服从，也不是单指本性的服从，而是为了天主的服从。\par

三章19节。\BibleQuote{丈夫，要爱妻子，不要对她们怀恨。}\BibleRef{哥 3:19}\par

看，他又一次敦促了互惠。如同在另一种情形中他吩咐了畏惧和爱，这里也是如此。因为即使爱一个人，也可能怀恨。他所说的是：不要争斗；因为没有什么比争斗更痛苦，尤其是丈夫对妻子的争斗。因为相爱之人之间的争斗是痛苦的，他表明，当一个人与自己的肢体不和时，就源于极大的痛苦。所以爱是丈夫的本分，顺从是另一方的本分。如果各尽己责，一切就稳固。因着被爱，妻子也变成有爱；因着服从，丈夫也变得顺从。看，在自然中也是如此安排的：一方应爱，另一方应服从。当治理的一方爱被治理的一方时，一切都稳固。来自被治理者的爱并不如此必要，而来自治理者对被治理者的爱却是必要的；因为从另一方要求的是服从。因为女人有美貌，男人有欲望，这无非表明，为了爱而如此安排。所以，不要因为你的妻子服从你，你就专横；也不要因为你丈夫爱你，你就自高。不要让丈夫的爱使妻子得意，也不要让妻子的服从使丈夫自傲。为此，他把妻子置于你之下，好使她更被爱。为此，他使你被爱，噢妻子，好使你容易承受你的服从。不要害怕服从；因为被爱者的服从没有困难。不要害怕去爱，因为你拥有她的顺从。再没有别的方式能形成纽带。你们有源自本性、必要的权威，也要保持源自爱的纽带，因为这使得软弱变得可忍受。\par

三章20节。\BibleQuote{儿女，要在一切事上听从父母，因为这是主所喜悦的。}\BibleRef{哥 3:20}\par

他又一次加上了\BibleQuote{在主内}，一面立下服从的律法，一面使他们惭愧并降服。因为这，他说，是主所喜悦的。看，他如何希望我们做一切事，不仅出于本性，更先于本性，出于使天主喜悦的事，好使我们也能得到赏报。\par

三章21节。\BibleQuote{父亲，不要激怒儿女，免得他们灰心丧气。}\BibleRef{哥 3:21}\par

看！这里同样有服从和爱。他没有说\BibleQuote{爱你们的儿女}，因为那是多余的，因为本性本身迫使如此；但他纠正了需要纠正的事，使爱在此处也更加强烈，因为服从更大。因为圣经没有以夫妇关系为例；而是什么？听先知说：\BibleQuote{如同父亲怜恤儿女，主也怜恤敬畏祂的人。}\BibleRef{咏 102:13} 基督又说：\BibleQuote{你们中间有谁，儿子向他求饼，反而给他石头？或者求鱼，反而给他蛇？}\BibleRef{玛 7:9}\par

\BibleQuote{父亲，不要激怒儿女，免得他们灰心丧气。}\par

他写下了他知道最能抓住他们的事；在命令他们时，他说话更像朋友；他完全没有提到天主，因为他想胜过父母，软化他们温柔的情感。就是说，\Quote{不要使他们更好争辩，有时候你们甚至应当让步。}\par

接下来，他谈到第三种权威。\par

这里也有某种爱，但不再是出于本性，像上面那样，而是出于习惯，出于权威本身和所做的事。既然在此处爱的范围缩小了，而服从的范围扩大了，他就详细论述这一点，希望从服从上给他们像前者从本性上得到的一样。因此，他单独对仆人所说的，不只是为了主人的缘故，也是为了他们自己的缘故，好使他们使自己成为主人温柔关爱的对象。但他没有明说；因为那样无疑会使他们懈怠。\par

Ver. 22. \Quote{僕人，}他說，\Quote{要凡事順從你們肉身的主人。}\par

請注意，他如何總是列出這些稱呼：\Quote{妻子、兒女、僕人，}既是對他們順從的正當要求。但為免有人感到痛苦，他補充說：\Quote{對你們肉身的主人。}你們更好的部分，即靈魂，是自由的，他說；你們的服務是暫時的。因此你要使肉身服從，好使你的服務不再出於被迫。\Quote{不要只在眼前事奉，像是討人喜歡的。}他說，要使你們按律法所作的服務出於對基督的畏懼。因為如果你的主人看不見你時，你仍盡本分並做榮耀他的事，顯然你是因那永不睡眠的眼睛而做。\Quote{不要只在眼前事奉，}他說，\Quote{像是討人喜歡的；}這話暗示：\Quote{受損害的將是你們自己。}因為請聽先知說：\Quote{天主驅散了討人喜歡者的骨頭。}\BibleRef{詠 52:6}請看他如何寬容他們，並將他們納入秩序。\Quote{卻要以心誠，}他說，\Quote{敬畏天主。}因為那不是誠實，而是虛偽：心口不一；主人無論在場不在場，都表現一致。因此他不僅說\Quote{以心誠，}還說\Quote{敬畏天主。}因為敬畏天主就是這樣：即使沒有人看見，我們也不做任何邪惡的事；但如果我們做了，我們不是敬畏天主，而是敬畏人。你看出他如何將他們納入秩序了嗎？\par

Ver. 23. \BibleQuote{你們無論做什麼，都要從心裡做工，如同為主，而不是為人。}\BibleRef{哥 3:23}\par

他渴望不僅使他們脫離虛偽，也脫離懶惰。他使他們不再是奴隸，而是自由人，因為他們不需要主人的監督；因為\Quote{從心裡}一詞的意思是：\Quote{出於好心好意}，不是出於奴隸般的不得已，而是出於自由和選擇。那報酬是什麼呢？\par

Ver. 24. \BibleQuote{知道，}他說，\BibleQuote{你們將從主那裡得到基業的賞報；因為你們事奉主基督。}\BibleRef{哥 3:24}\par

因為顯然你們也將從祂那裡得到賞報。而且你們事奉主，從此也可見。\par

Ver. 25. \BibleQuote{因為行不義的，}他說，\BibleQuote{必受他所行不義的報應。}\BibleRef{哥 3:25}\par

他在此確認了先前的話。為使他的話不顯得是奉承，他說：\BibleQuote{他必受他所行不義的報應，}亦即，他也將受懲罰，\BibleQuote{因為天主不徇情面。}因為即使你是僕人，這對你也不是羞恥。事實上，他本可對主人說這話，如同他在《致厄弗所人書》中所做的。\BibleRef{弗 6:9}但在這裡，他似乎暗示希臘主人。因為即使他是希臘人而你是一個基督徒，受審查的不是人，而是行為，因此在這種情況下，你也應該出於好心好意、誠心誠意地事奉。\par

第四章，Ver. 1. \BibleQuote{主人，要將正義與公平施行於你們的僕人。}\BibleRef{哥 4:1}\par

什麼是\Quote{正義}？什麼是\Quote{公平}？就是使他們在一切事上充足，不讓他們需要依靠別人，並為他們的勞作給予報償。因為，既然我說他們從天主那裡得到賞報，你不要因此剝奪他們的賞報。在另一處他說：\BibleQuote{含忍恐嚇}\BibleRef{弗 6:9}，希望使他們更溫和；因為那些人是成全的；也就是說：\BibleQuote{你們用什麼量器量給人，也必用什麼量器量給你們。}\BibleRef{瑪 7:2}而\Quote{不徇情面}這話是對這些主人說的，但卻歸於那些僕人，好使這些主人接受它。因為當我們對一個人說了適用於另一個人的話時，我們糾正的不是那個人，而是那個有過錯的人。\Quote{你們也}，他與他們一同說。他在此使服務成為共同的，因為他說：\Quote{知道你們也有一位主人在天上。}\par

Ver. 2. \BibleQuote{要恆心祈禱，在其中醒寤，並懷著感恩。}\BibleRef{哥 4:2}\par

因為，恆心祈禱常使人倦怠，因此他說\Quote{醒寤}，即清醒，不游移。因為魔鬼知道，它知道祈禱有多大益處，因此它重重壓迫。保祿也知道許多人在祈禱時多麼漫不經心，因此他說\Quote{恆心}祈禱，如同說一件有些勞苦的事，\BibleQuote{在其中醒寤，並懷著感恩}。因為他說，讓這成為你們的工作：在祈禱中感恩，為所見和未見的、為祂對願意和不願意的人的恩惠、為天國、為地獄、為磨難、為安慰。因為聖人們習慣這樣祈禱，並為眾人的共同恩惠而感恩。\par

我知道一位圣人，他如此祈祷。他在这些话之前通常什么也不说，而是这样说：\Quote{我们感谢祢，为祢赐予我们不配之人的一切恩惠，从起初直到如今，为我们知道的、为我们不知道的，为看得见的、为看不见的，为行为上的、为言语上的，为我们意愿的、为我们不愿的，为所有赐予不配之人，甚至我们的一切恩惠；为困苦，为安慰，为地狱，为惩罚，为天国。我们恳求祢，使我们的灵魂圣洁，拥有纯洁的良心；赐予我们一个堪当祢慈爱的结局。祢曾如此爱我们，甚至为我们赐下祢的独生子，求祢赏赐我们堪当承受祢的爱情；在祢的言语和祢的敬畏中赐予我们智慧。独生子基督，求祢激发来自祢的力量。祢曾为我们赐下独生子，并派遣祢的圣神来赦免我们的罪，若我们在任何事情上故意或无意犯了过犯，求祢宽恕，不要归咎。求祢记念所有真实呼求祢圣名的人；记念所有善待我们或反对我们的人，因为我们都是人。}然后他加上了信友祷词，就在那里结束；将那祈祷作为某种顶峰和统合一切的纽带。因为天主赐予我们许多恩惠，甚至违背我们的意愿；还有许多，甚至更多，是在我们不知不觉中。因为当我们为某事祈祷，而祂却为我们做相反的事，显然祂是在我们不知道时仍然施恩于我们。\par

第3节。\BibleQuote{同时也为我们祈祷。}\BibleRef{哥 4:3}请看他的谦卑；他把自己置于他们之后。\par

\Quote{愿天主为我们打开传道的门，宣讲基督的奥秘。}他意指入门和讲道的胆量。奇妙！这位伟大的运动员没有说\Quote{使我脱离锁链}，而是身在锁链中却劝勉他人；他劝勉他们为了一个伟大的目标，就是他自己能够获得讲道的胆量。这两者都是伟大的，无论是人的品质还是事情的品质。奇妙！这尊严何其大！他说：\Quote{基督的}，\Quote{奥秘。}他表明没有比这更令他渴望的，就是讲道。\BibleQuote{我也为此锁链，使我能够显明它，按我所应当说的。}\BibleRef{哥 4:4}他意指极大胆的言论，毫不保留。他的锁链显明了他，而非遮蔽他。他意指极大胆。请告诉我，你身陷锁链，却劝勉别人？是的，我的锁链给了我更大的胆量；但我祈求天主的帮助，因为我曾听见基督的声音说：\BibleQuote{当他们把你们交付时，不要忧虑怎样或说什么。}\BibleRef{玛 10:19}并且看，他如何用比喻表达自己，\Quote{愿天主为我们打开传道的门。}（看，他多么谦卑；即使在锁链中，他如何表达自己）这就是说，求祂软化他们的心。但他没有这样说，而是说\Quote{求祂赐我们胆量}；出于谦卑他这样说了，他所拥有的，他求着领受。\par

他在此书信中表明，基督为何不在那些时代来临，因为他称以前的事物为\Quote{阴影，但实体}，他说，\Quote{属于基督。}所以他们必须在阴影下养成习惯。同时他也显明了对他们的爱的最伟大证明：他说，\Quote{为了使你们，}因此，\Quote{听到，\InnerQuote{我身陷锁链。}}他再次将他的那些锁链摆在我们面前；我非常喜爱它们，它们振奋我的心，总是吸引我渴望看见保禄被捆绑，并在他的锁链中写作、宣讲、施洗和教导。他在锁链中为各处教会代求；他在锁链中建立的工程无可估量。那时他反而更自由。因为听他说道：\BibleQuote{因此多数弟兄在主内因我的锁链而信赖，更加勇敢地无所畏惧地传讲天主的话。}\BibleRef{斐 1:14}他又再次如此声称自己，说：\BibleQuote{因为我软弱时，正是我刚强时。}\BibleRef{格后 12:10}因此他也曾说：\BibleQuote{但天主的话并不被束缚。}\BibleRef{弟后 2:9}他与歹徒、囚犯、凶杀者一同被捆绑；他，世界的导师，曾被提到第三层天上，听见不可言传的话语，竟然被捆绑。\BibleRef{格后 12:4}但他的道路反而更加迅速。被捆绑的，现在得以释放；未被捆绑的，却被捆绑。因为他确实在做他所愿的事；而对方既不能阻止他，也不能达成自己的目的。\par

你在做什么，愚昧的人？你以为他是属肉体的赛跑者吗？他在我们的赛场上竞争吗？他的生活轨迹在天上；奔跑在天上的人，地上的事物不能捆绑或阻挡他。你看不见这太阳吗？用镣铐锁住它的光芒！阻止它的轨道！你不能。那么你也不能阻挡保禄！是的，这位比那一位更甚，因为这位比那位更享有天主的眷顾，因为他带给我们不是那样的光，而是真光。\par

如今那些不愿为基督承受任何苦难的人在哪里？但为什么我说\Quote{受苦}，因为他们甚至不愿放弃自己的财富？过去保禄也曾捆绑人，投入监狱；但自从他成为基督的仆人，他不再夸耀所作，而夸耀所受的。而且，这在宣讲中更是奇妙，当它被受苦者自己高举和增长时，而非被迫害者。哪里见过这样的竞赛？遭受恶的人得胜；行恶的人失败。这人比那人更光明。通过锁链，宣讲进入了。\BibleQuote{我不以福音为耻}\BibleRef{罗 1:16}，是的，他甚至夸耀宣讲被钉者。因为请想想：全世界离开了那些自由的人，转而投向那些被捆绑的人；远离监禁者，尊荣负枷者；憎恨钉死基督者，敬拜被钉者。\par

不仅是讲道者是渔夫、没有学识这件事令人惊叹；还有其他的阻碍，甚至是本性上的阻碍；但增长反而更加丰盛。不仅他们的无知不是阻碍；甚至它本身就彰显了宣讲。因为听\Person{路加}说：\BibleQuote{他们看出他们是没有学识的平常人，就希奇。}\BibleRef{宗 4:13}不仅是锁链不是阻碍，甚至它们本身就使他们更加大胆。\Person{保禄}在自由时，门徒们并没有像他被捆绑时那样勇敢。因为他说，他们\BibleQuote{他们更放胆传讲天主的话，没有惧怕}\BibleRef{斐 1:14}那些反对宣讲之神性的人在哪里？难道他们的无知不足以使他们被定罪吗？难道在这种情况下，它不会吓倒他们吗？因为你知道，大多数人被这两种情感所支配：虚荣和胆怯。假设他们的无知不让他们感到羞耻，但危险一定会使他们害怕。\par

但是，有人说，他们行了奇迹。那么你相信他们行了奇迹。但是难道他们没有行奇迹吗？若人们没有奇迹而被他们吸引，这比行奇迹更大。在希腊人中，\Person{苏格拉底}也曾被下在锁链里。结果怎样？他的门徒不是立刻逃到\Place{默伽拉}去了吗？当然，为什么不呢？他们接受了他关于不朽的论证。但是请看这里。\Person{保禄}被下在锁链里，他的门徒却更加胆壮，这是有理由的，因为他们看见宣讲并未被阻止。因为，你能把舌头锁住吗？正是藉着锁链，宣讲奔腾。因为，正如你若没有绑住跑者的脚，你并没有阻止他奔跑；同样，你若没有绑住传福音者的舌头，你就没有阻止他奔跑。而且，如同前者，若你绑住他的腰，他反而跑得更快，得到支持；同样，后者也宣讲得更多，更加大胆。\par

一个囚犯只有在锁链之外别无他物时才会恐惧；但一个轻看死亡的人，怎能被捆绑呢？他们所做的，就好像是把\Person{保禄}的影子锁住，堵住它的嘴。因为这是一场与影子的搏斗；因为他既被朋友们更温柔地怀念，又被敌人更尊敬地看待，如同在锁链中戴着勇气的奖赏。冠冕束缚着头，但不会使它蒙羞，反而使之璀璨。他们违反他的意愿，用锁链为他加冕。因为，告诉我，一个敢于闯入死亡之金刚石大门的人，难道会害怕铁吗？来吧，亲爱的，让我们效法这些锁链。你们当中那些用金饰装扮自己的妇女们，渴望\Person{保禄}的锁链吧。这些铁链的恩宠围绕着他的灵魂闪耀，远非你们颈上的项链所能比！如果有人渴望那些锁链，就让他恨这些吧。因为柔顺与勇气有什么相通？身体的装饰与哲学有什么相干？那些锁链，天使尊敬；这些锁链，天使甚至嘲笑。那些锁链惯于从地上提升到天上；这些锁链则从天上拖到地上。因为事实上，这些才是锁链，那些不是；那些是装饰，这些是锁链；这些，连同身体，也折磨灵魂；那些，连同身体，也装饰灵魂。\par

你要信服那些是装饰吗？告诉我，谁会更引起观众的注意？你还是\Person{保禄}？我为什么说\Quote{你}？那满身金饰的王后本身也不会如此吸引观众；但如果\Person{保禄}戴着锁链和王后同时进入教会，所有人都会把目光从她转向他；而且有充分的理由。因为看见一个本性超越人类、毫无人的成分、是地上天使的人，比看见一个穿着华服的女人更令人赞叹。因为这样的景象，一个人在剧院、游行、浴场和许多地方都可以见到；但是，谁若看见一个戴着锁链的人，并且认为自己拥有最大的装饰，不在锁链之下屈服，他所见的不是地上的景象，而是配得上天堂的景象。那以那种方式装饰的灵魂环顾四周——谁见过？谁没见过？——充满了骄傲，被焦虑的思想所占据，被无数其他情欲所束缚；但是那个戴着这些锁链的人，没有骄傲：他的灵魂欢跃，摆脱了一切焦虑，喜乐，目光凝视天堂，披上了翅膀。如果有人让我选择看见\Person{保禄}从天国俯身说话，还是从狱中说话，我宁愿选择狱中。因为当他身在狱中时，天上的使者来拜访他。\Person{保禄}的锁链是宣讲的锁链，那锁链是宣讲的基础。让我们渴望那些锁链！\par

有人会说：这怎么可能呢？如果我们打碎并销毁这些饰品，对我们并无益处，反而有害。这些束缚使我们在那里显为囚犯；但\Person{保禄}的锁链会解开那些束缚；凡在此地被这些束缚捆绑的人，也要在那里被那些不朽的锁链捆绑，双手双脚；凡被\Person{保禄}的锁链所捆绑的人，在那一天，那些锁链将如同饰品环绕她。解除你自身的束缚，也解除穷人的饥饿。为何你要钉牢你罪孽的锁链？有人问：何以如此？当你佩戴黄金而另一个人濒临死亡，当你为了虚荣而获取如此多的黄金，而另一个人连吃的都没有，你不正是把你的罪孽钉得更牢固了吗？让\Person{基督}围绕你，而非黄金；哪里有\OriginalTerm{玛门}{Mammon}，哪里就没有\Person{基督}；哪里有\Person{基督}，哪里就没有\OriginalTerm{玛门}{Mammon}。难道你不愿穿上万有的君王本身吗？如果有人给你紫袍和王冠，你难道不是先于世上所有的黄金而接受它们吗？我给你的不是君王的饰品，而是邀请你穿上君王本身。有人问：人如何穿上\Person{基督}呢？请听\Person{保禄}所说：\BibleQuote{凡你们受洗归于基督的，就是穿上了基督。}\BibleRef{迦 3:27} 请听宗徒的训诫：\BibleQuote{不要为肉体安排，去放纵私欲。}\BibleRef{罗 13:14} 这样，人若不为肉体安排去放纵私欲，就是穿上了\Person{基督}。如果你穿上了\Person{基督}，连魔鬼也要畏惧你；但如果你穿上黄金，连人也会嘲笑你；如果你穿上了\Person{基督}，人也会尊敬你。\par

你想显得美丽动人吗？当以造物主的塑造为满足。为何你要覆盖这些金片，仿佛要修正天主的创造？你想显得动人吗？穿上施舍，穿上慈善，穿上谦逊、谦卑。这些都远比黄金更珍贵；它们使美丽者更加动人；它们使形体不佳者变得秀美。因为当人善意地看一张面容时，他从爱情出发作判断；但一个邪恶的女人，即使美丽，也无人能称她为美丽；因为心志被扰乱，无法正确判断。\par

古时的那位埃及女子打扮得华丽；\Person{若瑟}也同样打扮；他们中谁更美丽呢？我不是说她在宫中而他在狱中时。他赤身，却穿着贞洁的衣服；她穿着衣服，却比赤身更不雅观，因为她没有谦逊。妇人啊，当你过分装饰自己时，你比赤身更加不雅；因为你剥去了自己美丽的装饰。\Person{厄娃}也曾赤身；但她穿上衣服后，却更加不雅，因为她赤身时，确实有天主的光荣作装饰；但当她穿上罪恶的衣服时，就不雅了。而你呢，当你穿着精心修饰的华服时，就显得更加不雅。因为昂贵并不能使人显得美丽，而且穿着盛装的人甚至可能比赤身更不雅，请告诉我；如果你穿上吹笛人或吹箫者的衣服，这岂不有失体统？然而那些衣服是金制的；但正因如此才是不雅，因为它们是金的。因为昂贵衣着适合舞台上的人，如悲剧演员、戏子、哑剧演员、舞者、斗兽士；但一位有信仰的妇女，天主已赐给她其他衣袍，即天主自己的独生子。他说：\BibleQuote{凡你们受洗归于基督的，就是穿上了基督。}\BibleRef{迦 3:27} 请告诉我：如果有人给你君王的衣袍，而你拿了一件乞丐的衣服穿在上面，你除了不雅之外，难道不会受罚吗？你已穿上了天堂与天使的主，却还在忙于地上的事吗？\par

我这样说，因为喜爱装饰本身就是一大恶，即使它不产生其他恶，并且可以没有危险地持有它（因为它煽动虚荣和骄傲），但如今许多其他恶都是由华服产生的：恶意的猜疑、不当的开支、恶言、贪婪的时机。因为你为什么要装饰自己？告诉我。是为了取悦你的丈夫吗？那就待在家里做。但在这里情况恰恰相反。因为如果你想取悦自己的丈夫，就不要取悦别人；如果你取悦了别人，你就不能取悦你自己的丈夫。所以当你去广场或前往教堂时，你应当除去所有的装饰。此外，不要用妓女所用的方法来取悦你的丈夫，而要用自由妇人所用的方法。因为请告诉我，妻子与妓女有何区别？在于一个只关注一件事，即凭身体的美丽吸引她所爱之人；而另一个既管理家庭，又分担子女之事，并参与所有其他事务。\par

你有一个小女儿吗？要当心，免得她继承这种坏风气，因为她们的习惯往往根据教养而形成，并效仿母亲的行为。为你的女儿作谦逊的榜样，用那种装饰打扮自己，并确保你轻视另一种；因为那才是真正的装饰，另一种是毁容。话已足够。现在，愿那创造世界、赐给我们灵魂装饰的天主，以祂的光荣装饰并着装我们，使我们在善行中闪耀，并为祂的光荣而生活，好让我们向圣父、圣子和圣神献上光荣，自今而后，直到永远。\par

\SourceRecord{Translated by John A. Broadus.；Nicene and Post-Nicene Fathers, First Series；Vol. 13.；Edited by Philip Schaff.；Buffalo, NY: Christian Literature Publishing Co.,；1889.；<http://www.newadvent.org/fathers/230310.htm>.}

% structure: p0001 source=h1 target=section reason=page-title

\section{第十一篇讲道辞 论《哥罗森人书》}

\BibleRef{哥 4:5,6}\par

\BibleQuote{你们要凭智慧在外人面前行事，把握时机。你们的言谈要时常带着恩宠，用盐调和，叫你们知道应当怎样回答各人。}\par

基督对祂的门徒所说的，保禄现在也劝告。基督说了什么？\BibleQuote{看，我派遣你们好像羊进入狼群，所以你们要机警如同蛇，纯朴如同鸽子。} \BibleRef{玛 10:16} 那就是，你们要谨慎，不给他们攻击你们的把柄。因此又加了\Quote{对外人}，为叫我们知道，对我们的成员不需要像对外人那样谨慎。因为弟兄们之间，有许多宽容和慈爱。确实，在这里也需要谨慎，但在外更需要，因为在敌人和仇敌中间与在朋友中间是不同的。\par

然后因为他使他们惊慌，看他如何再次鼓励他们：\Quote{把握}，他说，\Quote{时机}：就是说，现在的时候是短暂的。现在他说这话，并不是要他们狡猾或虚伪，（因为这不是智慧的一部分，而是愚昧），而是什么？在那些不伤害你们的事上，他意思是，不要给他们把柄；正如他写信给罗马人时也说：\BibleQuote{向众人还清所欠的：该税的，上税；该征收的，征收；该尊敬的，尊敬。} \BibleRef{罗 13:7} 只是为了宣讲，你才争战，他说，让这战争没有别的起因。因为即使他们因其他原因成为我们的敌人，我们也不会得到赏报，他们会变得更坏，并似乎有正当的抱怨攻击我们。例如，如果我们不纳税，如果我们不给予应得的尊敬，如果我们不谦卑。你们没有看见保禄吗？他在不损害宣讲的地方是何等顺从。因为听他向阿格黎帕王说：\BibleQuote{我自以为有福，因为今天能在你面前为自己辩护，尤其因为我晓得你对犹太人的一切风俗和争论都很精通。} \BibleRef{宗 26:2} 但如果他认为有责任侮辱统治者，他就会毁掉一切。再听真福伯多禄同伴们如何温和地回答犹太人，说：\BibleQuote{应服从天主，而不服从人。} \BibleRef{宗 5:29} 然而那些已经舍弃自己生命的人，既可以侮辱，也可以做任何事；但他们舍弃生命的目的，不是为了追求虚荣，（因为那条路是虚荣的），而是为了放胆宣讲一切。那另一种行为标志着缺乏节制。\par

\Quote{你们的言谈要时常带着恩宠，用盐调和}；就是说，这种恩宠不至于沦为不拘一格。因为单纯地令人愉快是可能的，恰如其分地也是可能的。\Quote{叫你们知道应当怎样回答各人。} 所以不应该对所有人都用同样的言谈，我是说对希腊人和弟兄们。断乎不可，因为这是极端的愚昧。\par

第7节。\BibleQuote{关于我的一切事，提希苛——这位亲爱的弟兄、忠信的仆役、在主内的同仆——会告诉你们。} \BibleRef{哥 4:7}\par

令人赞叹！保禄的智慧何等伟大！请注意，他没有把一切都写在书信里，只写了必要和迫切的事。首先，他不愿把书信拉得太长；其次，为了使他的使者更加受尊重，因为他也有事要传达；第三，显示他对使者的感情，否则他不会把这些事托付给他。然后，有些事情不应以书面形式宣布。他说\Quote{亲爱的弟兄}。如果是亲爱的，他就知道一切，没有向他隐瞒任何事。\Quote{忠信的仆役和在主内的同仆}。如果\Quote{忠信}，他不会说假话；如果\Quote{同仆}，他分担了他的苦难，所以他从各个方面汇集了可信赖的依据。\par

第8节。\BibleQuote{我派他到你们那里去，正是为了这个目的。} \BibleRef{哥 4:8}\par

这里他表明了他的大爱，因为他正是为此目的派他去，这就是他行程的原因；所以当写信给得撒洛尼人时，他说：\BibleQuote{因此，我们不能再忍耐，就决定独自留在雅典，并派了我们的弟兄弟茂德去。} \BibleRef{得前 3:1} 又给厄弗所人派了同一个人，为同样的原因：\BibleQuote{叫他知晓你们的情况，并安慰你们的心。} \BibleRef{弗 6:21} 请看他说什么，不是\Quote{叫你们知道我的事}，而是\Quote{叫我知道你们的事}。这表明他们也处于考验中，用\Quote{安慰你们的心}这句话。\par

第9节。\BibleQuote{一同去的还有敖乃息摩，这位亲爱的忠信弟兄，他是你们那里的人。他们会把这里一切所做的事都告诉你们。} \BibleRef{哥 4:9}\par

敖乃息摩就是他在写给费肋孟的书信中所说的那个：\BibleQuote{我很想留他在我这里，让他代你服事我，我因福音的缘故被囚，但没有你的同意，我什么也不愿做。} \EditorialNote{参看费肋孟书 1:1-3, 14} 然后他又加上对他们城市的赞美，使他们不仅不以为耻，反而以他为荣。他说\Quote{他是你们那里的人}。\Quote{他们会把这里一切所做的事都告诉你们。}\par

第10节。\BibleQuote{我的同囚者阿黎斯塔苛问候你们。} \BibleRef{哥 4:10}\par

没有比这更高的赞美了。这就是那个和保禄一起从耶路撒冷被带上来的人。这个人说了比先知更伟大的话；因为他们自称是\Quote{客旅和外人}，但这一位甚至自称是囚犯。就像一个战俘被拖来拖去，躺在每个人的意志之下受他们的恶待，甚至比战俘更糟。因为敌人的确在俘获他们之后，对他们十分照顾，关心他们如同自己的财产；但保禄，却像敌人和仇敌一样，所有人拖着他到处走，打他，鞭打他，侮辱他，诽谤他。这对那些（他写信给他们的人）也是一种安慰，因为他们的主人也处在这样的境况中。\par

\Quote{及巴尔纳伯的表弟马尔谷}；他仍凭此关系称赞此人，因为巴尔纳伯是一位伟大人物；\Quote{关于他，你们曾领受命令：他若到你们那里，要接待他。}为何？难道他们不会接待他吗？不，他意思是说，要以极大的关注接待他；这也显明此人是伟大的。他们从何处领受这些命令，他并未说明。\par

11节. \Quote{及号称犹斯托的耶稣。}\BibleRef{哥 4:11}\par

此人大概是位格林多人。其次，他对众人给以共同的称赞，之前已逐一称许了每一个人；\Quote{他们是受过割礼的人：这些人只是我的同工，为天主的国，他们曾安慰了我。}说了\Quote{同囚}之后，为免使听者的心灵沮丧，请看他是如何用这句话来振奋他们。他说：\Quote{同工，}\Quote{为天主的国。}如此，既同受苦，便也同享国。\Quote{他们曾安慰了我。}他显明他们是伟大的人物，因为连保禄也得了他们的安慰。\par

但让我们看看保禄的智慧。他说：\Quote{你们在外人当中要明智地行事，赎回时光。}\BibleRef{哥 4:5}即，时间不是你们的，而是他们的。所以不要固执己见，而要赎回时光。他不仅说\Quote{买，}而且说\Quote{赎回，}以另一种方式使之成为你们的。因为制造战争与敌意的借口，简直是极度的疯狂。除了遭受无谓而无益的危险之外，还有这额外的害处：希腊人不会归向我们。因为你在弟兄中间时，应该勇敢；但在外人中间时，就不应如此。\par

你看见他如何处处提及那些外人，即希腊人吗？因此，在写信给弟茂德时，他也说：\Quote{此外，他也必须从外人有好声望。}\BibleRef{弟前 3:7}又说：\Quote{因为审判外人，与我何干？}\BibleRef{格前 5:12}\Quote{你们要在外人当中明智地行事。}因为他们是\Quote{外人}，即使他们与我们同住在这世界，但他们仍在国度和父家之外。他也用称别人为\Quote{外人}来安慰他们，如先前所说：\Quote{你们的生命已与基督一同藏在天主内。}\BibleRef{哥 3:3}\par

然后，他说，你们寻求荣耀，寻求尊荣，寻求那些事物，但现在不可如此，而应让给那些外人。其次，免得你以为他在说钱财，他加说：\Quote{你们的言语要常带着恩宠，用盐调和，好叫你们知道该怎样回答各人。}这样就不会充满虚伪，因为这不是\Quote{恩宠}，也不是\Quote{用盐调和}。例如，如果需要在不招致危险的情况下向某人讨好，不要拒绝；若场合需要你礼貌地谈论，不要以为是奉承；凡关于尊敬的事，只要不损害虔诚，都可去做。你没有看见达尼尔如何向一个不虔敬的人讨好？你没有看见那三个青年如何明智地行事，表现出勇气和言语的胆量，却没有任何鲁莽或刺人的话？因为那样就不是胆量，而是虚荣。他说：\Quote{好叫你们知道该怎样回答各人。}因为对统治者应有一种回答，对被统治的应有另一种；对富有一种回答，对穷人有另一种。何故？因为富人和有权势者的灵魂更软弱、更易激动，因此对他们当用迁就；穷人和被统治者的灵魂更坚强、更聪明，因此对这些当用更大的直言。只着眼一点：他们的建树。并不是因为有人富、有人穷，就要更尊敬前者、更不尊敬后者，而是因为前者软弱，要扶持他，后者则不必。例如，若无必要，不要称希腊人为\Quote{不洁}，也不要侮辱；但若有人问及你的道理，你就回答它是不洁的、不虔敬的；但若无人问你，也没有逼你说话，你就不应无端挑起他的敌意。何必为你自己准备无故的敌视呢？又如，你若教导某人，就谈论当前的主题，否则就沉默。若言辞\Quote{用盐调和}，落在一个松散的心灵中，它会收紧其松弛；落在一个粗糙的心灵中，它会抚平其粗野。要让言语有恩宠，既不过于严厉，也不过于软弱，而要同时带有严肃和愉快。因为若过分严厉，弊多利少；若过分随和，则痛苦多于愉快。所以凡事都当有节制。不要垂头丧气、面容阴沉，因为那令人不快；也不要完全松懈，因为那易遭轻视和践踏。而要像蜜蜂那样，采集每一种德性，从前者取其愉快，从后者取其庄重，避免其缺点。因为医生若不能对所有的身体都使用同样的疗法，那么教师更是如此。不过，身体比灵魂更能承受不适当的药物。例如，一个希腊人来到你面前，成了你的朋友；暂时不要与他谈论此事，直到他成为密友；之后，逐渐为之。\par

请看保禄到了雅典时，如何与他们谈论。他没有说：\Quote{你们这些不洁的、全然不洁的人！}而是说什么？\Quote{雅典人啊，我看你们在各方面都有些迷信。}\BibleRef{宗 17:22}又如，当必须斥责时，他也不推辞；他用极大的严厉对厄吕玛说：\Quote{你充满各种欺诈和奸恶，魔鬼的儿子，一切正义的仇敌。}因为若对前者加以斥责，乃是愚蠢；而不斥责后者，则是软弱。又如，你因事务被带到一位统治者面前，要记得给予他应有的尊敬。\par

第9节。\BibleQuote{他们要把这里的一切事告诉你们，}他说，\BibleQuote{所有在这里所做的事。}为什么你没有和他们一起来，有人问？但\BibleQuote{他们要把一切事都告诉你们}是什么意思？那就是我的锁链，以及所有其他留住我的事。那么我，既祈祷见到他们，也派遣了别人，自己本不应留下，若非有什么重大的必要留住我。然而，这并非指责的话——是的，是强烈的指责。因为让他们确信他既已陷入试炼，并且正英勇地承受它们，乃是证实此事并提升他们心灵的一部分。\BibleRef{哥 4:9}\par

第9节。\BibleQuote{并与\Person{敖尼息摩}一起，}他说，\BibleQuote{这位亲爱的、忠信的弟兄。}\BibleRef{哥 4:9}\par

\Person{保禄}称一个奴隶为弟兄，这是有理由的；因为他称自己为信众的仆人。\BibleRef{格后 4:5}让我们全都放下骄傲，践踏我们的自夸。\Person{保禄}称自己为奴隶，他值得世界和千万重天；而你竟自高自大？他随意夺取一切作为掠物，他在天国中居首位，他得了冠冕，他升到第三层天，却称仆人为\Quote{弟兄}和\Quote{同仆}。你的疯狂在哪里？你的傲慢在哪里？\par

\Person{敖尼息摩}变得如此可靠，以至于连这样的事都托付给他。\par

第10节。\BibleQuote{还有\Person{马尔谷}，}他说，\BibleQuote{\Person{巴尔纳伯}的表弟，关于他你们曾领受命令，要接待他。}也许他们是从\Person{巴尔纳伯}领受了命令。\BibleRef{哥 4:10}\par

第11节。\BibleQuote{那些属于割损的人。}他抑制了\Person{犹太人}膨胀的骄傲，并激励这些（\Place{哥罗森}人）的心灵，因为他们中只有少数属于割损，大多数是外邦人。\BibleRef{哥 4:11}\par

\BibleQuote{这些人曾是我的安慰，}他说，\BibleQuote{是安慰我的人。}他表明自己正处于巨大的试炼之中。所以这也不是一件小事。当我们通过临在、言语、殷勤服侍来安慰圣人们时，当我们与他们一同受苦时（因为他说\BibleQuote{好像与被捆绑的人同受捆绑}；\BibleRef{希 13:3}），当我们把他们的苦难当作自己的苦难时，我们也必在他们的冠冕上有份。你没有被人拖到场上去吗？你没有进入竞技场吗？是别人赤身露体，别人摔跤；但如果你愿意，你也会成为分享者。给他抹油，成为他的支持者和拥护者，在赛场外为他高声呼喊，振奋他的精神，恢复他的精神。由此可知，在一切其他情况下也应如此行。因为\Person{保禄}并不需要，但为了激励他们，他说了这些话。所以你在一切别人身上，要堵住那些辱骂这样的人的口，为他争取支持者，当他出来时以极大的关注接待他，这样你就能在他的冠冕上有份，在他的荣耀上有份；即使你不做别的，只是对所做之事感到喜悦，即便这样你也并非普通地有份，因为你贡献了爱德——一切美善的总和。\par

因为如果那哭泣的似乎分担了忧伤者的悲伤，大大安慰了他们，并消除了他们过度的痛苦，那么那些与人同乐的人，更能使他们的快乐增大。因为没有人分担悲伤是何等大的恶事，听先知说：\BibleQuote{我期待有人与我一同哀悼，却没有一个。}因此\Person{保禄}也说：\BibleQuote{与喜乐的人一同喜乐，与哭泣的人一同哭泣。}\BibleRef{罗 12:15}增加他们的快乐。如果你看见你的弟兄受到好评，不要说：\Quote{这好评是他的，我为何要高兴？}这些话不是弟兄的话，而是仇敌的话。如果你这样想，那就不是他的，而是你的。你更有能力使它增大，如果你不垂头丧气，反而高兴，喜悦，快乐。事实正是如此，从这一点可以明显看出：嫉妒者不仅嫉妒那些受好评的人，也嫉妒那些为他们的好评而高兴的人，因为他们意识到这些人也与那好评有关，而正是这些人最以此为荣。因为后者甚至在受到过度赞扬时会脸红，但这些人却以极大的喜悦自豪。你们不见在运动员身上，一个戴冠冕，另一个没有被冠冕，但高兴和悲伤是在支持者和反对者中，正是他们跳跃、欢腾。\par

看，不嫉妒是多么伟大的事。辛劳是别人的，快乐是你的；别人戴冠冕，而你跳跃、欢快。因为告诉我，既然别人得胜，你为什么跳跃？但他们也深知，所做的事是共通的。因此他们并不确实指责这个人，而是试图击败这场胜利；你听见他们说出这样的话：\Quote{我曾在那里除掉你，}和\Quote{我曾打败你。}虽然行为是别人的，但赞美却是你的。但如果在外在事物上，不嫉妒而是将别人的好处变成自己的，是如此大的善事，那么当魔鬼对我们得胜时，他更加猛烈地呼吸，显然是因为我们更高兴。他虽是邪恶、刻薄的，却深知这快乐是伟大的。你想使他痛苦吗？就高兴欢喜。你想使他高兴吗？就面露愁容。他从你兄弟的胜利所得的痛苦，你以忧愁来抚慰；你与他站在一边，脱离了你的兄弟，你造成了比他更大的祸害。因为一个敌人做敌人的事，与一个朋友与敌人站在一起，并不相同；这样的人比敌人更可憎。如果你的兄弟因说话或光辉的成就而获得美名，你就要成为他声誉的分享者，表明他是你的肢体。\par

\Quote{怎麼會這樣？}有人說，\Quote{因為名聲不是我的。}絕不要這樣說。閉上你的嘴。你若在我身邊，說這樣話的人，我甚至會用手摀住你的嘴，免得敵人聽見。我們常常彼此有仇，卻不向敵人揭露；你倒向魔鬼揭露嗎？不要這樣說，不要這樣想；反倒要說：\InnerQuote{他是我的肢體，榮耀歸於身體。}\Quote{那麼，}有人說，\Quote{為什麼那些外人沒有這樣想？}因為你的過錯：當他們看見你不把他的喜樂當作自己的，他們也不把它當作你的；但若他們看見你將它視為己有，他們就不敢這樣做，你就會與他同樣顯赫。你沒有因說話獲得名聲；但因分享他的喜樂，你得到的榮耀比他的更大。因為如果愛是一件大事，是萬德的總和，你就得到了它所賜的冠冕；他得到的是演說的冠冕，你得到的是超絕之愛的冠冕；他展現了言辭的力量，你卻用行動踐踏了嫉妒，踩碎了惡眼。所以，按理你更應該受冠冕，你的競賽更輝煌；你不僅踐踏了嫉妒，還做了另一件事。他只有一個冠冕，你卻有兩個，而且兩個都比他那一個更明亮。哪兩個？一個是你戰勝嫉妒所得，另一個是你因愛而環繞的。分享他的喜樂不僅證明你沒有嫉妒，也證明你紮根於愛。他有時會被人性的情感嚴重困擾，例如虛榮；但你卻不受任何情感困擾，因為你為他人的好處喜樂，這不是出於虛榮。他重整了教會，告訴我？他增加了會眾？稱讚他；你又得到雙重冠冕；你擊倒了嫉妒；你用愛環繞了自己。是的，我懇求和祈求你們。你們還要聽第三個冠冕嗎？他受人世間的低聲讚揚，你卻受天使在天上的讚揚。因為展示口才和駕馭情感不是同一回事。這種讚揚是暫時的，那種是永恆的；這種是來自人的，那種是來自天主的；這個人公開受冠冕；你卻在暗中受冠冕，在那裡你的天父看見。若能剝去身體，看到每個人的靈魂，我會讓你們看到這個比那個更尊貴、更光輝。\par

我們要踐踏嫉妒的刺棒，我們是為自己謀益處，親愛的，我們要為自己戴上冠冕。嫉妒別人的人是在與天主作戰，而不是與那人作戰；因為當他看見那人得了恩寵，就難過，並希望教會被拆毀，他不是與那人作戰，而是與天主作戰。告訴我，如果有人裝飾國王的女兒，並因裝飾和美化她而為自己贏得聲譽；另一個人卻希望她衣著不整，並希望那人無法裝飾她；那麼他是在對誰圖謀不軌？是對那人嗎？還是對她和她的父親？照樣，你這個嫉妒的人，是在與教會作戰，是在與天主交戰。因為你兄弟的好聲譽也與教會的益處交織在一起，若一個被毀，另一個也必被毀。所以在這方面，你也做了撒殫的行為，因為你對基督的身體圖謀不軌。你為這個人難過嗎？錯，當他絲毫沒有得罪你；其實，你更是為基督難過。祂怎麼得罪了你，以至於你不願讓祂的身體被妝飾得美麗？你不願讓祂的新娘被裝飾？請想想懲罰有多大。你使你的敵人高興；而且，你嫉妒的那個有聲望的人，你想因嫉妒使他難過，你反倒使他更高興；你的嫉妒反而更顯出他有聲望，否則你不會嫉妒他。你反而顯出你是在受罰。\par

我實在羞於從這樣的動機來勸勉你們，但看到我們的軟弱如此之大，我們還是要從這些事上受教，使自己脫離這種毀滅性的情感。你因他聲望高而難過嗎？那你為什麼還要因嫉妒而助長他的聲望？你想懲罰他嗎？那你為什麼表現出痛苦？為什麼在他之前懲罰你自己，而你又不想讓他聲望高？此後他的快樂會加倍，你的懲罰也加倍；不僅因為你證明了他偉大，還因為你因懲罰自己而在他心中生出另一種快樂；而當你難過時，他卻高興，你卻嫉妒。看，我們如何不知不覺地重擊自己！他是敵人。然而，為什麼是敵人？他做了什麼錯事？但是，我們這樣做卻使我們的敵人更加顯赫，從而更懲罰自己。而且在這一點上，如果我們發現他知道，我們又懲罰自己。因為也許他不高興，但我們以為他是這樣，就又因此難過。那麼停止你的嫉妒吧。為什麼要給自己造成傷口？\par

我們要思想這些事，親愛的：那兩個不嫉妒的人所得的冠冕；那些來自人的讚揚，那些來自天主的讚揚；嫉妒所帶來的邪惡；這樣我們就能制服那獸性，在天主面前蒙悅納，並與那些有聲望的人得到同樣的賞報。因為或許我們會得到它們，即使得不到，也是對我們有益；然而，即便如此，如果我們為天主的榮耀而生活，我們就能得到祂為愛祂的人所預備的恩許，藉著我們的主耶穌基督的恩寵和對人的慈愛，願光榮歸於祂，等等。\par

\SourceRecord{Translated by John A. Broadus.；Nicene and Post-Nicene Fathers, First Series；Vol. 13.；Edited by Philip Schaff.；Buffalo, NY: Christian Literature Publishing Co.,；1889.；<http://www.newadvent.org/fathers/230311.htm>.}

% structure: p0001 source=h1 target=section reason=page-title

\section{《哥罗森书》第十二篇讲道}

\BibleRef{哥 4:12-13}\par

\BibleQuote{你们的同乡\Person{厄帕夫辣}，\Person{基督耶稣}的仆人，问候你们，他在祈祷中时常为你们奋斗，使你们在\Person{天主}的一切旨意上得以完善，并且坚定不移。我为他作证，他对你们以及在\Place{劳狄刻雅}和\Place{希耶拉波利斯}的人，怀有很大的热忱。}\par

在这封信的开端，他也称赞此人的爱心；因为称赞也是爱心的标记。所以他在开头说：\BibleQuote{他也曾在圣神中向我们报告了你们的爱心。}\BibleRef{哥 1:8}。为别人祈祷也是爱心的标记，并且能再次激发爱心。他称赞他，也是为了给其教导打开一扇门，因为教师之可敬是门徒的益处。同样，他说\BibleQuote{你们的同乡}\BibleRef{哥 4:12}，是为了让他们以这人为荣，因为产生了这样的人。他又说：\BibleQuote{他在祈祷中时常为你们奋斗}\BibleRef{哥 4:12}。他没有简单地说\BibleQuote{祈祷}，而是说\BibleQuote{奋斗}，表明颤抖和恐惧。\BibleQuote{我为他作证}\BibleRef{哥 4:13}，他说，\BibleQuote{他对你们怀有很大的热忱}。这是一个可信的见证。\BibleQuote{他对你们}，他说，\BibleQuote{怀有很大的热忱}，也就是说，他非常爱你们，并且对你们充满炽热的感情。\BibleQuote{以及他们\Place{劳狄刻雅}和\Place{希耶拉波利斯}的人}。他也向那些人称赞他。但那些人如何知道这事？他们当然会听说；不过，当这封信被宣读时，他们也会知道。因为他说：\BibleQuote{也请此信在\Place{劳狄刻雅}教会宣读。}\BibleRef{哥 4:16}\BibleQuote{使你们完善}\BibleRef{哥 4:12}，他说。他既指责他们，又毫无冒犯地给予劝告和指导。因为人可能既完善，却又没有站稳，例如知道一切却仍然摇摆不定；也可能不完善，却站着，例如只知道一部分，却没有坚定地站着。但此人祈求两者兼得：\BibleQuote{使你们完善}\BibleRef{哥 4:12}，他说。请看他是如何再次提醒他们关于天使和生命的话。\BibleQuote{并且坚定不移}\BibleRef{哥 4:12}，他说，\BibleQuote{在\Person{天主}的一切旨意上}。仅仅遵循\Person{天主}的旨意还不够。那\BibleQuote{充满}的人不允许任何其他旨意存于自己内，因为如果那样，他就没有完全充满。\BibleQuote{我为他作证}\BibleRef{哥 4:13}，他说，\BibleQuote{他怀有很大的热忱}。\BibleQuote{热忱}和\BibleQuote{很大的}都是强调。正如他本人在给\Place{格林多}人写信时所说：\BibleQuote{我以\Person{天主}的嫉妒为你们嫉妒。}\BibleRef{格后 11:2}\par

\BibleRef{哥 4:14}。此人是圣史。他把他放在后面，并不是贬低这个人，而是为了抬高另一个人，即\Person{厄帕夫洛迪托}。可能还有其他人叫这个名字。他又说：\BibleQuote{还有\Person{德玛斯}}。在说了\BibleQuote{亲爱的医生\Person{路加}问候你们}之后，他加上了\BibleQuote{亲爱的}。这不是小小的称赞，而是极大的称赞，因为成为\Person{保禄}所爱的人是了不起的。\par

\BibleQuote{请问候在\Place{劳狄刻雅}的弟兄们，以及\Person{宁法}，还有在他们家里的教会。}\BibleRef{哥 4:15}\par

请看他是如何将他们彼此联结和融合在一起的——不仅通过问候，还通过交换书信。然后他又通过单独提及此人来表示敬意。他这样做不是没有原因的，而是为了引导其他人也效仿他的热忱。因为不被视为与众人一样，这不是小事。再看他如何显明这人的伟大：他的家就是一座教会。\par

\BibleRef{哥 4:16}。\BibleQuote{当这封信在你们中间宣读之后，也请它在\Place{劳狄刻雅}教会宣读。}\par

\BibleQuote{你们也要宣读从\Place{劳狄刻雅}来的信。}有人说这不是\Person{保禄}写给他们的，而是他们写给\Person{保禄}的，因为他没有说\BibleQuote{写给\Place{劳狄刻雅}人的}，而是说\BibleQuote{从\Place{劳狄刻雅}来的信}。\par

\BibleQuote{请告诉\Person{阿尔希波}：要谨慎你在主内所接受的职务，务要完成它。}\BibleRef{哥 4:17}\par

\BibleQuote{\Person{保禄}亲笔问候。请记念我的锁链。愿恩宠与你们同在。阿们。}\BibleRef{哥 4:18}。这是他们真诚和爱心的证明；他们既注视他的笔迹，又激动不已。\BibleQuote{请记念我的锁链}。奇妙！这是何等的安慰！因为这就足以鼓励他们去做一切事，并在磨难中更勇敢地担当；但他不仅使他们更勇敢，还使他们更关切。\BibleQuote{愿恩宠与你们同在。阿们。}\par

这是极大的赞扬，且超越一切其余，他提到\Person{厄帕弗辣}\Quote{他是你们中的一位，基督的仆人}。他又称他为他们的仆役，如同他自称为教会的仆役一样，正如他说：\Quote{我保禄成了这福音的仆役}\BibleRef{哥 1:23}。他使此人升至同样的尊位；前面他称他为\Quote{同为仆役}\BibleRef{哥 1:7}，而此处称为\Quote{仆人}。他说\Quote{他是你们中的一位}，仿佛对一位母亲说话，说\Quote{他是你腹中所生的}。但这赞扬可能引起嫉妒；因此他不只从这些事情称赞他，也从与他们自身有关的事称赞；如此他消除了嫉妒，无论在之前还是此处。他说：\Quote{常常为你们奋斗}，不只在现在，当与我们在一起时，为炫耀；也不只在与你们在一起时，为在你们面前炫耀。通过说\Quote{奋斗}，他显示了他的极大热忱。然后，为了不显得是在奉承他们，他加上：\Quote{因他对你们，并对在\Place{劳狄刻雅}和\Place{希耶辣颇里}的人有极大的热忱。}这句话：\Quote{使你们得以成全地站立}，不是奉承之词，而是出自可敬的教师。他既说了\Quote{确信无疑}，又说了\Quote{成全}。前者是他赐予他们的，后者他说是欠缺的。他没有说\Quote{使你们不动摇}，而是\Quote{使你们站立得住}。然而，他们被许多人问候，使他们欣喜，因为不仅他们中间的朋友，而且其他的人也纪念他们。\par

他说：\Quote{并且对\Person{阿尔希颇}说：\InnerQuote{务要谨慎你从主所领受的职务，以完成它。}}他的主要目的是使他们完全服从他。因为他们不能再因他责备他们而抱怨，当时他们自己承担了一切；因为对门徒谈论师父是不合理的。但为堵住他们的口，他这样写信给他们：他说：\Quote{对\Person{阿尔希颇}说：务要谨慎。}这个词处处用于警告；正如他说：\Quote{务要谨慎狗类。}\BibleRef{斐 3:2}\Quote{务要谨慎，免得有人用哲学掳掠你们。}\BibleRef{哥 2:8}\Quote{务要谨慎，免得你们这自由竟成了软弱人的绊脚石。}\BibleRef{格前 8:9}当他想要恐吓时，他总是这样表达自己。他说：\Quote{务要谨慎从主所领受的职务，以完成它。}他甚至不给他选择的权力，正如他自己所说：\Quote{我若甘心作这事，就有赏报；若不甘心，管家的职分却已托付给我了。}\BibleRef{格前 9:17}\Quote{你要完成它}，常常尽力。\Quote{你从主所领受的，要完成它。}再者，\Quote{在}这个词的意思是\Quote{藉着主}。他说，是主赐给你的，不是我们。他也使他们服从他，当他表明他们已被天主交在他手中时。\par

\Quote{你们要纪念我的锁链。愿恩宠与你们同在。阿们。}他解除了他们的恐惧。因为虽然他们的老师被锁链束缚，然而\Quote{恩宠}释放了他。这也是恩宠，即赐予他被锁链束缚。请听路加所说：宗徒们\Quote{离开公议会，心里欢喜，因被算是配为这名受辱。}\BibleRef{宗 5:41}因为既受辱又被锁链束缚，实在就是\Quote{被算为配得}。因为，如果有人爱他的所爱，就认为为他受苦是一种收获，那么为基督受苦更是如此。所以我们不要因基督的缘故而抱怨我们的苦难，反而要纪念保禄的锁链，让这成为我们的激励。例如：你劝诫人为基督的缘故施舍穷人？提醒他们保禄的锁链，哀叹你自己和他们的可怜，因为那一位为了他竟把自己的身体交付锁链，而你却连你食物的一部分都不愿给出。你因善行而自高？纪念保禄的锁链，你并没有遭受那种事，你就不会再自高。你贪恋邻人的任何东西？纪念保禄的锁链，你就会看到这是多么不合理：他在危险中，你却应在享乐中。再者，你心系放纵？在你的脑海中想象保禄的监狱；你是他的门徒，他的战友。你的战友被锁链束缚，而你却奢侈，这合理吗？你处于苦难中？你认为自己被抛弃？听保禄的锁链，你就会明白处于苦难并非被抛弃的证明。你想穿丝绸长袍？纪念保禄的锁链；这些事物在你看来将比那坐在角落的污秽破布更不值钱。你想佩戴金饰？在脑海中想象保禄的锁链，这些事物在你看来将不比一根干枯的灯心草更好。你想梳理头发，变得美丽可看？想想保禄在监狱中的肮脏，你就会渴望那种美丽，认为这是极端的丑陋，并因渴望那些锁链而痛苦呻吟。你想用膏油和颜料涂抹自己？想想他的眼泪：三年之久，昼夜不停哭泣。\BibleRef{宗 20:31}用这装饰来妆点你的脸颊；这些眼泪使它明亮。我不是说你要为别人哭泣（我确实希望如此，但这对你太高了），而是我劝你为自己的罪这样做。你命令将你的奴隶锁起来，你就生气、恼怒？纪念保禄的锁链，你将立即平息怒气；记住我们是被束缚的人，而非束缚者，是心灵受伤的人，而非伤人者。你失去控制，大笑大叫？想想他的哀哭，你将呻吟；这样的眼泪将使你远为明亮。你看到有人狂欢跳舞？纪念他的眼泪。有什么泉源曾涌出如他眼睛泪水的洪流？\Quote{纪念我的眼泪}\BibleRef{宗 20:31}，他说，正如这里说\Quote{锁链}。他这样对他们说话是有理由的，那时他从厄弗所打发人来到米肋托。因为那时他对教师们说话。因此他要求他们也要同情，但这些人只应面对危险。\par

有什么泉源你能比作这些眼泪？乐园中的那泉，灌溉全地？但你将说不到类似的事。因为这眼泪之泉灌溉了灵魂，而非大地。若有人向我们展示保禄泪流满面、呻吟哭泣，这岂不远远胜过看到无数队列欢歌戴冠？我现在不是对你们说；但是，如果有人从戏院和舞台拉来一个放荡的家伙，燃烧着并沉醉于肉体的爱情，然后向他展示一位正值青春芳华的年轻贞女，她在各方面都超越同伴，特别是在面容上更胜于身体其余部分，有一双温柔而柔软的眼睛，安详地注视又安详地转动，湿润、柔和、平静地微笑，穿着极度的谦逊和许多优雅，上下睫毛如羽，眼球，可说是活生生的，前额光彩；下面，脸颊泛出恰到好处的红晕，如大理石般光滑平坦；然后有人向我展示保禄哭泣；我将离开那位少女，急切地奔向看他；因为从他的眼中闪耀出属灵之美。因为那少女会令青年的灵魂兴奋，灼烧并点燃他们；但这相反，却使他们降服。这使灵魂的眼睛更美丽，克制肚腹；充满智慧之爱，充满同情；它能软化金刚石般的灵魂。教会因这些眼泪而受灌溉，灵魂因这些眼泪而被种植；是的，即使有实质的、有形的火，这些眼泪也能将其熄灭；这些眼泪熄灭了恶者的火箭。\par

因此我们要纪念他的这些眼泪，我们就会嘲笑一切现世的事物。这些眼泪基督曾称之为有福，说：\Quote{哀恸的人是有福的，哭泣的人是有福的，因为他们将欢笑。}\BibleRef{玛 5:4}依撒意亚和耶肋米亚也曾流过这样的眼泪；前者说：\Quote{容我独自痛哭，痛哭不已}\BibleRef{依 22:4}；后者说：\Quote{谁给我的头水，给我的眼睛泪水的泉源？}\BibleRef{耶 9:1}；仿佛自然的源泉不够。\par

没有什么比这些眼泪更甘甜；它们比任何欢笑都甘甜。哀恸的人知道这眼泪蕴含着多大的安慰。我们不要认为这是一件该被嫌弃的事，反要极其祈求它；不是祈求别人犯罪，而是当他们犯罪时，我们为他们心碎。让我们记住这些眼泪，这些锁链。确实，那些锁链上也落下了眼泪；但那些灭亡者的死——就是那些捆绑他的人——使他不能尝到锁链的甘甜。因为他为他们忧伤，他是那位为犹太司铎哀哭者的门徒；不是因为他们将要把祂钉在十字架上，而是因为他们自己正在灭亡。祂不仅自己如此行，也这样劝勉别人说：\BibleQuote{耶路撒冷女子，不要为我哭。}\BibleRef{路 23:28} 这双眼睛见过乐园，见过第三层天：但我不认为他们因这个异象而有福，正如因那些眼泪而有福，藉着眼泪他们看见了基督。诚然，那个异象是有福的；因为他自己甚至以此夸口说：\BibleQuote{我岂没有见过我们的主耶稣基督吗？}\BibleRef{格前 9:1}；但这样哭泣更为有福。\par

在那个异象中，许多人都有分；但那些没有如此看见的人，基督反倒称他们有福，说：\BibleQuote{那些没有看见而相信的，是有福的。}\BibleRef{若 20:29} 但达到这地步的不多。因为如果为基督的缘故留在此世比离开祂去\BibleRef{斐 1:23}更为必要，为了他人的救恩；那么为了他人的缘故而叹息，甚至比看见祂更为必要。因为如果为祂的缘故而在地狱里，比与祂同在更可向往；并且为祂的缘故与祂分离，比与祂同在更可向往（这正是他所说的：\BibleQuote{我宁愿我自己被诅咒，与基督隔绝}\BibleRef{罗 9:3}），那么为祂的缘故而哭泣更是如此。他说：\BibleQuote{我昼夜不停地流泪劝诫你们每一个人。}\BibleRef{宗 20:31} 为什么？不是惧怕危险；不；而是如同一个人坐在病人身旁，不知道结局如何，出于深情而哭泣，恐怕他失去生命；他也是如此：当他看见一个人患病，用责斥无法奏效时，便从此流泪。基督也是如此，希望他们能敬畏祂的眼泪：这样，有人犯罪，祂责斥他；被责斥的人朝祂吐唾沫，并远远跳开；祂哭泣，希望藉此能赢得他。\par

让我们记住这些眼泪：这样教养我们的女儿，这样教养我们的儿子；当我们看见他们陷于罪恶时哭泣。凡希望被爱的女子，让她们记住保禄的眼泪而叹息；你们中凡被认为有福的、凡在新娘洞房中的、凡在享乐中的，记住这些；凡在哀恸中的，用眼泪交换眼泪。他不是为死人哀哭，而是为那些活着却正在灭亡的人哀哭。我要提及别的眼泪吗？弟茂德也曾哭泣；因为他是这个人的门徒；因此写信给他时也说：\BibleQuote{忆及你的眼泪，使我充满喜乐。}\BibleRef{弟后 1:4} 许多人甚至因喜乐而流泪。所以这也是一种喜乐，而且是极其强烈的喜乐。所以眼泪并非痛苦：是的，从这种忧伤流出的眼泪，甚至远胜于世俗喜乐的眼泪。听先知说：\BibleQuote{上主听见了我哭泣的声音，听见了我的恳求。}\BibleRef{咏 6:8} 因为何处眼泪没有用处？在祈祷中？在劝诫中？我们给它们带来了坏名声，因为没有用在它们被赐予我们的目的上。当我们恳求一个犯罪的兄弟时，我们应该哭泣，悲伤叹息；当我们劝诫某人，而他毫不理会，继续走向灭亡时，我们应该哭泣。这些是天上的智慧的眼泪。但是当一个人处于贫困、疾病或死亡时，却不必如此；因为这些不是值得流泪的事。\par

正如当我们不恰当地使用欢笑时，我们给欢笑带来坏名声；同样，如果我们不合时宜地使用眼泪，也给眼泪带来坏名声。因为每件事的美德就在它被用于其适当工作时才显现，若用于不当之处，便不再如此。例如，酒是为使人喜乐，而非醉酒；食物是为营养；性交是为生育子女。正如这些事已经得了坏名声，眼泪也是如此。若有这样一条法则：它们只用于祈祷和劝诫，那么看它们会成为多么可渴望的事。没有什么比眼泪更能抹去罪过。眼泪甚至使这身体的面容显得美丽；因为它们赢得旁观者的怜悯，使它在我们的眼中受到尊重。没有什么比含泪的眼睛更甘甜。因为这是我们最高贵的器官，最美丽的，且是灵魂所特有的。因此我们如此被它打动，仿佛看见了灵魂本身在哀叹。\par

我说这些话并非没有理由；而是为了使你们不再参加婚礼、舞蹈和撒殚的表演。因为看魔鬼发明了什么。既然自然本身使妇女远离舞台及那里上演的丑事，他便把戏院的设施引入了闺房，我说的是放荡的男人和娼妓。这种灾祸是由婚姻习俗引入的，或者更确切地说，不是婚姻造成的，绝不然，而是我们自己的愚蠢。你做的是什么事，人啊？你不知道你在干什么吗？你娶妻是为了贞洁和生育子女；那么这些娼妓是做什么的？有人回答说：为了更大的欢乐。难道这不是更疯狂吗？你侮辱了你的新娘，侮辱了被邀请的妇女。如果她们对这样的表演感到高兴，那正是侮辱。如果看娼妓做出不雅行为能带来任何荣誉，那你为什么不把你的新娘也带去，让她也看看呢？将淫荡之人、舞者和所有撒殚的排场引进家中，是非常不体面和可耻的。\par

\Quote{记住，}他说，\Quote{我的锁链。}婚姻是一种锁链，是天主所定的锁链；娼妓是一种分离与拆散。允许你们用别的事物装饰婚姻，例如丰盛的筵席和服饰。我并不禁止这些，免得我显得过分粗野；然而\Person{黎贝加}只满足于她的面纱\BibleRef{创 24:65}，但我仍然不禁止它们。允许你们藉着服饰、藉着可敬男女的临在来装饰和点缀婚姻。你为何引进那些嘲弄？那些怪物？请告诉我们，你从他们那里听到什么？什么？你羞于启齿？你既感到羞耻，为何还强迫他们做这事？如果这事是光荣的，为何你自己不做？如果是可耻的，为何你强迫别人？一切应充满纯洁、庄重、秩序；但我看到的却是相反，人们像骆驼和骡子一样蹦跳。至于贞女，她的内室才是唯一合适的地方。有人说：\Quote{但她是贫穷的。}因为她贫穷，她也应当谦逊；让她的品德代替财富。难道她没有嫁妆可带？那你为何透过她的生活和举止使她更受轻视？我称赞这个习俗，贞女们出席以尊敬她们的同伴；已婚妇女们出席以尊敬那被纳入她们行列的人。这样的安排是正确的。因为这里有两群人，一是贞女，一是已婚者；前者将她交出，后者将她接收。新娘处于两者之间，既非贞女，亦非妻子，因为她正从那里出来，进入与这些人的共融。但那些娼妓，她们是什么意思？当婚礼举行时，她们应当掩面，应当被埋入地里（因为娼妓是婚姻的败坏），但我们却在我们的婚礼上引进她们。当你从事任何工作时，你视作不祥，不愿说出哪怕一个与它相反的词语；例如，当你播种、当你从酒槽中取酒时，你即使被问及，也不会说出关于醋的一个音节；但在这里，目标是纯洁，你们却引进醋？因为娼妓就是这样。当你准备香膏时，你不允许任何恶臭靠近。婚姻是香膏。那么，你为何在准备你的香膏时引进粪堆的恶臭？你说什么？贞女可以跳舞，而在她的同伴面前不感到羞耻？因为她应当比另一位更庄重；她至少是从（奶妈的）臂膀中出来，而非从\OriginalTerm{角力场}{palæstra}出来。因为贞女根本不应在婚姻中公开出现。\par

难道你没有看见，在王宫里，受尊敬的人在内，靠近君王，不受尊敬的人在外？你也应在内，靠近新娘。但你要在纯洁中留在家中，不要暴露你的童贞。两群人都在旁站立，一群是为了显示他们所交出的新娘是怎样的，另一群是为了守护她。你为何羞辱贞女的身份？因为如果你是这样的，新郎也会怀疑她如此。如果你想让人爱上你，这是卖货女、菜贩和手艺人的作为。这不是羞辱吗？行为不端就是羞辱，即使是国王的女儿。难道她的贫穷妨碍她吗？或是她的生活方式？即使贞女是奴隶，也让她持守谦逊。\BibleQuote{因为在基督耶稣内，没有奴隶或自由人之分。}\BibleRef{迦 3:28}\par

什么？婚姻是个剧场？它是一个奥秘，是一件伟大事物的预像；即使你不尊敬它，也请尊敬它所预像的那一位。\BibleQuote{这个奥秘是伟大的，但我是指基督和教会说的。}\BibleRef{弗 5:32}它是教会的预像，是基督的预像，而你在其中引进娼妓？有人说，如果贞女不跳舞，已婚者也不跳舞，那么谁来跳舞？无人，因为跳舞有何必要？在希腊人的奥秘中有舞蹈，但在我们的奥秘中，是沉默与端庄、谦逊与羞涩。一个伟大的奥秘正在举行：娼妓出去！亵渎的人出去！它怎么是奥秘？他们结合，两人成为一体。为何在他进入时，没有舞蹈，没有铙钹，而是极大的沉默，极大的宁静；但当他们结合，制造出不是无生命的像，也不是地上任何事物的像，而是天主本身的像，并按照祂的肖像时，你却引进如此大的喧闹，扰乱在场的人，使灵魂感到羞耻，困惑它？他们前来，将要成为一体。看，再一个爱的奥秘！如果两人不成为一体，只要他们继续是两个，他们就不会生出许多；但当他们进入一体时，他们才生出许多。我们从中学到什么？即合一的力量是伟大的。起初天主的智慧将那一个分为两个；并想要显示即使在分开之后它仍然是一个，祂不容许那一个单独足以繁殖。因为尚未（结合）的那一个并不是一个，而是半个；这一点从他不产生后代可以清楚看出，正如在先前也是如此。你看见婚姻的奥秘了吗？祂从一造出一，然后又使这两个成为一，祂就这样造出一，使得现在人也从一产生。因为男人和妻子不是两个人，而是一个人。这可以从许多来源得到证实：例如，从\Person{雅各伯}，从基督的母亲\Person{玛利亚}，从这话：\BibleQuote{祂造了他们，有男有女。}\BibleRef{创 1:27}如果他是头，她是身体，他们怎能是两个？因此，一个处于门徒的地位，另一个处于教师的地位；一个处于统治者的地位，另一个处于服从的地位。此外，从她身体的形成本身，人可以看到他们是一体的，因为她是由他的肋骨造成的，他们就如两个半体。\par

为此，祂也称她为助手，以显示他们是一体；\BibleRef{创 2:18}为此，祂尊崇他们的同居超过父母，以显示他们是一体。\BibleRef{创 2:24}同样，父亲在儿子和女儿结婚时欢喜，好像身体急于连接自己的肢体；虽然要承担如此大的费用和开支，他仍不能漠不关心地看她不嫁。因为好像她自己的肉本身被切断了一样，每个人单独对于生育子女都是不完全的，每个人对于今生生活的构成都是不完全的。因此先知也说：\BibleQuote{你精神的残余。}\BibleRef{拉 2:15}他们如何成为一体？好像你取了金子最纯净的部分，与别的金子混合；同样，这里的女人好像接受了最丰富的部分，被快感熔化，滋养它、珍惜它，并且贡献她自己的份额，送回一个人。孩子像一座桥，使三者成为一体，孩子连接两边，彼此结合。因为如同两座城市被一条河流分割，如果有桥连接两边，就成为一体；这里更是如此，而且这里的桥是由双方的实体构成的。如同身体与头是一体；它们被脖子分隔；但分隔不如连接紧密，因为脖子在它们中间将彼此连接。又如同一个被分割的合唱团，通过从这边取一部分，再从右边取另一部分，合为一体；或者如同它们紧密排列，伸出手臂，成为一体；因为伸出的手臂不允许它们是两个。所以祂说得准确，不是说\BibleQuote{他们将成为一体}，而是结合在一起\BibleQuote{成为一体}\BibleRef{创 2:2}，即孩子的身体。那么，如果没有孩子，他们不是两个吗？不，因为他们的结合产生这样的效果：它扩散并混合双方的肉体。如同一个人把油膏倒入油中，使整体成为一体；这里也是如此。\par

我知道许多人因所说的感到羞耻，其原因就是我所说的，你们自己的好色和不洁。婚姻被如此执行、如此败坏的事实使这事蒙上恶名：因为\BibleQuote{婚姻是尊贵的，床榻是无玷污的。}\BibleRef{希 13:4}为什么你们为尊贵的感到羞耻，为什么你们为无玷污的面红耳赤？这是为了异端者，这是为了那些引进娼妓的人。为此我渴望彻底净化它，使其恢复它应有的尊贵，堵住异端者的口。天主的恩赐被侮辱，我们世代的根源；因为那根源有许多粪土和污秽。所以让我们用我们的言论清除它们。忍耐一会儿，因为拿粪土的人必须忍受臭味。我愿意向你们显示，你们不应当对这些事感到羞耻，而应对你们所做的事感到羞耻；但是你们，忽视对那些事的羞耻，对这些事感到羞耻；那么你们确实谴责了这样规定的天主。\par

我是否要告诉你们婚姻也是教会的奥秘？正如基督来到教会，她是由祂而造，祂与她以精神的交合结合，\BibleQuote{因为，}有人说，\BibleQuote{我已把你们许配给一个丈夫，一个纯洁的贞女。}\BibleRef{格后 11:2}而且我们是出于祂，祂说，祂的肢体，\BibleQuote{和祂的肉。}那么，想到这一切，我们不要让这伟大奥秘蒙羞。婚姻是基督临在的预象，而你们在它上面醉酒吗？告诉我；如果你们看见国王的肖像，会侮辱它吗？绝不。\par

如今婚礼上的做法似乎是无关紧要的事，但它们是大恶行的原因。一切充满不法。\BibleQuote{污秽、愚妄的谈吐和戏谑，不可出自，}他说，\BibleQuote{你们的口。}\BibleRef{弗 5:4}现在所有这些是污秽、愚妄的谈吐和戏谑；而且不单如此，还变本加厉，因为这事已成为一种艺术，追求它的人大受赞扬。罪恶已成为艺术！我们不是随意地追求它们，而是认真、有学问地，从此魔鬼掌握了自己的阵列。哪里有醉酒，哪里就有不洁：哪里有污秽谈吐，哪里就有魔鬼在近旁加入他的贡献；用这样的娱乐，告诉我，你们庆祝基督的奥秘吗？你们邀请魔鬼吗？\par

我敢说你们认为我冒犯。因为这也是极端乖张的特性，以至于责备你们的人反而被你们嘲笑为严峻。你们难道没听见保禄说：\BibleQuote{你们无论做什么，或吃或喝，或无论做什么，都要为光荣天主而做}？\BibleRef{格前 10:31}但你们做一切是为了坏名声和羞辱。你们难道没听见先知说：\BibleQuote{应以敬畏侍奉上主，战战兢兢向祂欢乐}？\BibleRef{咏 2:11}但你们完全不受约束。既享受快乐又安全地做，难道不可能吗？你们想听优美的歌曲吗？其实最好不听；但若你们执意，我迁就：不要听那些撒殚的歌曲，而要听精神的歌曲。你们想看舞蹈的合唱队吗？看天使的合唱队。有人问：怎么可能看见他们？如果你们驱除所有这些事，基督甚至也会来到这样的婚姻，基督临在，天使的合唱队也同在。如果你们愿意，祂现在也行奇迹，如同那时一样；祂甚至现在也可以把水变成酒\EditorialNote{若 2}；更奇妙的是，祂会转变这种不稳定且消散的快乐、这种冷淡的欲望，转化为精神的。这就是变水为酒。哪里有吹笛手，基督绝不会在那里；即使祂进入，祂也会先赶出这些，然后行奇迹。有什么比这种撒殚的排场更讨厌？其中一切都无意义，一切无意义；如果有什么有意义，也都是可耻的、有害的。\par

没有比德行更令人愉悦的，没有比条理更甜美的，没有比庄重更可爱的。让任何人庆祝我所说的这样的婚姻；他必会找到快乐；但要留心这些婚姻是怎样的。首先，为那位童贞女寻找一个真正能作丈夫和保护的人；好像你打算为身体安上一个头；好像你即将把她交托出去的不是一个奴隶，而是一个女儿。不要寻求金钱，也不要家世的显赫，更不要国家的伟大；这一切都是多余的；而是要寻求灵魂的虔诚、温柔、真正的智慧、敬畏天主，如果你希望你的宝贝快乐地生活。因为如果你寻找一个更富有的丈夫，你不仅不会使她受益，反而会伤害她，使她成为奴隶而不是自由人。因为她从金饰中获得的快乐，不会比她因奴役而来的烦恼更大。我恳求你，不要寻求这些东西，而最重要的是，找一个条件相等的人；如果做不到，那就找一个贫穷的，而不是条件更好的；至少如果你不是想把女儿卖给一个主人，而是想把她嫁给一个丈夫。当你彻底调查了那个人的德行，并准备把她交给他时，恳求基督临在；因为祂不会耻于如此；这是祂临在的奥秘。不，不如首先恳求祂赐给她这样一个求婚者。不要比亚巴郎的仆人更差，他在如此重要的朝圣之旅中，看到了他应该求助的地方；因此他也得到了一切。当你焦虑不安，为她寻找丈夫时，要祈祷；对天主说：\Quote{祢所愿意的，求祢预备。}将这件事交在祂手中；祂因你如此尊崇祂，必会以尊荣回报你。\par

确实有两件事必须做：将事情交在祂手中，并寻求一个祂所认可的规矩人。\par

那么，当你举办婚礼时，不要挨家挨户地借镜子和礼服；因为这事不是为了炫耀，你也不是带领女儿去参加游行；而是用你家里的东西装饰你的房子，邀请你的邻居、朋友和亲戚。所有你认识的有良好品德的人，都邀请他们，并请他们满足于现有的。不要让乐师在场，因为这样的开销是多余的，而且不相称。最重要的是，邀请基督。你知道怎样邀请祂吗？祂说：\BibleQuote{你们对我这些最小兄弟中的一个所做的，就是对我做的。}\BibleRef{玛 25:40}不要以为为了基督的缘故邀请穷人是一件烦人的事；邀请娼妓才是一件烦人的事。因为邀请穷人是致富之道，而邀请娼妓则是毁灭之道。不要用这些黄金制成的饰物来装饰新娘，而是用温柔、谦逊和惯常的衣裳；不用一切金饰和编发，而用羞涩、腼腆和不渴望这些东西来装扮她。不要有喧闹，不要有混乱；让新郎被召来，让他迎接童贞女。晚餐和宴席，不要充满醉酒，而要充满丰盛和快乐。看，当我们看到这样的婚姻时，会有多少好事发生；但从现在所庆祝的婚姻中（如果至少应该称它们为婚姻而不是游行的话），有多少罪恶啊！宴席一散，立刻就有忧虑和恐惧，唯恐借来的东西丢失了，快乐之后便随之而来难以忍受的忧郁。但这种痛苦属于岳母——不，甚至新娘本人也不自由；至少所有后续的事都属于新娘本人。因为看到一切都散了，是悲伤的理由，看到房子空荡荡的。\par

那里有基督，这里有撒殚；那里有快乐，这里有忧虑；那里有愉悦，这里有痛苦；那里有花费，这里没有类似的事；那里有失态，这里有谦逊；那里有嫉妒，这里没有嫉妒；那里有醉酒，这里有清醒，这里有健康，这里有节制。牢记这一切，让我们从这一点上止住罪恶，好使我们能中悦天主，并配得获得为爱祂的人所预许的美善事物，借着我们的主耶稣基督的恩宠和慈爱，与祂一同，偕同圣神，归于圣父，愿光荣、权能、尊崇，从现在直到永远，世世无尽。阿们。\par

\SourceRecord{Translated by John A. Broadus.；Nicene and Post-Nicene Fathers, First Series；Vol. 13.；Edited by Philip Schaff.；Buffalo, NY: Christian Literature Publishing Co.,；1889.；<http://www.newadvent.org/fathers/230312.htm>.}
